% !TeX program = lualatex
% Deterministically generated from the frozen D019 corpus.
% Compile only with LuaLaTeX or XeLaTeX. No shell escape is required.
% D019-REPAIR-MAP-SHA256: DFC270E5AEB4D9AE859B215CAA9BB168DDFB67B1194D9DD4D294050DD7785831
% D019-INHERITED-CORPUS-SNAPSHOT-SHA256: 3C1F81CFF7BC6EF172780059366B24761F3E69EE65AD32F0B31D6CFB5692A658
\documentclass[11pt,twoside]{article}

\usepackage{iftex}
\ifPDFTeX
  \PackageError{D019 canonical build}{pdfTeX is unsupported}{Use LuaLaTeX or XeLaTeX.}
\fi
\usepackage{fontspec}
\usepackage{amsmath}
\usepackage{unicode-math}
\setmainfont{STIXTwoText}[Extension=.otf,UprightFont=*-Regular,ItalicFont=*-Italic,BoldFont=*-Bold,BoldItalicFont=*-BoldItalic]
\setmathfont{STIXTwoMath-Regular.otf}

\usepackage[main=english]{babel}
\usepackage[a4paper,inner=27mm,outer=23mm,top=25mm,bottom=27mm,headheight=14pt]{geometry}
\usepackage{graphicx}
\usepackage{booktabs}
\usepackage{longtable}
\usepackage{array}
\usepackage{needspace}
\usepackage{microtype}
\usepackage{parskip}
\usepackage{fancyhdr}
\usepackage{xcolor}
\usepackage[hidelinks,unicode,bookmarksopen=true,bookmarksnumbered=false]{hyperref}

% LuaTeX omits volatile producer/date/trailer fields. XeTeX/xdvipdfmx is
% assigned a fixed trailer ID and must be launched with SOURCE_DATE_EPOCH=0.
\ifLuaTeX
  \pdfvariable suppressoptionalinfo 767\relax
\fi
\ifXeTeX
  \special{pdf:trailerid [<D019D019D019D019D019D019D019D019><D019D019D019D019D019D019D019D019>]}
\fi

\hypersetup{
  pdftitle={Hodge Theory, III - Canonical English Edition},
  pdfauthor={},
  pdfsubject={D019 canonical scholarly edition},
  pdfkeywords={D019, Hodge theory, source edition, translation, critical apparatus}
}

\graphicspath{{../assets/}}
\setlength{\parindent}{0pt}
\setlength{\parskip}{0.55\baselineskip}
\setlength{\emergencystretch}{3em}
\raggedbottom

\pagestyle{fancy}
\fancyhf{}
\fancyhead[LE,RO]{\small Canonical English Edition}
\fancyhead[RE,LO]{\small D019}
\fancyfoot[C]{\thepage}

\newcommand{\DArticleAnchor}[6]{%
  \clearpage
  \phantomsection
  \hypertarget{article-#1}{}%
  \pdfbookmark[1]{Article page #2 - printed #5}{bookmark-article-#1}%
  \section*{Article page #2}%
  \addcontentsline{toc}{section}{Article page #2}%
  \noindent\small\textsc{Page anchor}\quad
  authority PDF #3; collected-volume #4; printed #5; prompt #6.
  \normalsize\par\medskip
}

\newcommand{\DApparatusAnchor}[6]{%
  \Needspace{7\baselineskip}%
  \phantomsection
  \hypertarget{article-#1}{}%
  \pdfbookmark[1]{Article page #2 - printed #5}{bookmark-article-#1}%
  \subsection*{Article page #2}%
  \addcontentsline{toc}{subsection}{Article page #2}%
  \noindent\small\textsc{Page anchor}\quad
  authority PDF #3; collected-volume #4; printed #5; prompt #6.
  \normalsize\par
}

\newcommand{\DRecordMeta}[1]{%
  {\footnotesize\color{black!70}#1\par}\medskip
}

\newcommand{\DFormula}[1]{%
  \par\smallskip\noindent\hspace*{0pt}\(\displaystyle #1\)\par\smallskip
}

\newcommand{\DAsset}[2]{%
  \par\medskip
  \begin{minipage}{\linewidth}
    \centering
    \includegraphics[width=\linewidth,height=.44\textheight,keepaspectratio]{#1_presentation.png}\par
    \vspace{2pt}{\footnotesize\textbf{#1.} #2\par}
  \end{minipage}
  \par\medskip
}

\newcommand{\DAssetReference}[2]{%
  \par\smallskip{\footnotesize\textit{#1 (same local fallback):} #2\par}\smallskip
}

\newcommand{\DAssetList}[1]{%
  {\footnotesize\textsc{Fallbacks:} #1\par}
}

\newcommand{\DSetminus}{\ensuremath{\mathbin{\backslash}}}
\newcommand{\DLBrace}{\ensuremath{\lbrace}}
\newcommand{\DRBrace}{\ensuremath{\rbrace}}
\newcommand{\DArrow}[1]{\ensuremath{\xrightarrow{#1}}}
\newcommand{\DLeftArrow}[1]{\ensuremath{\xleftarrow{#1}}}

\title{Hodge Theory, III - Canonical English Edition}
\author{}
\date{}

\begin{document}
\maketitle
\thispagestyle{empty}

\begin{quote}\small This standalone source is generated deterministically from the preserved English reading records (SHA-256 8AEB2BC7\allowbreak{}00BAD5AC\allowbreak{}71193EA1\allowbreak{}71F868A9\allowbreak{}B5DE6D84\allowbreak{}C586B82C\allowbreak{}0707F0AA\allowbreak{}2A5B9D5C\allowbreak{}). Exact, fail-closed repair IDs applied to this layer: D019-RM-001, D019-RM-002, D019-RM-003, D019-RM-004, D019-RM-005, D019-RM-006, D019-RM-007, D019-RM-008, D019-RM-009, D019-RM-011, D019-RM-013, D019-RM-015, D019-RM-017, D019-RM-019, D019-RM-021, D019-RM-022, D019-RM-024, D019-RM-026, D019-RM-030, D019-RM-032, D019-RM-047, D019-RM-050, D019-RM-053, D019-RM-058. Every numbered entry preserves the corpus four-way page anchor. Complex diagrams are supplied only from local, byte-verified presentation fallbacks. The asset-override receipt records canonical-only expansions of clipped authority crops; inherited record asset paths retain their frozen provenance. No external URL, script, or network resource is used. \end{quote}

\begin{quote}\small Reading convention. Mathematical notation is emitted in Unicode math mode with STIX Two Math. Labelled arrows are executable TeX, while complex arrow topology remains controlled by the adjacent local fallback images. \end{quote}

% Article-level navigation is supplied by the complete PDF bookmark tree.
% The source work's own sommaire remains in article page 1.

\DArticleAnchor{0001}{1}{2}{668}{5}{P01}

\begingroup\small\setlength{\parskip}{0.1\baselineskip}

HODGE THEORY, III\par
by PIERRE DELIGNE\par
\par\medskip
CONTENTS\par
\par\medskip
INTRODUCTION\dotfill 6\par
TERMINOLOGY AND NOTATION\dotfill 7\par
5. Cohomological descent\dotfill 8\par
\hspace*{1em}1. Simplicial topological spaces\dotfill 8\par
\hspace*{1em}2. Cohomology of simplicial topological spaces\dotfill 12\par
\hspace*{1em}3. Cohomological descent\dotfill 14\par
6. Examples of simplicial topological spaces\dotfill 16\par
\hspace*{1em}1. Classifying spaces\dotfill 16\par
\hspace*{1em}2. Construction of hypercoverings\dotfill 19\par
\hspace*{1em}3. Relative cohomology\dotfill 21\par
\hspace*{1em}4. Multisimplicial spaces\dotfill 22\par
7. Supplements to § 1\dotfill 23\par
\hspace*{1em}1. Filtered derived category\dotfill 23\par
\hspace*{1em}2. Supplements to the two-filtrations lemma\dotfill 25\par
8. Hodge theory of algebraic spaces\dotfill 28\par
\hspace*{1em}1. Hodge complexes\dotfill 28\par
\hspace*{1em}2. Separated algebraic spaces\dotfill 38\par
\hspace*{1em}3. Hodge theory of simplicial schemes\dotfill 41\par
9. Examples and applications\dotfill 45\par
\hspace*{1em}1. Cohomology of groups and classifying spaces\dotfill 45\par
\hspace*{1em}2. Hodge theory of smooth hypersurfaces, following Griffiths\dotfill 48\par
\hspace*{1em}3. Construction of complexes of first-order differential operators\dotfill 50\par
10. Hodge theory in level \(\leq \) 1\dotfill 53\par
\hspace*{1em}1. 1-motives\dotfill 53\par
\hspace*{1em}2. 1-motives and biextensions\dotfill 60\par
\hspace*{1em}3. Algebraic interpretation of mixed \(H^{1}\): the case of curves\dotfill 69\par
\hspace*{1em}4. Translation of a theorem of Picard\dotfill 75\par
BIBLIOGRAPHY\dotfill 77\par

\endgroup

\DArticleAnchor{0002}{2}{3}{669}{6}{P01}

Introduction\par
\par\medskip
This article follows [4] and [5], referred to as I and II. The numbering of its sections continues that of II (which contains §§ 1 through 4). To refer to a result of II, we simply give its number. For example: (3.2.17).\par
\par\medskip
In I, we set out the yoga underlying II and the present article III. Nevertheless, II and III are logically independent of I. They do not contain all the results announced in I.\par
\par\medskip
In II, we set out the Hodge theory of nonsingular algebraic varieties (not necessarily complete). Here we treat the case of arbitrary singularities. Beginning with § 7, essential use will be made of the results of §§ 1 through 3 of II.\par
\par\medskip
In the case of a complete singular algebraic variety X, the fundamental idea is to use resolution of singularities to “replace” X by a simplicial system of smooth projective schemes \(X_{\bullet }\).\par
\par\medskip
\DAsset{P0002-A01}{simplicial diagram}
\par\medskip
(or, as we shall say, by a smooth projective simplicial scheme). The results of §§ 5 and 6 make it possible to find such simplicial schemes \(X_{\bullet }\) having, in a suitable sense, the same cohomology as X. One may then “express,” by means of a spectral sequence, the cohomology of X in terms of that of the \(X_{n}\). In accordance with the general principles of I, one may then use classical Hodge theory on each \(X_{n}\) and thereby deduce a mixed Hodge structure on the cohomology of X.\par
\par\medskip
In the case of an arbitrary algebraic variety X, one “replaces” X by a smooth simplicial scheme \(X_{\bullet }\), compactified by a smooth projective simplicial scheme \(\overline{X}_{\bullet }\). One further arranges that \(\overline{X}_{n}\) \(−\) \(X_{n}\) is a normal-crossings divisor, a union of smooth divisors \(D_{n,i}\). One then “expresses,” by means of a spectral sequence, the cohomology of X in terms of the cohomology of the p-fold intersections of the \(D_{n,i}\) (n \(\geq \) 0, p \(\geq \) 0), and deduces a mixed Hodge structure on \(H^{\bullet }(X)\).\par
\par\medskip
Sections 5 and 6 constitute a summary (without proofs) of the theory of “cohomological descent.” This theory, in the more general framework of topoi, is set out by B. Saint-Donat in S.G.A.4, V bis.\par
\par\medskip
In § 7, no. 1, certain previously obtained results are recast in the language of filtered derived categories. No. 2 contains a proof of the two-filtrations lemma simpler than the one in § 1, no. 3.\par
\par\medskip
The forbidding no. 1 of § 8 is the heart of this work. The essential result is Theorem (8.1.15). In its proof, in order to apply the two-\par

\DArticleAnchor{0003}{3}{4}{670}{7}{P01}

filtrations lemma, one uses the fact that every morphism of mixed Hodge structures is strictly compatible with the filtrations W and F (2.3.5). In no. 2, the mixed Hodge structure on \(H^{n}(X\), \(\mathbb{Z})\) is defined (X a separated scheme). It is shown that the Hodge numbers \(h^{pq}\) of \(H^{n}(X\), \(\mathbb{Z})\) can be nonzero only for (p, q) \(\in \) [0, n] \(\times \) [0, n]. For n \(>\) dim X, there is a more precise result (8.2.4). If X is smooth (respectively, complete), they can be nonzero only if, in addition, p \(+\) q \(\geq \) n (respectively, p \(+\) q \(\leq \) n). If X is a complete closed subscheme with complement U in a complete smooth scheme Z of dimension n, then, following N. Katz, one may interpret this “duality” between the complete and smooth cases as arising from “Alexander duality”\par
\par\medskip
\DFormula{…\,H^{i}(Z,\,\mathbb{Q})\,\to \,H^{i}(X,\,\mathbb{Q})\,\to \,(H^{2n−i−1}(U,\,\mathbb{Q}))^{\ast }\,\to \,H^{i+1}(Z,\,\mathbb{Q})\,…}
\par\medskip
From the functoriality properties of the theory constructed, one deduces the useful Corollaries (8.2.5) through (8.2.8), which are stated in terms independent of any Hodge theory.\par
\par\medskip
In § 8, no. 3, the cohomology of every simplicial scheme \(X_{\bullet }\) is endowed with a mixed Hodge structure. This generality is not illusory.\par
\par\medskip
a) Relative cohomology spaces can be interpreted as the cohomology of suitable simplicial schemes.\par
\par\medskip
b) Let \(B_{G}\) be the classifying space of the Lie group underlying an algebraic group G. The cohomology of \(B_{G}\) can be interpreted as that of a suitable simplicial scheme.\par
\par\medskip
In § 9, no. 1, after computing the mixed Hodge structure on the cohomology of \(B_{G}\) (for G linear), one deduces that of G. As a corollary, one finds that if a linear group G acts on a complete nonsingular variety X, then the map\par
\par\medskip
\DFormula{H^{\bullet }(X,\,\mathbb{Q})\,\to \,H^{\bullet }(X,\,\mathbb{Q})\,\otimes \,H^{\bullet }(G,\,\mathbb{Q})}
\par\medskip
\DFormula{factors\,through\,H^{\bullet }(X,\,\mathbb{Q})\,\otimes \,H^{0}(G,\,\mathbb{Q})\,\subset \,H^{\bullet }(X,\,\mathbb{Q})\,\otimes \,H^{\bullet }(G,\,\mathbb{Q}).}
\par\medskip
In § 10, mixed Hodge structures purely of type \(\lbrace (−1,\,−1),\,(−1,\,0),\,(0,\,−1),\,(0,\,0)\rbrace \) (III.0.5) are interpreted in terms of abelian schemes, and \(H^{1}\) of curves is considered in detail.\par
\par\medskip
The theory developed thus far is an absolute theory (the functors \(Rf_{\ast }\) are not studied), and concerns only constant coefficients. I conjecture that if \(\mathcal{H}\) is a polarizable variation of Hodge structures (in the sense of Griffiths [6]) on a scheme X, then the cohomology of X with coefficients in the local system \(\mathcal{H}\) carries a natural mixed Hodge structure. I can prove this only when X is complete.\par
\par\medskip
Terminology and notation\par
\par\medskip
(III.0.1) Let u : X \(\to \) Y be a continuous map between topological spaces. We shall say that u is proper if it is proper in the sense of Bourbaki (i.e. universally closed) and, in addition, separated (i.e. if the diagonal of X \(\times _{Y}\) X is closed).\par
\DAsset{P0003-A01}{Alexander-duality exact sequence}
\DAsset{P0003-A02}{cohomology map and factorization through \(H^{0}(G\), \(\mathbb{Q})\)}

\DArticleAnchor{0004}{4}{5}{671}{8}{P01}

(III.0.2) Following Gabriel and Zisman in this, we shall say simplicial where formerly one said semisimplicial.\par
\par\medskip
(III.0.3) We shall denote by A a Noetherian subring of \(\mathbb{R}\) such that A \(\otimes \) \(\mathbb{Q}\) is a field. The useful cases are A \(=\) \(\mathbb{Z}\), \(\mathbb{Q}\), or \(\mathbb{R}\).\par
\par\medskip
(III.0.4) An A-mixed Hodge structure consists of a finite-type A-module \(H_{A}\), a finite increasing filtration W on the A \(\otimes \) \(\mathbb{Q}-vector\) space \(H_{A\otimes \mathbb{Q}}\) \(=\) \(H_{A}\) \(\otimes \) \(\mathbb{Q}\), and a finite decreasing filtration F on the \(\mathbb{C}-vector\) space \(H_{\mathbb{C}}\) \(=\) \(H_{A}\) \(\otimes _{A}\) \(\mathbb{C}\). It is required that the \((\operatorname{Gr}_{n}^{W}(H_{A\otimes \mathbb{Q}})\), \(\operatorname{Gr}_{n}^{W}(F))\) be A \(\otimes \) \(\mathbb{Q}-Hodge\) structures. For A \(=\) \(\mathbb{Z}\) (respectively, \(\mathbb{Q})\), one recovers (2.3.1) (respectively, 2.3.8); the results of (2.3) carry over unchanged.\par
\par\medskip
(III.0.5) Let \(\mathcal{E}\) be a subset of \(\mathbb{Z}\) \(\times \) \(\mathbb{Z}\). A mixed Hodge structure H is of type \(\mathcal{E}\) if its Hodge numbers \(h^{pq}\) vanish for (p, q) \(\notin \) \(\mathcal{E}\).\par
\par\medskip
(III.0.6) We shall henceforth write \(\Omega _{X}^{\bullet }(\operatorname{log}\) D) for what we denoted by \(\Omega _{X}^{\ast }⟨D⟩\) (3.1.2).\par
\par\medskip
(III.0.7) Unless explicitly stated otherwise, scheme means “scheme of finite type over \(\mathbb{C}\),” and a sheaf on a scheme X is a sheaf on the topological space underlying \(X_{\mathrm{an}}\).\par
\par\medskip
5. COHOMOLOGICAL DESCENT\par
\par\medskip
5.1. Simplicial topological spaces.\par
\par\medskip
This subsection begins with reminders for which one may refer to the Strasbourg homotopy seminar, 63/64.\par
\par\medskip
(5.1.1) We shall use the following notation.\par
\par\medskip
\(𝒜^{o}\) \(=\) the category opposite to a category \(𝒜\).\par
\(\operatorname{Hom}(𝒜\), \(\mathcal{B})\) \(=\) the category of functors from \(𝒜\) to \(\mathcal{B}\).\par
\par\medskip
Let n and k be integers \(\geq \) \(−1\).\par
\par\medskip
\(\Delta _{n}\) \(=\) the finite totally ordered set [0, n].\par
\(\delta _{i}\) : \(\Delta _{n}\) \(\to \) \(\Delta _{n+1}\) \(=\) the increasing injection such that i \(\notin \) \(\delta _{i}(\Delta _{n})\) (0 \(\leq \) i \(\leq \) n \(+\) 1).\par
\(s_{i}\) : \(\Delta _{n+1}\) \(\to \) \(\Delta _{n}\) \(=\) the increasing surjection such that \(s_{i}(i)\) \(=\) \(s_{i}(i\) \(+\) 1) (0 \(\leq \) i \(\leq \) n).\par
\(\varepsilon \) : \(\Delta _{−1}\) \(\to \) \(\Delta _{n}\) \(=\) the unique map from \(\Delta _{−1}\) to \(\Delta _{n}\).\par
\((\Delta ^{+})\) \(=\) the category whose objects are the \(\Delta _{n}\) (n \(\geq \) \(−1)\), and whose arrows are the increasing maps between the \(\Delta _{n}\).\par
\((\Delta )\) \(=\) the full subcategory of \((\Delta ^{+})\) whose objects are the \(\Delta _{n}\) (n \(\geq \) 0).\par
\((\Delta ^{+})_{k}\) \(=\) the full subcategory of \((\Delta ^{+})\) whose objects are the \(\Delta _{n}\) (k \(\geq \) n).\par
\((\Delta )_{k}\) \(=\) the full subcategory of \((\Delta ^{+})\) whose objects are the \(\Delta _{n}\) (k \(\geq \) n \(\geq \) 0).\par
\DAsset{P0004-A01}{mixed-Hodge-structure notation, (III.0.2)-(III.0.7)}
\DAsset{P0004-A02}{simplex-category notation, (5.1.1)}

\DArticleAnchor{0005}{5}{6}{672}{9}{P01}

For every category \(\mathcal{C}\), a simplicial object (respectively, a k-truncated simplicial object) of \(\mathcal{C}\) is an object of \(\operatorname{Hom}((\Delta )^{o}\), \(\mathcal{C})\) (respectively, \(\operatorname{Hom}((\Delta )_{k}^{o}\), \(\mathcal{C}))\). Likewise, a cosimplicial object (respectively, …) of \(\mathcal{C}\) is an object of \(\operatorname{Hom}((\Delta )\), \(\mathcal{C})\) (respectively, …).\par
\par\medskip
For k \(\geq \) \(−1\), the skeleton functor \(\operatorname{sq}_{k}\) is the restriction functor\par
\par\medskip
\DFormula{\operatorname{sq}_{k}\,:\,\operatorname{Hom}((\Delta )^{o},\,\mathcal{C})\,\to \,\operatorname{Hom}((\Delta )_{k}^{o},\,\mathcal{C}).}
\par\medskip
The coskeleton functor \(\operatorname{cosq}_{k}\) is the right adjoint to \(\operatorname{sq}_{k}\)\par
\par\medskip
\DFormula{\operatorname{cosq}_{k}\,:\,\operatorname{Hom}((\Delta )_{k}^{o},\,\mathcal{C})\,\to \,\operatorname{Hom}((\Delta )^{o},\,\mathcal{C}).}
\par\medskip
Let Y be a simplicial object of \(\mathcal{C}\). The functor\par
\par\medskip
\DFormula{\operatorname{sq}_{k}\,:\,\operatorname{Hom}((\Delta )^{o},\,\mathcal{C})/Y\,\to \,\operatorname{Hom}((\Delta )_{k}^{o},\,\mathcal{C})/\operatorname{sq}_{k}(Y)}
\par\medskip
is also called the skeleton functor. Its right adjoint is denoted by \(\operatorname{cosq}_{k}^{Y}\) (the coskeleton relative to Y)\par
\par\medskip
\DFormula{\operatorname{cosq}_{k}^{Y}\,:\,\operatorname{Hom}((\Delta )_{k}^{o},\,\mathcal{C})/\operatorname{sq}_{k}\,Y\,\to \,\operatorname{Hom}((\Delta )^{o},\,\mathcal{C})/Y.}
\par\medskip
The coskeleton functors are defined if finite projective limits exist in \(\mathcal{C}\). The functors \(\operatorname{cosq}_{k}^{Y}\) exist as soon as fiber products exist. One has\par
\par\medskip
\DFormula{\operatorname{cosq}_{k}^{Y}(X)\,=\,\operatorname{cosq}_{k}(X)\,\times _{\operatorname{cosq}_{k}\,\operatorname{sq}_{k}(Y)}\,Y.}
\par\medskip
(5.1.2) If \(X_{\bullet }\) : \((\Delta )^{o}\) \(\to \) \(\mathcal{C}\) is a simplicial object of \(\mathcal{C}\), set \(X_{n}\) \(=\) \(X_{\bullet }(\Delta _{n})\), \(\delta _{i}\) \(=\) \(X_{\bullet }(\delta _{i}\) : \(\Delta _{n}\) \(\to \) \(\Delta _{n+1})\) : \(X_{n+1}\) \(\to \) \(X_{n}\), and \(s_{i}\) \(=\) \(X_{\bullet }(s_{i}\) : \(\Delta _{n+1}\) \(\to \) \(\Delta _{n})\) : \(X_{n}\) \(\to \) \(X_{n+1}\).\par
\par\medskip
\DAsset{P0005-A02}{face and degeneracy diagram}
\par\medskip
(5.1.3) Let S \(\in \) Ob \(\mathcal{C}\). The constant simplicial object \(S_{\bullet }\) is the simplicial object for which \(S_{n}\) \(=\) S and \(\delta _{i}\) \(=\) \(s_{i}\) \(=\) \(Id_{S}\). A simplicial object (respectively, a k-truncated simplicial object) of \(\mathcal{C}\) augmented to S is a morphism a : \(X_{\bullet }\) \(\to \) \(S_{\bullet }\) (respectively, a : \(X_{\bullet }\) \(\to \) \(\operatorname{sq}_{k}(S_{\bullet }))\). We identify simplicial objects (respectively, k-truncated simplicial objects) of \(\mathcal{C}\) augmented to S, a : \(X_{\bullet }\) \(\to \) \(S_{\bullet }\) (respectively, …), with objects \(X_{\bullet }^{+}\) of \(\operatorname{Hom}((\Delta ^{+})^{o}\), \(\mathcal{C})\) (respectively, \(\operatorname{Hom}((\Delta ^{+})_{k}^{o}\), \(\mathcal{C}))\) such that \(X_{\bullet }^{+}(\Delta _{−1})\) \(=\) S. One has \(a_{n}\) \(=\) \(X_{\bullet }^{+}(\varepsilon \) : \(\Delta _{−1}\) \(\to \) \(\Delta _{n})\). These objects will be denoted by notation of the form a : \(X_{\bullet }\) \(\to \) S. Relative coskeletons will be used mainly in this setting and will be denoted by \(\operatorname{cosq}_{k}^{S}\), or simply \(\operatorname{cosq}_{k}\).\par
\par\medskip
(5.1.4) The coskeleton of the 0-truncated simplicial object \(a_{0}\) : X \(\to \) S is the simplicial object augmented to S whose components are the Cartesian powers of X in \(\mathcal{C}/S\)\par
\par\medskip
\DFormula{\operatorname{cosq}_{0}(X\,\to \,S)\,=\,(((X/S)^{\Delta _{n}})_{n\geq 0}\,\to \,S).}
\par\medskip
(5.1.5) Let u : X \(\to \) Y be a continuous map between topological spaces, let F be a sheaf on X, and let G be a sheaf on Y. The set \(\operatorname{Hom}_{u}(G\), F) of u-morphisms from G to F is the set \(\operatorname{Hom}(u^{\ast }G\), F) \(\simeq \) Hom(G, \(u_{\ast }F)\).\par
\DAsset{P0005-A01}{skeleton and coskeleton functor formulas}
\DAsset{P0005-A03}{relative coskeleton and the adjunction for u-morphisms, (5.1.4)-(5.1.5)}

\DArticleAnchor{0006}{6}{7}{673}{10}{P01}

To give a u-morphism f from G to F amounts to giving, for every pair of open sets U \(\subset \) X and V \(\subset \) Y such that u(U) \(\subset \) V, a map \(f_{UV}\) : G(V) \(\to \) F(U); these maps must satisfy the condition\par
\par\medskip
(*) For U′ \(\subset \) U \(\subset \) X, V′ \(\subset \) V \(\subset \) Y, u(U) \(\subset \) V, and u(U′) \(\subset \) V′, the diagram\par
\par\medskip
\DAsset{P0006-A01}{commutative square}
\par\medskip
is commutative.\par
\par\medskip
(5.1.6) A simplicial topological space is a simplicial object in the category whose objects are topological spaces and whose morphisms are continuous maps.\par
\par\medskip
A sheaf \(F^{\bullet }\) on a simplicial topological space \(X_{\bullet }\) consists of\par
\par\medskip
a) a family \(F^{n}\) of sheaves on the \(X_{n}\);\par
\par\medskip
\DFormula{b)\,for\,f\,:\,\Delta _{n}\,\to \,\Delta _{m},\,\mathrm{an}\,X_{\bullet }(f)-morphism\,F^{\bullet }(f)\,from\,F^{n}\,to\,F^{m}.}
\par\medskip
It is further required that \(F^{\bullet }(f\) \(∘\) g) \(=\) \(F^{\bullet }(f)\) \(∘\) \(F^{\bullet }(g)\).\par
\par\medskip
A morphism u from \(F^{\bullet }\) to \(G^{\bullet }\) is a family of morphisms \(u^{n}\) : \(F^{n}\) \(\to \) \(G^{n}\) such that, for f \(\in \) \(\operatorname{Hom}_{\Delta }(\Delta _{n}\), \(\Delta _{m})\), one has \(u^{m}\) \(F^{\bullet }(f)\) \(=\) \(G^{\bullet }(f)\) \(u^{n}\).\par
\par\medskip
(5.1.7) For U open in \(X_{n}\), set \(F^{\bullet }(U)\) \(=\) \(F^{n}(U)\). For f : \(\Delta _{n}\) \(\to \) \(\Delta _{m}\), U \(\subset \) \(X_{n}\), and V in \(X_{m}\) such that \(X_{\bullet }(f)(V)\) \(\subset \) U, let \(F^{\bullet }(f\), V, U) : \(F^{\bullet }(U)\) \(\to \) \(F^{\bullet }(V)\) be the map induced by \(F^{\bullet }(f)\). It is immediate that the system of sets \(F^{\bullet }(U)\) (indexed by n and U \(\subset \) \(X_{n})\) and the maps \(F^{\bullet }(f\), V, U) determines \(F^{\bullet }\) uniquely (cf. (5.1.5)). For such a system to arise from a sheaf on \(X_{\bullet }\), it is necessary and sufficient that:\par
\par\medskip
\DFormula{a)\,F^{\bullet }(fg,\,U,\,W)\,=\,F^{\bullet }(f,\,U,\,V)\,F^{\bullet }(g,\,V,\,W)\,whenever\,this\,makes\,sense;}
\par\medskip
b) for every n, the \(F^{\bullet }(U)\), for U \(\subset \) \(X_{n}\), form a sheaf on \(X_{n}\).\par
\par\medskip
(5.1.8) This shows that sheaves on \(X_{\bullet }\) can be interpreted as sheaves on a suitable site and, in particular, form a topos \((X_{\bullet })^{\sim }\). We shall freely use, for sheaves on \(X_{\bullet }\), the terminology used for sheaves on a site. Thus a sheaf of abelian groups (respectively, a sheaf of rings, …) \(F^{\bullet }\) is a system of sheaves of abelian groups (respectively, sheaves of rings, …) \(F^{n}\) on the \(X_{n}\), together with morphisms \(F^{\bullet }(f)\).\par
\par\medskip
(5.1.9) Examples. (I) Let \(X_{\bullet }\) be a simplicial analytic space. The structure sheaves \(\mathcal{O}_{X_{n}}\) form a sheaf of rings on \(X_{\bullet }\).\par
\par\medskip
(II) Let \(X_{\bullet }\) be a simplicial analytic space augmented to S. The sheaves \(\Omega ^{1}_{X_{n}/S}\) form a sheaf of \(\mathcal{O}-modules\) on \(X_{\bullet }\). Its i-th exterior power is the sheaf\par
\DAsset{P0006-A02}{structural equations for sheaves on a simplicial space, (5.1.6)-(5.1.8)}

\DArticleAnchor{0007}{7}{8}{674}{11}{P02}

of \(\mathcal{O}-modules\) \((\Omega ^{i}_{X_{n}/S})_{n\geq 0}\), denoted \(\Omega ^{i}_{X_{\bullet }/S}\). The De Rham complexes \((\Omega ^{\bullet }_{X_{n}/S})_{n\geq 0}\) form a complex of sheaves on \(X_{\bullet }\).\par
\par\medskip
(III) Let \(F^{\bullet }\) be a sheaf of abelian groups on \(X_{\bullet }\). The canonical flasque resolutions of Godement \(\mathcal{C}^{\bullet }(F^{n})\) form a complex of sheaves on \(X_{\bullet }\) and a resolution of \(F^{\bullet }\).\par
\par\medskip
(IV) Let S be a topological space. Sheaves \(F^{\bullet }\) on the constant simplicial space \(S_{\bullet }\) identify with cosimplicial sheaves on S. In particular, an abelian sheaf \(F^{\bullet }\) on \(S_{\bullet }\) defines a differential complex \((F^{n}\), d \(=\) \(\sum _{i}\) \((−1)^{i}\) \(\delta _{i})\). A complex K of abelian sheaves on \(S_{\bullet }\) defines a double complex \((K^{n,m})\), also denoted K (m \(=\) cosimplicial degree), with associated simple complex sK:\par
\par\medskip
\DFormula{(5.1.9.1)\,(sK)^{n}\,=\,⨁_{p+q=n}\,K^{pq}}
\par\medskip
and\par
\par\medskip
\DFormula{(5.1.9.2)\,d(x^{pq})\,=\,d_{K}(x^{pq})\,+\,(−1)^{p}\,\sum _{i}\,(−1)^{i}\,\delta _{i}\,x^{pq}.}
\par\medskip
We denote by L the second filtration of sK:\par
\par\medskip
\DFormula{(5.1.9.3)\,L^{r}(sK)\,=\,⨁_{q\geq r}\,K^{pq}.}
\par\medskip
(5.1.10) Let u : \(X_{\bullet }\) \(\to \) \(Y_{\bullet }\) be a morphism of simplicial topological spaces, with components \(u_{n}\) : \(X_{n}\) \(\to \) \(Y_{n}\). If G (respectively, F) is a sheaf on \(Y_{\bullet }\) (respectively, on \(X_{\bullet })\), then \((u_{n}^{\ast }G^{n})_{n\geq 0}\) (respectively, \((u_{n\ast }F^{n})_{n\geq 0})\) is a sheaf on \(X_{\bullet }\) (respectively, on \(Y_{\bullet })\); it is denoted \(u^{\ast }G\) (respectively, \(u_{\ast }F)\). The functors \(u_{\ast }\) and \(u^{\ast }\) are adjoint; they are the direct-image and inverse-image morphisms of a morphism of topoi u : \((X_{\bullet })^{\sim }\) \(\to \) \((Y_{\bullet })^{\sim }\).\par
\par\medskip
(5.1.11) Let a : \(X_{\bullet }\) \(\to \) S be a simplicial topological space augmented to S. If F is a sheaf on S, \(a^{\ast }F\) \(=\) \((a_{n}^{\ast }F)_{n\geq 0}\) “is” a sheaf on \(X_{\bullet }\). The functor \(a^{\ast }\) has a right adjoint\par
\par\medskip
\DFormula{(5.1.11.1)\,a_{\ast }\,:\,F^{\bullet }\,↦\,\operatorname{Ker}(a_{0\ast }F^{0}\,\to \to [\delta _{0}^{\ast };\,\delta _{1}^{\ast }]\,a_{1\ast }F^{1}).}
\par\medskip
The functors \(a_{\ast }\) and \(a^{\ast }\) are the direct-image and inverse-image functors of a morphism of topoi a : \((X_{\bullet })^{\sim }\) \(\to \) \((S)^{\sim }\).\par
\par\medskip
(5.1.12) Let \(S_{\bullet }\) be the constant simplicial space associated with S, and let \(a_{\bullet }\) : \(X_{\bullet }\) \(\to \) \(S_{\bullet }\) be the morphism defined by a. For an abelian sheaf \(F^{\bullet }\) on \(X_{\bullet }\), \(a_{\bullet \ast }(F^{\bullet })\) is identified with a cosimplicial sheaf on S ((5.1.9) (IV)). For K a complex of abelian sheaves on \(X_{\bullet }\), if \(K^{pq}\) is the degree-p component of the restriction of K to \(X_{q}\), the components of the simple complex associated with the double complex defined by \(a_{\bullet \ast }K\) are\par
\par\medskip
\DFormula{(5.1.12.1)\,(s\,a_{\bullet \ast }K)^{n}\,=\,⨁_{p+q=n}\,a_{q\ast }K^{pq}.}
\DAsset{P0007-A01}{simple-complex formulas and the L filtration, (5.1.9.1)-(5.1.9.3)}
\DAsset{P0007-A02}{direct/inverse image functors and formulas (5.1.11.1)-(5.1.12.1)}

\DArticleAnchor{0008}{8}{9}{675}{12}{P02}

The spectral sequence defined by the filtration L ((5.1.9.3)) of s \(a_{\bullet \ast }K\) is\par
\par\medskip
\DFormula{(5.1.12.2)\,E_{1}^{pq}\,=\,H^{q}(a_{p\ast }(K\,\mid \,X^{p}))\,\Rightarrow \,H^{p+q}(s\,a_{\bullet \ast }K).}
\par\medskip
(5.1.13) Take S \(=\) \(P^{t}\) (the topological space consisting of one point). For a sheaf \(F^{\bullet }\) on \(X_{\bullet }\), the \(\Gamma (X_{n},F^{n})\) are the components of a cosimplicial set \(\Gamma ^{\bullet }(X_{\bullet },F^{\bullet })\). The functor \(\Gamma \) (global sections) is the functor\par
\par\medskip
\DFormula{(5.1.13.1)\,\Gamma \,:\,F^{\bullet }\,↦\,\operatorname{Ker}(\Gamma (X_{0},F^{0})\,\rightrightarrows \,\Gamma (X_{1},F^{1})).}
\par\medskip
If K is a complex of abelian sheaves, \(\Gamma ^{\bullet }(X_{\bullet },K)\) denotes the differential complex whose components are the cosimplicial abelian groups \(\Gamma ^{\bullet }(X_{\bullet },K^{n})\), and \(s\Gamma ^{\bullet }(X_{\bullet },K)\) denotes the associated simple complex.\par
\par\medskip
5.2. Cohomology of simplicial topological spaces.\par
\par\medskip
(5.2.1) To every simplicial topological space \(X_{\bullet }\) there is attached an ordinary topological space \(\mid X_{\bullet }\mid \), called its geometric realization (see below). In the special case where the \(X_{n}\) are discrete, one recovers the usual notion of the geometric realization of a simplicial set.\par
\par\medskip
In [14], G. Segal defines the cohomology of \(X_{\bullet }\) with values in an abelian group A to be \(H^{\bullet }(\mid X_{\bullet }\mid ,A)\). Under suitable hypotheses, the filtration of \(\mid X_{\bullet }\mid \) by its successive skeleta gives a spectral sequence.\par
\par\medskip
\DFormula{(5.2.1.1)\,E_{1}^{pq}\,=\,H^{q}(X_{p},A)\,\Rightarrow \,H^{p+q}(X_{\bullet },A).}
\par\medskip
One advantage of this definition is that it applies equally well, for example, to the definition of groups \(K(X_{\bullet })\). We shall adopt another definition, better suited to sheaf-theoretic techniques.\par
\par\medskip
Geometric realization. — For an integer n \(\geq \) 0, let \(\mid \Delta _{n}\mid \) denote the simplex in \(\mathbb{R}^{\Delta _{n}}\) whose vertices are the set of basis vectors. We identify \(\Delta _{n}\) with the set of vertices of \(\mid \Delta _{n}\mid \). Every function f : \(\Delta _{n}\) \(\to \) \(\Delta _{m}\) extends linearly to \(\mid f\mid \) : \(\mid \Delta _{n}\mid \) \(\to \) \(\mid \Delta _{m}\mid \).\par
\par\medskip
Let \(X_{\bullet }\) be a simplicial topological space. Let Y \(=\) \(⨿_{n\geq 0}(X_{n}\) \(\times \) \(\mid \Delta _{n}\mid )\). Let R be the finest equivalence relation on Y for which, for every f : \(\Delta _{n}\) \(\to \) \(\Delta _{m}\) in \((\Delta )\), \(x_{m}\) \(\in \) \(X_{m}\), and a \(\in \) \(\mid \Delta _{n}\mid \), one has\par
\par\medskip
\DFormula{(x_{m},\,\mid f\mid (a))\,\equiv \,(f(x_{m}),\,a)\,\,\,(\operatorname{mod}\,R).}
\par\medskip
The geometric realization of \(X_{\bullet }\) is, by definition,\par
\par\medskip
\DFormula{\mid X_{\bullet }\mid \,=\,Y/R.}
\par\medskip
Definition (5.2.2). — Let \(X_{\bullet }\) be a simplicial topological space. The cohomology functors with coefficients in the abelian sheaf \(F^{\bullet }\) on \(X_{\bullet }\), \(H^{i}(X_{\bullet },F^{\bullet })\), are the derived functors of the functor \(\Gamma \) (5.1.13.1).\par
\par\medskip
This definition is equivalent to the following one, which is sometimes more convenient.\par
\DAsset{P0008-A01}{spectral sequence (5.1.12.2) and the global-sections functor (5.1.13.1)}
\DAsset{P0008-A02}{spectral sequence and geometric-realization formulas, (5.2.1.1)-(5.2.2)}

\DArticleAnchor{0009}{9}{10}{676}{13}{P02}

(5.2.3) Let \(F^{\bullet }\) be an abelian sheaf on \(X_{\bullet }\). One shows (for example with the aid of (5.1.9) (III)) that \(F^{\bullet }\) always admits right resolutions K such that\par
\par\medskip
\(H^{r}(X_{q},K^{p,q})\) \(=\) 0   for p,q \(\geq \) 0 and r \(>\) 0.\par
\par\medskip
If K is such a resolution, there is a canonical isomorphism\par
\par\medskip
\DFormula{(5.2.3.1)\,H^{n}(X_{\bullet },F^{\bullet })\,\simeq \,H^{n}(s\Gamma ^{\bullet }(X_{\bullet },K)).}
\par\medskip
It is easy to verify directly that the right-hand side is independent (up to unique isomorphism) of the choice of K; for what follows, one could define \(H^{\bullet }(X_{\bullet },F^{\bullet })\) by (5.2.3.1). Moreover, the spectral sequence of the complex \(s\Gamma ^{\bullet }(X_{\bullet },K)\), filtered by L ((5.1.9.3) and (5.1.12.2) for S \(=\) \(P^{t})\),\par
\par\medskip
\DFormula{(5.2.3.2)\,E_{1}^{pq}\,=\,H^{q}(X_{p},F^{p})\,\Rightarrow \,H^{p+q}(X_{\bullet },F^{\bullet }),}
\par\medskip
is independent (up to unique isomorphism) of the choice of K (cf. (5.2.1.1)).\par
\par\medskip
(5.2.4) The preceding construction should be made precise by passing to derived categories, and a relative variant should be given.\par
\par\medskip
Let a : \(X_{\bullet }\) \(\to \) S, let \(S_{\bullet }\) be the constant simplicial space defined by S, let \(a_{\bullet }\) : \(X_{\bullet }\) \(\to \) \(S_{\bullet }\) be induced by a, and let \(\varepsilon \) : \(S_{\bullet }\) \(\to \) S be the augmentation morphism. One has\par
\par\medskip
\DFormula{(5.2.4.1)\,a_{\ast }\,=\,\varepsilon _{\ast }\,a_{\bullet \ast }\,:\,(X_{\bullet })^{\sim }\,\to \,(S)^{\sim }.}
\par\medskip
This formula derives to\par
\par\medskip
\DFormula{(5.2.4.2)\,Ra_{\ast }\,=\,R\varepsilon _{\ast }\,Ra_{\bullet \ast }\,:\,D^{+}(X_{\bullet })\,\to \,D^{+}(S)}
\par\medskip
\((D^{+}(X_{\bullet })\) is the bounded-below derived category of the category of abelian sheaves on \(X_{\bullet })\). Let us compute \(Ra_{\bullet \ast }\) and \(R\varepsilon _{\ast }\).\par
\par\medskip
(5.2.5) Let u : \(X_{\bullet }\) \(\to \) \(Y_{\bullet }\) be a morphism of simplicial topological spaces. The functor \(Ru_{\ast }\) : \(D^{+}(X_{\bullet })\) \(\to \) \(D^{+}(Y_{\bullet })\) is computed “component by component”: if K is a complex of abelian sheaves on \(X_{\bullet }\), to compute \(Ru_{\ast }K\) one takes a resolution K \(\simeq \) K′ such that the components \(F^{p}\) of K′ satisfy \(R^{i}\) \(u_{p\ast }(F^{p})\) \(=\) 0 for i \(>\) 0, p \(\geq \) 0 (cf. (5.1.9) (III)); then \(Ru_{\ast }K\) \(\simeq \) \(u_{\ast }K′\).\par
\par\medskip
(5.2.6) The functor s : \(C^{+}(S_{\bullet })\) \(\to \) \(C^{+}(S)\) sends an acyclic complex to an acyclic complex. It therefore derives trivially to\par
\par\medskip
\DFormula{s\,:\,D^{+}(S_{\bullet })\,\to \,D^{+}(S).}
\par\medskip
For an injective sheaf F on \(S_{\bullet }\), the differential complex sF is a resolution of \(\varepsilon _{\ast }F\). Hence there is an isomorphism \(R\varepsilon _{\ast }\) \(\simeq \) s, and (5.2.4.2) becomes\par
\par\medskip
\DFormula{(5.2.6.1)\,Ra_{\ast }\,=\,s\,Ra_{\bullet \ast }.}
\DAsset{P0009-A01}{resolutions and spectral sequence, (5.2.3.1)-(5.2.3.2)}
\DAsset{P0009-A02}{derived-functor formulas, (5.2.4.1)-(5.2.6.1)}

\DArticleAnchor{0010}{10}{11}{677}{14}{P02}

(5.2.7) Together with (5.2.5), and specialized to the case S \(=\) \(P^{t}\), this formula proves (5.2.3.1). Concretely, it means that \(Ra_{\ast }K\) can be computed in two stages:\par
\par\medskip
a) take a resolution K \(\simeq \) K′ such that the components \(F^{p}\) of K′ satisfy \(R^{i}\) \(a_{p\ast }(F^{p})\) \(=\) 0 for i \(>\) 0. The complex \(a_{\bullet \ast }K′\) \(\in \) \(D^{+}(S_{\bullet })\) (the derived category of the category of cosimplicial abelian sheaves on S) identifies with \(Ra_{\bullet \ast }K\).\par
\par\medskip
b) \(Ra_{\ast }K\) is the simple complex s \(a_{\bullet \ast }K′\) associated with the double complex \(a_{\bullet \ast }K′\).\par
\par\medskip
The spectral sequence (5.2.3.2) generalizes to a spectral sequence deduced from (5.1.12.2):\par
\par\medskip
\DFormula{(5.2.7.1)\,E_{1}^{pq}\,=\,R^{q}\,a_{p\ast }(K\,\mid \,X_{p})\,\Rightarrow \,R^{p+q}a_{\ast }K.}
\par\medskip
5.3. Cohomological descent.\par
\par\medskip
(5.3.1) Let a : \(X_{\bullet }\) \(\to \) S be an augmented simplicial topological space. For every sheaf F on S, there is an adjunction morphism\par
\par\medskip
\DFormula{\varphi \,:\,F\,\to \,a_{\ast }a^{\ast }F.}
\par\medskip
This morphism derives to a morphism of functors from \(D^{+}(S)\) to \(D^{+}(S)\):\par
\par\medskip
\DFormula{(5.3.1.1)\,\varphi \,:\,Id\,\to \,Ra_{\ast }a^{\ast }.}
\par\medskip
Definition (5.3.2). — We say that a is of cohomological descent if, for every abelian sheaf F on S,\par
\par\medskip
\DFormula{F\,\simeq \,\operatorname{Ker}(a_{0\ast }a_{0}^{\ast }F\,\to \,a_{1\ast }a_{1}^{\ast }F)}
\par\medskip
and\par
\par\medskip
\DFormula{R^{i}\,a_{\ast }a^{\ast }F\,=\,0\,\,\,for\,i\,>\,0.}
\par\medskip
Equivalently, (5.3.1.1) is required to be an isomorphism.\par
\par\medskip
(5.3.3) If a is of cohomological descent, then, for K \(\in \) Ob \(D^{+}(S)\), the canonical map\par
\par\medskip
\DFormula{(5.3.3.1)\,R\Gamma (S,K)\,\to \,R\Gamma (S,Ra_{\ast }a^{\ast }K)\,\simeq \,R\Gamma (X_{\bullet },a^{\ast }K)}
\par\medskip
is an isomorphism. In particular, for an abelian sheaf F on S, there is a spectral sequence (5.2.3.2)\par
\par\medskip
\DFormula{(5.3.3.2)\,E_{1}^{pq}\,=\,H^{q}(X_{p},a_{p}^{\ast }F)\,\Rightarrow \,H^{p+q}(S,F).}
\par\medskip
For a complex, there is likewise, in hypercohomology, a spectral sequence\par
\par\medskip
\DFormula{(5.3.3.3)\,E_{1}^{pq}\,=\,H^{q}(X_{p},a_{p}^{\ast }K)\,\Rightarrow \,H^{p+q}(S,K).}
\par\medskip
In both cases, the \(E_{1}^{\bullet q}\) (q fixed) form a simplicial group, and\par
\par\medskip
\DFormula{d_{1}\,=\,\sum _{i}\,(−1)^{i}\,\delta _{i}\,:\,E_{1}^{pq}\,\to \,E_{1}^{p+1,q}.}
\DAsset{P0010-A01}{calculation of \(Ra_{\ast }\) and definition of cohomological descent, (5.2.7.1)-(5.3.2)}
\DAsset{P0010-A02}{descent isomorphism and spectral sequences, (5.3.3.1)-(5.3.3.3)}

\DArticleAnchor{0011}{11}{12}{678}{15}{P02}

Definition (5.3.4). — A continuous map a : X \(\to \) S is of cohomological descent if the augmentation morphism of cosq(X \(\to \) S),\par
\par\medskip
\DFormula{((X/S)^{\Delta _{n}})_{n\geq 0}\,\to \,S,}
\par\medskip
is of cohomological descent. We say that a is of universal cohomological descent if, for every u : S′ \(\to \) S, the continuous map a′ : X \(\times _{S}\) S′ \(\to \) S′ is of cohomological descent.\par
\par\medskip
(5.3.5) The fundamental results, proved in S.G.A.4, V bis, are as follows.\par
\par\medskip
(I) Continuous maps of universal cohomological descent form a Grothendieck topology on the category of topological spaces.\par
\par\medskip
It is called the topology of universal cohomological descent.\par
\par\medskip
(II) A proper (III.0.1) surjective map is of universal cohomological descent.\par
\par\medskip
(III) A map a : X \(\to \) S that admits local sections over S is of universal cohomological descent.\par
\par\medskip
(IV) Let a : \(X_{\bullet }\) \(\to \) S be an augmented k-truncated simplicial space \((−1\) \(\leq \) k \(\leq \) \(\infty )\). For k \(\geq \) n \(\geq \) \(−1\), let \(\varphi _{n}\) : cosq \(X_{\bullet }\) \(\to \) cosq \(\operatorname{sq}_{n}\) \(X_{\bullet }\) be the evident arrow. We say that \(X_{\bullet }\) is a k-truncated hypercover of S for the topology of universal cohomological descent if the maps\par
\par\medskip
\DFormula{(5.3.5.1)\,(\varphi _{n})_{n+1}\,:\,X_{n+1}\,\to \,(\operatorname{cosq}\,\operatorname{sq}_{n}\,X_{\bullet })_{n+1}\,\,\,(−1\,\leq \,n\,\leq \,k−1)}
\par\medskip
are of universal cohomological descent. If \(X_{\bullet }\) is such a hypercover, then the simplicial space augmented to S, \(\operatorname{cosq}(X_{\bullet })\), is of cohomological descent.\par
\par\medskip
(V) Let a be a morphism of simplicial topological spaces augmented to S\par
\par\medskip
\DAsset{P0011-A02}{diagram \(X_{\bullet }\) \(\to \) \(Y_{\bullet }\) over S}
\par\medskip
We say that a is a hypercover for the topology of universal cohomological descent if the evident maps \(X_{n}\) \(\to \) \((\operatorname{cosq}_{n−1}^{Y_{\bullet }}\) \(\operatorname{sq}_{n−1}X_{\bullet })_{n}\) are of universal cohomological descent. If a is such a hypercover, then, for every K \(\in \) Ob \(D^{+}(S)\),\par
\par\medskip
\DFormula{a^{\ast }\,:\,Ry_{\ast }y^{\ast }K\,\simeq \,Rx_{\ast }x^{\ast }K.}
\par\medskip
(5.3.6) For k \(=\) \(\infty \), (IV) states that a : \(X_{\bullet }\) \(\to \) S is of cohomological descent if the \((\varphi _{n})_{n+1}\) are of universal cohomological descent. For n \(=\) \(−1\) or 0, these maps are\par
\par\medskip
\DFormula{n\,=\,−1\,:\,X_{0}\,\xrightarrow{a}\,S}
\par\medskip
\DFormula{n\,=\,0\,:\,X_{1}\,\xrightarrow{(\delta _{0},\delta _{1})}\,X_{0}\,\times _{S}\,X_{0}.}
\DAsset{P0011-A01}{coskeleton, hypercover condition, and formula (5.3.5.1)}

\DArticleAnchor{0012}{12}{13}{679}{16}{P02}

For n \(=\) 1, (cosq \(\operatorname{sq}_{1}(X_{\bullet }))_{1}\) is the subspace of \(X_{1}\) \(\times _{S}\) \(X_{1}\) \(\times _{S}\) \(X_{1}\) formed by triples (x,y,z) such that \(\delta _{0}x\) \(=\) \(\delta _{0}y\), \(\delta _{1}x\) \(=\) \(\delta _{0}z\), and \(\delta _{1}y\) \(=\) \(\delta _{1}z\). The map \((\varphi _{1})_{2}\) sends x to \((\delta _{0}x\), \(\delta _{1}x\), \(\delta _{2}x)\).\par
\par\medskip
\DFormula{For\,k\,=\,0,\,(IV)\,is\,Definition\,(5.3.4).}
\par\medskip
Example (5.3.7). — Let \(𝒰\) \(=\) \((U_{i})_{i\in I}\) be an open cover, or a locally finite closed cover, of S. Let X \(=\) \(⨿_{i\in I}\) \(U_{i}\). Then a : X \(\to \) S is of cohomological descent. The spectral sequence (5.3.3.2) for \(X_{\bullet }\) \(=\) cosq(X \(\to \) S) is precisely the Leray spectral sequence of the cover \(𝒰\).\par
\par\medskip
(5.3.8) Let a : \(X_{\bullet }\) \(\to \) S be as in (5.3.5) (IV). We say that \(X_{\bullet }\) is a proper k-truncated hypercover of S if the arrows (5.3.5.1) are proper and surjective. For k \(=\) \(\infty \), we simply speak of a proper hypercover.\par
\par\medskip
6. EXAMPLES OF SIMPLICIAL TOPOLOGICAL SPACES\par
\par\medskip
6.1. Classifying spaces.\par
\par\medskip
(6.1.1) Let u : X \(\to \) S be a continuous map. For every sheaf F on S, the sheaf \(u^{\ast }F\) is equipped with “descent data” relative to u: there is an isomorphism between the two inverse images of \(u^{\ast }F\) on X \(\times _{S}\) X, and this isomorphism satisfies a cocycle condition. If u locally admits a section over S, this construction defines an equivalence between the category of sheaves on S and the category of sheaves on X equipped with descent data.\par
\par\medskip
Take X to be a principal homogeneous space (on the left) for a group G over a space S. The G-equivariant sheaves on X are precisely the sheaves equipped with descent data: every equivariant sheaf on X is, in one and only one way (as an equivariant sheaf), the inverse image of a sheaf on S \(=\) X/G.\par
\par\medskip
(6.1.2) If a topological group G acts on a space X, then G acts on \(G^{\Delta _{n}}\) \(\times \) X by\par
\par\medskip
\DFormula{g·(g_{0},…,g_{n},x)\,=\,(g_{0}g^{−1},…,g_{n}g^{−1},gx).}
\par\medskip
We denote by \([X/G]_{\bullet }\) the simplicial space\par
\par\medskip
\DFormula{(6.1.2.1)\,[X/G]_{\bullet }\,=\,((G^{\Delta _{n}}\,\times \,X)/G)_{n\geq 0}.}
\par\medskip
a) If X is a principal homogeneous space for G over S \(=\) X/G, then the map\par
\par\medskip
\DFormula{G^{\Delta _{n}}\,\times \,X\,\to \,X^{\Delta _{n}}\,:\,(g_{0},…,g_{n},x)\,↦\,(g_{0}x,…,g_{n}x)}
\par\medskip
identifies \([X/G]_{n}\) with the iterated fiber product \((X/S)^{\Delta _{n}}\):\par
\par\medskip
\DFormula{\operatorname{cosq}(X\,\to \,S)\,=\,([X/G]_{\bullet }\,\to \,S).}
\DAsset{P0012-A01}{coskeleton formula printed with subscript 1 and Example (5.3.7)-(5.3.8)}
\DAsset{P0012-A02}{G-action, simplicial classifying space, and identification (6.1.2.1)}

\DArticleAnchor{0013}{13}{14}{680}{17}{P03}

In particular, by (5.3.5 (III)), we have\par
\par\medskip
\DFormula{(6.1.2.2)\,H^{\bullet }([X/G]_{\bullet })\,\simeq \,H^{\bullet }(X/G)\,\,\,(for\,a\,principal\,homogeneous\,space).}
\par\medskip
b) For every n, \(G^{\Delta _{n}}\) \(\times \) X is a principal homogeneous space for G over \([X/G]_{n}\). For every equivariant sheaf F on X, \(\operatorname{pr}_{2}^{\ast }F\) is an equivariant sheaf on \(G^{\Delta _{n}}\) \(\times \) X; by (6.1.1), the latter is the inverse image of \(F^{n}\) on \([X/G]_{n}\). Thus every equivariant sheaf F on X defines a sheaf on \([X/G]_{\bullet }\). It is easy to verify that this gives an equivalence between the category of equivariant sheaves on X and the category of sheaves \(F^{\bullet }\) on \([X/G]_{\bullet }\) satisfying\par
\par\medskip
(*) For every f : \(\Delta _{n}\) \(\to \) \(\Delta _{m}\), the structural morphism \(f^{\ast }F^{n}\) \(\to \) \(F^{m}\) is an isomorphism.\par
\par\medskip
c) The construction in b) is natural in (G,X,F). We set\par
\par\medskip
\DFormula{(6.1.2.3)\,H^{\bullet }(X,G;F)\,=\,H^{\bullet }([X/G]_{\bullet },F^{\bullet })}
\par\medskip
(the mixed cohomology of X, G with coefficients in F). Under the hypotheses of a), if F is the inverse image of \(F^{−1}\) on S \(=\) X/G, then\par
\par\medskip
\DFormula{(6.1.2.4)\,H^{\bullet }(X,G;F)\,=\,H^{\bullet }(X/G,F^{−1})\,\,\,(for\,a\,principal\,homogeneous\,space).}
\par\medskip
This generalizes (6.1.2.2) (which is the case F \(=\) \(\mathbb{Z})\).\par
\par\medskip
(6.1.3) Let \(P^{t}\) be the topological space consisting of one point. The simplicial classifying space of G, denoted \(B_{\bullet G}\), is the simplicial space\par
\par\medskip
\DFormula{B_{\bullet G}\,=\,[P^{t}/G]_{\bullet }.}
\par\medskip
Let P be a principal homogeneous space for G over S. The evident morphism\par
\par\medskip
\DFormula{\operatorname{cosq}(P\,\to \,S)\,=\,[P/G]_{\bullet }\,\to \,[P^{t}/G]_{\bullet }\,=\,B_{\bullet G}}
\par\medskip
defines the composite morphism\par
\par\medskip
\DFormula{(6.1.3.1)\,[P]\,:\,H^{\bullet }(B_{\bullet G})\,\to \,H^{\bullet }([P/G]_{\bullet })\,\xleftarrow[(6.1.2.2)]{\sim }\,H^{\bullet }(S).}
\par\medskip
We shall see below that in the good cases \(H^{\bullet }(B_{\bullet G})\) \(=\) \(H^{\bullet }(B_{G})\), and that the image of [P] consists of the characteristic classes of P.\par
\par\medskip
(6.1.4) Let G be a Lie group, let \(B_{G}\) be a classifying space for G, and let a : \(U_{G}\) \(\to \) \(B_{G}\) be the universal principal homogeneous G-space. Let X be a G-space; X \(\times \) \(U_{G}\) is a principal homogeneous G-space over X \(\times \) \(U_{G}/G\), so that, for every equivariant sheaf F on X, \(\operatorname{pr}_{1}^{\ast }F\) is the inverse image of a sheaf \(F^{G}\) on X \(\times \) \(U_{G}/G\). Since \(U_{G}\) is contractible, and by (6.1.2.2), we have\par
\par\medskip
\DFormula{(6.1.4.1)\,H^{\bullet }(X,G;F)\,\xrightarrow{\sim }\,H^{\bullet }(X\,\times \,U_{G},G;\operatorname{pr}_{1}^{\ast }F)\,\xleftarrow{\sim }\,H^{\bullet }(X\,\times \,U_{G}/G,F^{G}).}
\par\medskip
\DFormula{In\,particular,\,for\,X\,=\,P^{t},}
\par\medskip
\DFormula{(6.1.4.2)\,H^{\bullet }(B_{\bullet G})\,=\,H^{\bullet }(B_{G}).}
\DAsset{P0013-A01}{equivariant mixed cohomology and formulas (6.1.2.2)-(6.1.2.4)}
\DAsset{P0013-A02}{simplicial classifying space, characteristic map [P], and identifications (6.1.3)-(6.1.4.2)}

\DArticleAnchor{0014}{14}{15}{681}{18}{P03}

One verifies that the isomorphism (6.1.4.2) is the special case of (6.1.3.1) with S \(=\) \(B_{G}\) and P \(=\) \(U_{G}\).\par
\par\medskip
(6.1.5) The spectral sequence (5.2.3.2) for \(B_{\bullet G}\),\par
\par\medskip
\DFormula{E_{1}^{pq}\,=\,H^{q}(G^{\Delta _{p}}/G)\,\Rightarrow \,H^{p+q}(B_{\bullet G})\,=\,H^{p+q}(B_{G}),}
\par\medskip
is essentially the Eilenberg-Moore spectral sequence. Let us briefly recall how it relates the rational cohomologies of G and \(B_{G}\) when G is connected.\par
\par\medskip
a) The algebra \(H^{\bullet }(G,\mathbb{Q})\) is a connected graded Hopf algebra of finite dimension over \(\mathbb{Q}\). If \(P^{\bullet }(G)\) is the graded module of its primitive elements, then\par
\par\medskip
\DFormula{H^{\bullet }(G,\mathbb{Q})\,=\,\Lambda P^{\bullet }(G),}
\par\medskip
and the generators of \(P^{\bullet }(G)\) have odd degrees.\par
\par\medskip
b) The simplicial algebra \((E_{1}^{p\bullet })_{p\geq 0}\) is\par
\par\medskip
\DFormula{E_{1}^{p\bullet }\,=\,\Lambda (P^{\bullet }(G)^{\Delta _{p}}/P^{\bullet }(G));}
\par\medskip
it is the exterior algebra of the suspension of the constant cosimplicial module \(P^{\bullet }(G)\). Hence (Quillen [12]) \(E_{2}^{p\bullet }\) \(=\) \(\operatorname{Sym}^{p}(P^{\bullet }(G))\). The terms \(E_{2}^{pq}\) can be nonzero only when p \(+\) q is even; therefore \(E_{2}^{pq}\) \(=\) \(E_{\infty }^{pq}\), and, for a suitable filtration, there is a canonical isomorphism\par
\par\medskip
\DFormula{\operatorname{Gr}\,H^{\bullet }(B_{G},\mathbb{Q})\,\simeq \,\operatorname{Sym}^{\bullet }(P^{\bullet }(G)[−1])}
\par\medskip
and a noncanonical isomorphism\par
\par\medskip
\DFormula{H^{\bullet }(B_{G},\mathbb{Q})\,\simeq \,\operatorname{Sym}^{\bullet }(P^{\bullet }(G)[−1]).}
\par\medskip
(6.1.6) Take G to be a (complex) linear algebraic group.\par
\par\medskip
a) If T is a maximal torus of G, with Weyl group W, then\par
\par\medskip
\DFormula{(6.1.6.1)\,H^{\bullet }(B_{G},\mathbb{Q})\,\simeq \,H^{\bullet }(B_{T},\mathbb{Q})^{W}.}
\par\medskip
b) For a torus T with character group X(T),\par
\par\medskip
\DFormula{(6.1.6.2)\,H^{\bullet }(T,\mathbb{Z})\,\simeq \,\Lambda ^{\bullet }X(T)}
\par\medskip
(an isomorphism of graded Hopf algebras).\par
\par\medskip
We shall use a) only in the following weaker form:\par
\par\medskip
a′) The map \(H^{\bullet }(B_{G},\mathbb{Q})\) \(\to \) \(H^{\bullet }(B_{T},\mathbb{Q})\) is injective (splitting principle).\par
\par\medskip
For completeness, let us recall a proof of a′).\par
\par\medskip
If B is a Borel subgroup of G, the bundle \(U_{G}/B\) over \(B_{G}\) is a bundle of flag manifolds. By ([3] 2.1 \(+\) 2.6.3), explained more fully in P. A. Griffiths, Periods III, Publ. Math. I.H.E.S., 38, prop. (3.1), or by [1], the Leray spectral sequence of \(U_{G}/B\) \(\to \) \(B_{G}\) degenerates in rational cohomology. Thus \(H^{\bullet }(B_{G},\mathbb{Q})\) \(\hookrightarrow \) \(H^{\bullet }(U_{G}/B,\mathbb{Q})\). The conclusion follows from \(U_{G}/B\) \(\sim \) \(B_{B}\) \(\sim \) \(B_{T}\).\par
\DAsset{P0014-A01}{Eilenberg-Moore spectral sequence and Hopf-algebra formulas, (6.1.5)}
\DAsset{P0014-A02}{maximal-torus formulas and the splitting principle, (6.1.6)}

\DArticleAnchor{0015}{15}{16}{682}{19}{P03}

6.2. Construction of hypercovers.\par
\par\medskip
(6.2.1) Let \(X_{\bullet }\) be a simplicial set. Denote by \(D(\Delta _{n},\Delta _{m})\) the set of order-preserving surjections from \(\Delta _{n}\) onto \(\Delta _{m}\) (degeneracy operators), and set\par
\par\medskip
\DFormula{(6.2.1.1)\,N(X_{n})\,=\,X_{n}\,−\,⋃_{s\in D(\Delta _{n},\Delta _{n−1})}\,s(X_{n−1}).}
\par\medskip
Recall that, for every n, the map\par
\par\medskip
\DFormula{(6.2.1.2)\,⨿s\,:\,⨿_{m\leq n}\,⨿_{s\in D(\Delta _{n},\Delta _{m})}\,N(X_{m})\,\to \,X_{n}}
\par\medskip
is bijective.\par
\par\medskip
Definition (6.2.2). — A simplicial topological space is called s-split if the maps (6.2.1.2) are homeomorphisms.\par
\par\medskip
Let \(X_{\bullet }\) be a k-truncated simplicial set. For n \(\leq \) k, \(N(X_{n})\) is again defined by (6.2.1.1), and (6.2.1.2) is a bijection. A k-truncated simplicial topological space is called s-split if (6.2.1.2) is a homeomorphism for n \(\leq \) k.\par
\par\medskip
(6.2.3) For an s-split \((n+1)-truncated\) simplicial topological space X augmented over S, let \(\alpha (X)\) be the triple consisting of \(\operatorname{sq}_{n}(X)\), \(NX_{n+1}\), and the evident map from \(NX_{n+1}\) to (cosq \(\operatorname{sq}_{n}\) \(X)_{n+1}\).\par
\par\medskip
\DFormula{This\,triple\,(Y,N,\beta )\,satisfies}
\par\medskip
(*) Y is an s-split n-truncated simplicial topological space augmented over S, and \(\beta \) is a continuous map from N to (cosq \(Y)_{n+1}\).\par
\par\medskip
\DFormula{Proposition\,(6.2.4)\,(S.G.A.4,\,V\,bis\,(5.1.3)).\,—\,Let\,(Y,N,\beta )\,satisfy\,(\ast )\,above.}
\par\medskip
(i) Up to unique isomorphism, there exists one and only one s-split \((n+1)-truncated\) simplicial topological space X augmented over S such that \(\alpha (X)\) \(\simeq \) \((Y,N,\beta )\).\par
\par\medskip
(ii) To give f : X \(\to \) Z is equivalent to giving:\par
\par\medskip
\DFormula{a)\,a\,morphism\,f′\,:\,Y\,\to \,\operatorname{sq}_{n}(Z),}
\par\medskip
b) a morphism f″ : N \(\to \) \(Z_{n+1}\) making the diagram commute:\par
\par\medskip
\DAsset{P0015-A02}{commutative square N \(\to \) (cosq \(Y)_{n+1}\), \(Z_{n+1}\) \(\to \) (cosq \(\operatorname{sq}_{n}\) \(Z)_{n+1}\)}
\par\medskip
Proposition (6.2.4) also applies to simplicial objects in categories \(\mathcal{C}\) other than the category of topological spaces. It is enough that \(\mathcal{C}\) satisfy\par
\DAsset{P0015-A01}{nondegenerate simplices, s-split spaces, and formulas (6.2.1.1)-(6.2.3)}

\DArticleAnchor{0016}{16}{17}{683}{20}{P03}

(6.2.4.1) Finite inverse limits exist. Finite sums exist and are disjoint and universal.\par
\par\medskip
(6.2.5) This proposition permits the following inductive construction of proper hypercovers of S.\par
\par\medskip
0) Choose \(f_{0}\) : \(X_{0}\) \(\to \) S, proper and surjective. \(\lbrace X_{0}\rbrace \) is a proper 0-truncated hypercover of S (5.3.8), and it is s-split.\par
\par\medskip
1) Choose \(f_{1}\) : \(N_{1}\) \(\to \) \(\operatorname{cosq}(\lbrace X_{0}\rbrace )_{1}\), i.e. \(f_{1}\) : \(N_{1}\) \(\to \) \(X_{0}\) \(\times _{S}\) \(X_{0}\). Applying (6.2.4), associate with \(f_{1}\) the s-split augmented 1-truncated simplicial topological space\par
\par\medskip
\DAsset{P0016-A01}{diagram \({}_{1}X_{\bullet }\) : \(N_{1}\) \(⨿\) \(X_{0}\) \(\rightleftarrows \) \(X_{0}\) \(\to \) S}
\par\medskip
Suppose \(f_{1}\) has been chosen so that\par
\par\medskip
\DFormula{f′_{1}\,:\,N_{1}\,⨿\,X_{0}\,\to \,X_{0}\,\times _{S}\,X_{0}}
\par\medskip
is proper and surjective (for example, \(f_{1}\) proper and surjective). Then \({}_{1}X_{\bullet }\) is an s-split proper 1-truncated hypercover of S.\par
\par\medskip
……\par
\par\medskip
\(k+1)\) Suppose that an s-split proper k-truncated hypercover \({}_{k}X_{\bullet }\) \(\to \) S has already been constructed. Choose \(f_{k+1}\) : \(N_{k+1}\) \(\to \) \((\operatorname{cosq}({}_{k}X_{\bullet }))_{k+1}\); applying (6.2.4), associate with \(f_{k+1}\) an s-split augmented \(k+1-truncated\) semisimplicial topological space \({}_{k+1}X_{\bullet }\). Suppose that \(f_{k+1}\) is such that\par
\par\medskip
\DFormula{f′_{k+1}\,:\,{}_{k+1}X_{k+1}\,\to \,\operatorname{cosq}({}_{k}X_{\bullet })_{k+1}}
\par\medskip
is proper and surjective (for example, \(f_{k+1}\) proper and surjective). Then \({}_{k+1}X_{\bullet }\) is an s-split proper \(k+1-truncated\) hypercover of S.\par
\par\medskip
……\par
\par\medskip
\(\infty )\) The \({}_{k}X_{\bullet }\) so constructed are the successive skeleta of an s-split proper hypercover of S.\par
\par\medskip
(6.2.6) A simplicial scheme \(X_{\bullet }\) over \(\mathbb{C}\) is called smooth if the \(X_{n}\) are smooth; it is called proper if the \(X_{n}\) are compact. A normal-crossings divisor \(D_{\bullet }\) in \(X_{\bullet }\), assumed smooth, is a family \(D_{n}\) \(\subset \) \(X_{n}\) of normal-crossings divisors (3.1.2) such that \(U_{n}\) \(=\) \(X_{n}\) \(−\) \(D_{n}\) form a simplicial subscheme \(U_{\bullet }\) of \(X_{\bullet }\). This definition is justified by the following lemma.\par
\par\medskip
Lemma (6.2.7). — If \(D_{\bullet }\) is a normal-crossings divisor in \(X_{\bullet }\), then the logarithmic de Rham complexes \((\Omega _{X_{n}}^{\bullet }(\operatorname{log}\) \(D_{n}))_{n\geq 0}\), equipped with the weight filtration (3.1.5), form a filtered complex on \(X_{\bullet }\).\par
\par\medskip
This follows from (3.1.3 (ii)). The complex \((\Omega _{X_{n}}^{\bullet }(\operatorname{log}\) \(D_{n}))_{n\geq 0}\) will be denoted \(\Omega _{X_{\bullet }}^{\bullet }(\operatorname{log}\) \(D_{\bullet })\).\par
\par\medskip
(6.2.8) Using (6.2.5), one proves that for every separated scheme S over \(\mathbb{C}\) there exist:\par
\DAsset{P0016-A02}{successive skeleta, normal-crossings divisors, and logarithmic de Rham complexes}

\DArticleAnchor{0017}{17}{18}{684}{21}{P03}

\(\alpha )\) a simplicial scheme \(X_{\bullet }\) over \(\mathbb{C}\), proper and smooth, which may be chosen s-split;\par
\par\medskip
\DFormula{\beta )\,a\,normal-crossings\,divisor\,D_{\bullet }\,in\,X_{\bullet };\,set\,U_{\bullet }\,=\,X_{\bullet }\,−\,D_{\bullet };}
\par\medskip
\(\gamma )\) an augmentation a : \(U_{\bullet }\) \(\to \) S making \(U_{\bullet }^{\mathrm{an}}\) a proper hypercover of \(S^{\mathrm{an}}\).\par
\par\medskip
Moreover, any two such systems are dominated by a common third one, and a morphism u : S \(\to \) T can be dominated by a morphism\par
\par\medskip
\DAsset{P0017-A01}{diagram \(X_{\bullet }\) \(\to \) \(Y_{\bullet }\), \(U_{\bullet }\) \(\to \) \(V_{\bullet }\), S \(\to \) T}
\par\medskip
of such systems... (see S.G.A.4, V bis).\par
\par\medskip
6.3. Relative cohomology.\par
\par\medskip
(6.3.1) The mapping-cone construction for morphisms of simplicial sets carries over unchanged to simplicial objects in any category \(\mathcal{C}\) having a final object e and finite sums. For u : \(Y_{\bullet }\) \(\to \) \(X_{\bullet }\), the cone C(u) satisfies\par
\par\medskip
\DFormula{C(u)_{n}\,=\,X_{n}\,⨿\,⨿_{i<n}\,Y_{i}\,⨿\,e.}
\par\medskip
\DFormula{We\,shall\,take\,\mathcal{C}\,to\,be:}
\par\medskip
a) the category of topological spaces and continuous maps (final object: \(P^{t})\);\par
\par\medskip
b) the category of pairs (X,F) consisting of a topological space X and an abelian sheaf F on X, an arrow (u,f) : (Y,F) \(\to \) (X,G) being a continuous map u : Y \(\to \) X together with a u-morphism (5.1.5) f : G \(\to \) F (final object: \((P^{t},0))\).\par
\par\medskip
(6.3.2) Let u : \(Y_{\bullet }\) \(\to \) \(X_{\bullet }\) be a morphism of simplicial topological spaces, with cone C(u); let F be an abelian sheaf on \(X_{\bullet }\), G an abelian sheaf on \(Y_{\bullet }\), and f : G \(\to \) F a u-morphism. The cone C(f) of f is an abelian sheaf on C(u), and we set\par
\par\medskip
\DFormula{(6.3.2.1)\,H^{n}(C(u),C(f))\,=\,H^{n}(X_{\bullet }\,\operatorname{mod}\,Y_{\bullet },F\,\operatorname{mod}\,G).}
\par\medskip
These are the relative cohomology groups. They fit naturally into a long exact sequence\par
\par\medskip
\DFormula{(6.3.2.2)\,…\,\to \,H^{i}(X_{\bullet }\,\operatorname{mod}\,Y_{\bullet },F\,\operatorname{mod}\,G)\,\to \,H^{i}(X_{\bullet },F)\,\to \,H^{i}(Y_{\bullet },G)\,\to \,…}
\par\medskip
(6.3.3) More generally, let L and K be bounded-below complexes of abelian sheaves on \(Y_{\bullet }\) and \(X_{\bullet }\), respectively, and let f : K \(\to \) L be a u-morphism. This gives a complex C(f) on C(u). In hypercohomology we again set\par
\par\medskip
\DFormula{H^{n}(C(u),C(f))\,=\,H^{n}(X_{\bullet }\,\operatorname{mod}\,Y_{\bullet },K\,\operatorname{mod}\,L).}
\DAsset{P0017-A02}{mapping cone and relative-cohomology formulas (6.3.1)-(6.3.3)}

\DArticleAnchor{0018}{18}{19}{685}{22}{P03}

These groups occur in an exact sequence analogous to (6.3.2.2), arising, in the appropriate derived category, from a distinguished triangle\par
\par\medskip
(6.3.3.1)\par
\DAsset{P0018-A01}{triangle \(R\Gamma (Y_{\bullet },K)\), \(R\Gamma (C(u),C(f))\), \(R\Gamma (X_{\bullet },L)\)}
\par\medskip
(6.3.4) The construction presented above is not the only possible one. It has the disadvantage that, even when one starts from genuine topological spaces X and Y (i.e. constant simplicial spaces), one is led to consider nonconstant simplicial spaces.\par
\par\medskip
6.4. Multisimplicial spaces.\par
\par\medskip
(6.4.1) Let r be an integer \(\geq \) 0. An r-simplicial object \(Z_{\bullet }\) of a category \(\mathcal{C}\) is a contravariant functor from \((\Delta )^{r}\) to \(\mathcal{C}\). The diagonal simplicial object \(\delta Z_{\bullet }\) is the composite functor\par
\par\medskip
\DFormula{(\Delta )\,\to \,(\Delta )^{r}\,\to \,\mathcal{C}.}
\par\medskip
(6.4.2) As in (5.1.6), one defines the topos of sheaves on an r-simplicial topological space. Let \(\Gamma ^{\bullet }\) be the functor\par
\par\medskip
\DFormula{F\,↦\,(\Gamma (X_{n_{1},…,n_{r}},F^{n_{1},…,n_{r}}))_{n_{1},…,n_{r}}}
\par\medskip
from sheaves on \(X_{\bullet }\) to r-cosimplicial sets. For small r it will often instead be denoted \(\Gamma ^{\bullet \bullet …\bullet }\) (r dots). There is a coaugmentation \(\Gamma (X_{\bullet },F^{\bullet })\) \(\to \) \(\Gamma ^{\bullet }(X_{\bullet },F^{\bullet })\). The cohomology functors with values in an abelian sheaf F are the derived functors of the “global sections” functor \(\Gamma \). They may be calculated by a procedure parallel to (5.2.3):\par
\par\medskip
\DFormula{(6.4.2.1)\,R\Gamma \,=\,sR\Gamma ^{\bullet }\,:\,D^{+}(X_{\bullet })\,\to \,D^{+}((Ab)).}
\par\medskip
A sheaf F on an r-simplicial topological space \(X_{\bullet }\) induces a sheaf \(\delta F\) on the diagonal simplicial space \(\delta X_{\bullet }\). The Cartier-Eilenberg-Zilber theorem gives\par
\par\medskip
\DFormula{(6.4.2.2)\,R\Gamma (X_{\bullet },K)\,\simeq \,R\Gamma (\delta X_{\bullet },\delta K).}
\par\medskip
(6.4.3) Restrict to r \(=\) 2. A bisimplicial object 1-augmented over a simplicial object \(S_{\bullet }\) is a contravariant functor from \((\Delta ^{+})\) \(\times \) \((\Delta )\) to \(\mathcal{C}\) such that \(S_{\bullet }\) is the composite functor\par
\par\medskip
\DFormula{(\Delta )\,\xrightarrow{(\Delta _{−1},\bullet )}\,(\Delta ^{+})\,\times \,(\Delta )\,\to \,\mathcal{C}.}
\par\medskip
To denote a bisimplicial object 1-augmented over \(S_{\bullet }\), with underlying bisimplicial object \(X_{\bullet \bullet }\), we shall use notation of the form\par
\par\medskip
\DFormula{a\,:\,X_{\bullet \bullet }\,\to \,S_{\bullet }.}
\DAsset{P0018-A02}{multicosimplicial global sections and 1-augmented bisimplicial objects}

\DArticleAnchor{0019}{19}{20}{686}{23}{P04}

For n \(\geq \) 0, a : \((X_{\bullet n})\) \(\to \) \(S_{n}\) is a simplicial object augmented toward \(S_{n}\). If F is a sheaf on \(X_{\bullet \bullet }\), the family \((a_{n\ast }(F\mid X_{\bullet n})\) on \(S_{n})_{n\geq 0}\) forms a sheaf on \(S_{\bullet }\); this defines a morphism of topoi\par
\par\medskip
\DFormula{a\,:\,X_{\bullet \bullet }^{\sim }\,\to \,S_{\bullet }^{\sim }.}
\par\medskip
It is explained in S.G.A.4, V bis that \(Ra_{\ast }\) is computed “component by component”:\par
\par\medskip
\DFormula{Ra_{\ast }K\mid S_{n}\,=\,Ra_{n\ast }(K\mid X_{\bullet n}).}
\par\medskip
It follows that if, for every n, a : \((X_{\bullet n})\) \(\to \) \(S_{n}\) is of cohomological descent, then a : \(X_{\bullet \bullet }\) \(\to \) \(S_{\bullet }\) is of cohomological descent: for every complex of abelian sheaves K \(\in \) \(D^{+}(S_{\bullet })\), one has\par
\par\medskip
\DFormula{K\,\widetilde{\to }\,Ra_{\ast }a^{\ast }K.}
\par\medskip
(6.4.4) In S.G.A.4, V bis it is shown that for every separated simplicial scheme \(S_{\bullet }\) there exist a bisimplicial scheme \(X_{\bullet \bullet }\), 1-augmented toward \(S_{\bullet }\), and i : \(X_{\bullet \bullet }\) \(\hookrightarrow \) \(\overline{X}_{\bullet \bullet }\), such that:\par
\par\medskip
a) The \(\overline{X}_{nm}\) are projective and smooth; \(X_{nm}\) is the complement of a normal-crossings divisor \(D_{nm}\) in \(\overline{X}_{nm}\). One may suppose that \(D_{nm}\) is a union of smooth divisors.\par
\par\medskip
b) For n \(\geq \) 0, \(X_{\bullet n}\) is a proper hypercover of \(S_{n}\). It may be taken to be s-split.\par
\par\medskip
The construction proceeds as in (6.2), but the recursion is more complicated. The uniqueness assertions of (6.2.8) remain valid, mutatis mutandis.\par
\par\medskip
7. COMPLEMENTS TO § 1\par
\par\medskip
7.1. Filtered derived category.\par
\par\medskip
This subsection supplements § 1, no. 4.\par
\par\medskip
(7.1.1) Let \(𝒜\) be an abelian category. Set:\par
\par\medskip
\(F𝒜\) (respectively \(F_{2}𝒜)\) \(=\) the category of filtered (respectively bifiltered) objects of \(𝒜\) with finite filtration(s).\par
\par\medskip
\(K^{+}F𝒜\) (respectively \(K^{+}F_{2}𝒜)\) \(=\) the category of bounded-below filtered (respectively bifiltered) complexes of objects of \(𝒜\), modulo homotopies respecting the filtration (respectively the filtrations).\par
\par\medskip
\(D^{+}F𝒜\) (respectively \(D^{+}F_{2}𝒜)\) \(=\) the triangulated category obtained from \(K^{+}F𝒜\) (respectively \(K^{+}F_{2}𝒜)\) by inverting filtered (respectively bifiltered) quasi-isomorphisms (1.3.6). This is the filtered (respectively bifiltered) derived category.\par
\par\medskip
(7.1.2) A filtered quasi-isomorphism u : (K,F) \(\to \) (K′,F′) induces an isomorphism of spectral sequences u : \(E_{r}^{\bullet \bullet }(K,F)\) \(\to \) \(E_{r}^{\bullet \bullet }(K′,F′)\). An object K of \(D^{+}F𝒜\) therefore defines\par
\DAsset{P0019-A01}{end of (6.4.3), componentwise derived direct image, and opening of (6.4.4)}
\DAsset{P0019-A02}{end of (6.4.4), §7 headings, and definitions (7.1.1)-(7.1.2)}

\DArticleAnchor{0020}{20}{21}{687}{24}{P04}

a spectral sequence \(E_{r}^{\bullet \bullet }(K)\). Likewise, an object L of \(D^{+}F_{2}𝒜\) defines an assemblage of spectral sequences of the type considered in (1.4.9).\par
\par\medskip
(7.1.3) Let T be a left exact functor from \(𝒜\) to an abelian category \(\mathcal{B}\). Suppose that every object of \(𝒜\) embeds into an injective object. The functor T then “derives” to functors\par
\par\medskip
\DFormula{(7.1.3.1)\,RT\,:\,D^{+}(𝒜)\,\to \,D^{+}(\mathcal{B})\,\,(for\,reference)}
\par\medskip
\DFormula{(7.1.3.2)\,RT\,:\,D^{+}F(𝒜)\,\to \,D^{+}F(\mathcal{B}),}
\par\medskip
\DFormula{(7.1.3.3)\,RT\,:\,D^{+}F_{2}(𝒜)\,\to \,D^{+}F_{2}(\mathcal{B}).}
\par\medskip
They are computed as follows: if K′ is a T-acyclic resolution (respectively a filtered resolution, respectively a bifiltered resolution) of K ((1.4.5), (1.4.11)), then RT(K) \(=\) T(K′).\par
\par\medskip
The hypercohomology spectral sequence (for T) of K \(\in \) Ob \(D^{+}F(𝒜)\) is the spectral sequence of RTK \(\in \) Ob \(D^{+}F(\mathcal{B})\) (cf. (1.4.5)).\par
\par\medskip
(7.1.4) We shall need more precise results for the functors \(Ra_{\ast }\), where a : \(X_{\bullet }\) \(\to \) S is an augmentation of a simplicial topological space. The case S \(=\) \(P^{t}\), \(Ra_{\ast }\) \(=\) \(R\Gamma \) would suffice for us.\par
\par\medskip
Return to the notation of (5.1.12), and recall formula (5.2.6.1):\par
\par\medskip
\DFormula{Ra_{\ast }\,=\,sRa_{\bullet \ast }.}
\par\medskip
For every complex K \(\in \) Ob \(C^{+}(S_{\bullet })\), the simple complex sK carries a natural filtration L (5.1.9.3). A quasi-isomorphism u : K′ \(\widetilde{\to }\) K″ induces a filtered quasi-isomorphism u : (sK′,L) \(\to \) (sK″,L). Thus s factors through\par
\par\medskip
\DFormula{s\,:\,D^{+}(S_{\bullet })\,\to \,D^{+}F(S)}
\par\medskip
and \(Ra_{\ast }\) factors through\par
\par\medskip
\DFormula{Ra_{\ast }\,:\,D^{+}(X_{\bullet })\,\to \,D^{+}F(S).}
\par\medskip
The spectral sequence of the filtered complex \((Ra_{\ast }K,L)\) \(\in \) \(D^{+}F(S)\) is precisely (5.2.7.1).\par
\par\medskip
(7.1.5) If K is filtered (respectively bifiltered), then \(Ra_{\bullet \ast }K\) is filtered (bifiltered): one has\par
\par\medskip
\DFormula{(7.1.5.1)\,Ra_{\bullet \ast }\,:\,D^{+}F(X_{\bullet })\,\to \,D^{+}F(S_{\bullet })}
\par\medskip
\DFormula{(7.1.5.2)\,Ra_{\bullet \ast }\,:\,D^{+}F_{2}(X_{\bullet })\,\to \,D^{+}F_{2}(S_{\bullet }).}
\par\medskip
(7.1.6) Let K be a complex of cosimplicial sheaves on S, equipped with an increasing filtration W. The diagonal filtration \(\delta (W,L)\) of W and L is the following increasing filtration of sK:\par
\par\medskip
\DFormula{(7.1.6.1)\,\delta (W,L)_{n}(sK)\,=\,⨁_{p,q}\,W_{n+p}(K^{p,q})}
\par\medskip
\DFormula{=\,\sum _{p}\,s(W_{n+p}(K))\,\cap \,L^{p}(sK).}
\DAsset{P0020-A01}{continuation of (7.1.2), derived functors (7.1.3), and opening of (7.1.4)}
\DAsset{P0020-A02}{factorizations in (7.1.4), (7.1.5), and diagonal filtration (7.1.6.1)}

\DArticleAnchor{0021}{21}{22}{688}{25}{P04}

One has\par
\par\medskip
\DFormula{(7.1.6.2)\,\operatorname{Gr}_{n}^{\delta (W,L)}(sK)\,\simeq \,⨁_{p}\,\operatorname{Gr}_{n+p}^{W}(K^{\bullet p})[−p].}
\par\medskip
The functor (K,W) \(↦\) \((sK,\delta (W,L))\) takes filtered quasi-isomorphisms to filtered quasi-isomorphisms and defines\par
\par\medskip
\DFormula{(7.1.6.3)\,(s,\delta )\,:\,(K,W)\,↦\,(sK,\delta (W,L))\,:\,D^{+}F(S_{\bullet })\,\to \,D^{+}F(S),}
\par\medskip
whence, by composition with (7.1.5.1),\par
\par\medskip
\DFormula{(7.1.6.4)\,(R\Gamma ,\delta (\,,L))\,:\,(K,W)\,↦\,(R\Gamma K,\delta (W,L))\,:\,D^{+}F(X_{\bullet })\,\to \,D^{+}F(S).}
\par\medskip
It follows from (7.1.6.2) that\par
\par\medskip
\DFormula{(7.1.6.5)\,\operatorname{Gr}_{n}^{\delta (W,L)}(Ra_{\ast }K)\,=\,⨁_{p}\,Ra_{p\ast }(\operatorname{Gr}_{n+p}^{W}\,K)[−p].}
\par\medskip
(7.1.7) If (K,W,F) is a bifiltered complex of cosimplicial sheaves, then sK carries the three filtrations W, F, and L and defines various bifiltered complexes.\par
\par\medskip
For example, with W increasing, the functor (K,W,F) \(↦\) \((K,\delta (W,L),F)\) takes bifiltered quasi-isomorphisms to bifiltered quasi-isomorphisms, and hence defines\par
\par\medskip
\DFormula{(7.1.7.1)\,(K,W,F)\,↦\,(K,\delta (W,L),F)\,:\,D^{+}F_{2}(S_{\bullet })\,\to \,D^{+}F_{2}(S).}
\par\medskip
By composition with (7.1.5.2), one obtains\par
\par\medskip
\DFormula{(7.1.7.2)\,(K,W,F)\,↦\,(R\Gamma K,\delta (W,L),F)\,:\,D^{+}F_{2}(X_{\bullet })\,\to \,D^{+}F_{2}(S),}
\par\medskip
and\par
\par\medskip
\DFormula{(7.1.7.3)\,\operatorname{Gr}_{n}^{\delta (W,L)}(R\Gamma K,F)\,=\,⨁_{p}\,Ra_{p\ast }(\operatorname{Gr}_{n+p}^{W}\,K,F)[−p]}
\par\medskip
\DFormula{in\,D^{+}F(S).}
\par\medskip
7.2. Complements to the lemma of the two filtrations.\par
\par\medskip
In this subsection we give a new proof of the lemma of the two filtrations (1.3.16) and some complements.\par
\par\medskip
(7.2.1) Let (K,F,W) be a bounded-below bifiltered complex of objects of an abelian category \(𝒜\). The filtration F is assumed biregular.\par
\par\medskip
We say that (K,F,W) is F-splittable if the filtered complex (K,W) can be represented as a sum of filtered complexes \((K_{n},W_{n})_{n\in \mathbb{Z}}\)\par
\par\medskip
\DFormula{(K,W)\,=\,⨁_{n}\,(K_{n},W_{n})}
\par\medskip
with\par
\par\medskip
\DFormula{F^{n}K\,=\,⨁_{n′\geq n}\,K_{n′}.}
\DAsset{P0021-A01}{formulas (7.1.6.2)-(7.1.7.3) for diagonal filtrations}
\DAsset{P0021-A02}{§7.2 and the F-splittable decomposition in (7.2.1)}

\DArticleAnchor{0022}{22}{23}{689}{26}{P04}

Let \(r_{0}\) be an integer \(\geq \) 0, or \(+\infty \). The following condition was considered in (1.3.16), (1.3.17):\par
\par\medskip
(7.2.2; \(r_{0})\) For every nonnegative integer r \(<\) \(r_{0}\), the differentials \(d_{r}\) of the graded complex \(E_{r}(K,W)\) are strictly compatible with the recurrent filtration defined by F.\par
\par\medskip
(7.2.3) It is clear that if (K,F,W) is F-splittable, condition (7.2.2; \(\infty )\) holds. Conversely, it appears that when condition (7.2.2; \(\infty )\) holds, everything behaves as though the functor \(\operatorname{Gr}_{F}\) were exact. For example, we shall show that the \(\operatorname{Gr}_{F}\) of the spectral sequence E(K,W) is then identified with the spectral sequence \(E(\operatorname{Gr}_{F}\) K,W), and that the spectral sequence E(K,F) degenerates \((E_{1}\) \(=\) \(E_{\infty })\).\par
\par\medskip
(7.2.4) It follows immediately from definition (1.3.8) that the first direct filtration of \(E_{r}(K,W)\) is the filtration \(F_{d}\) of \(E_{r}(K,W)\) by the images\par
\par\medskip
\DFormula{F_{d}^{p}(E_{r}(K,W))\,=\,\operatorname{Im}(E_{r}(F^{p}K,W)\,\to \,E_{r}(K,W)).}
\par\medskip
Dually, the second direct filtration (1.3.9) of \(E_{r}(K,W)\) is the filtration \(F_{d\ast }\) of \(E_{r}(K,W)\) by the kernels\par
\par\medskip
\DFormula{F_{d\ast }^{p}(E_{r}(K,W))\,=\,\operatorname{Ker}(E_{r}(K,W)\,\to \,E_{r}(K/F^{p}K,W)).}
\par\medskip
The recurrent filtration \(F_{\mathrm{rec}}\) of \(E_{r}(K,W)\) lies between these two filtrations (1.3.13 (iii)).\par
\par\medskip
Proposition (7.2.5). — Suppose that (K,F,W) satisfies (7.2.2; \(r_{0})\). Then\par
\par\medskip
\DFormula{(i)\,(F^{a}K/F^{b}K,F,W)\,again\,satisfies\,(7.2.2;\,r_{0}).}
\par\medskip
(ii) For r \(\leq \) \(r_{0}\), the sequence\par
\par\medskip
\DFormula{0\,\to \,E_{r}(F^{p}K,W)\,\to \,E_{r}(K,W)\,\to \,E_{r}(K/F^{p}K,W)\,\to \,0}
\par\medskip
is exact; for r \(=\) \(r_{0}\) \(+\) 1, the sequence\par
\par\medskip
\DFormula{E_{r}(F^{p}K,W)\,\to \,E_{r}(K,W)\,\to \,E_{r}(K/F^{p}K,W)}
\par\medskip
is exact. In particular, for r \(\leq \) \(r_{0}\) \(+\) 1, the two direct filtrations and the recurrent filtration of \(E_{r}(K,W)\) coincide.\par
\par\medskip
Fix an integer p and prove by induction on r the assertion\par
\par\medskip
\((\ast )_{r}\) If r \(\leq \) \(r_{0}\), \(E_{r}(F^{p}K,W)\) injects into \(E_{r}(K,W)\); if r \(\leq \) \(r_{0}\) \(+\) 1, its image is \(F_{\mathrm{rec}}^{p}(E_{r}(K,W))\).\par
\par\medskip
Assertion \((\ast )_{0}\) is always true. Assume \((\ast )_{r}\) and prove \((\ast )_{r+1}\). We may suppose that 0 \(\leq \) r \(\leq \) \(r_{0}\). We have a diagram\par
\par\medskip
\DAsset{P0022-A02}{two rows \(E_{r}(K,W)\) and \(E_{r}(F^{p}K,W)\), horizontal differentials \(d_{r}\), and vertical inclusions, the last labeled i}
\DAsset{P0022-A01}{condition (7.2.2; \(r_{0})\), direct filtrations, and Proposition (7.2.5)}

\DArticleAnchor{0023}{23}{24}{690}{27}{P04}

and the image of the vertical inclusions is \(F_{\mathrm{rec}}^{p}(E_{r}(K,W))\). Because i is injective,\par
\par\medskip
\DFormula{F_{\mathrm{rec}}^{p}(E_{r+1}(K,W))\,=\,\operatorname{Im}(\operatorname{Ker}(F_{\mathrm{rec}}^{p}\,E_{r}(K,W)\,\xrightarrow{d_{r}}\,E_{r}(K,W))\,\to \,E_{r+1}(K,W))}
\par\medskip
\DFormula{=\,\operatorname{Im}(\operatorname{Ker}(E_{r}(F^{p}K,W)\,\xrightarrow{d_{r}}\,E_{r}(F^{p}K,W))\,\to \,E_{r+1}(K,W))}
\par\medskip
\DFormula{=\,\operatorname{Im}(E_{r+1}(F^{p}K,W)\,\to \,E_{r+1}(K,W)).}
\par\medskip
If r \(<\) \(r_{0}\), \(d_{r}\) is strictly compatible with \(F_{\mathrm{rec}}\), whence\par
\par\medskip
\DFormula{d_{r}E_{r}(K,W)\,\cap \,E_{r}(F^{p}K,W)\,=\,d_{r}E_{r}(F^{p}K,W),}
\par\medskip
and therefore \(E_{r+1}(F^{p}K,W)\) injects into \(E_{r+1}(K,W)\). This proves \((\ast )_{r+1}\).\par
\par\medskip
The assertions \((\ast )_{r}\), together with the dual assertions, imply (ii).\par
\par\medskip
It follows from (ii) that for r \(<\) \(r_{0}\) and a \(\leq \) b, \(E_{r}(F^{b}K,W)\) injects into \(E_{r}(F^{a}K,W)\), and the differential \(d_{r}\) of \(E_{r}(F^{a}K,W)\) is strictly compatible with the first direct filtration of \(E_{r}(F^{a}K,W)\). Induction on r \(<\) \(r_{0}\) then shows that the first direct filtration of \(E_{r}(F^{a}K,W)\) coincides with the recurrent filtration, and (i) follows.\par
\par\medskip
(7.2.6) If condition (7.2.2; \(r_{0})\) holds and r \(\leq \) \(r_{0}\) \(+\) 1, we shall denote by F the filtration \(F_{d}\) \(=\) \(F_{d\ast }\) \(=\) \(F_{\mathrm{rec}}\) of \(E_{r}(K,W)\). The exact sequences (7.2.5) (ii)\par
\par\medskip
\DAsset{P0023-A02}{autodual diagram of the exact sequences associated with \(F^{p+1}K\) \(\subset \) \(F^{p}K\), containing \(E_{r}(K,W)\) and \(E_{r}(\operatorname{Gr}_{F}^{p}K,W)\)}
\par\medskip
define, for r \(\leq \) \(r_{0}\), an autodual isomorphism compatible with the differentials \(d_{r}\):\par
\par\medskip
\DFormula{(7.2.6.1)\,\operatorname{Gr}_{F}^{p}(E_{r}(K,W))\,\simeq \,E_{r}(\operatorname{Gr}_{F}^{p}K,W)\,\,(r\,\leq \,r_{0}).}
\par\medskip
(7.2.7) If condition (7.2.2; \(\infty )\) holds, the sequences\par
\par\medskip
\DFormula{0\,\to \,E_{r}(F^{p}K,W)\,\to \,E_{r}(K,W)\,\to \,E_{r}(K/F^{p}K,W)\,\to \,0}
\par\medskip
are exact for every r. If the filtration W is biregular, the sequence\par
\par\medskip
\DFormula{0\,\to \,E_{\infty }(F^{p}K,W)\,\to \,E_{\infty }(K,W)\,\to \,E_{\infty }(K/F^{p}K,W)\,\to \,0}
\par\medskip
is therefore exact. This sequence can be rewritten as\par
\par\medskip
\DFormula{0\,\to \,\operatorname{Gr}_{W}(H(F^{p}K))\,\to \,\operatorname{Gr}_{W}(H(K))\,\to \,\operatorname{Gr}_{W}(H(K/F^{p}K))\,\to \,0.}
\par\medskip
Lemma (7.2.7.1). — Let L be a complex equipped with a biregular filtration G. For \(\operatorname{Gr}_{G}(L)\) to be acyclic, it is necessary and sufficient that L be acyclic and that the differentials of L be strictly compatible with the filtration G.\par
\par\medskip
This is a special case of (1.3.2).\par
\DAsset{P0023-A01}{kernel identities and strict compatibility in the proof of (7.2.5)}

\DArticleAnchor{0024}{24}{25}{691}{28}{P04}

Apply this lemma, taking L to be the complex\par
\par\medskip
\DFormula{0\,\to \,H(F^{p}K)\,\to \,H(K)\,\to \,H(K/F^{p}K)\,\to \,0.}
\par\medskip
One finds that:\par
\par\medskip
a) this sequence is exact, i.e. the differentials of K are strictly compatible with the filtration F;\par
\par\medskip
b) the filtration on \(H(F^{p}K)\) deduced from the filtration W on \(F^{p}K\) coincides with the filtration induced by the filtration W on H(K);\par
\par\medskip
c) an analogous statement holds for \(K/F^{p}K\).\par
\par\medskip
In conclusion:\par
\par\medskip
Proposition (7.2.8). — If (K,F,W) satisfies (7.2.2; \(\infty )\), the spectral sequence E(K,F) degenerates \((E_{1}\) \(=\) \(E_{\infty })\). Moreover, there is an isomorphism of spectral sequences\par
\par\medskip
\DFormula{\operatorname{Gr}_{F}^{p}(E_{\bullet }(K,W))\,\simeq \,E_{\bullet }(\operatorname{Gr}_{F}^{p}K,W).}
\par\medskip
8. HODGE THEORY OF ALGEBRAIC SPACES\par
\par\medskip
8.1. Hodge complexes.\par
\par\medskip
(8.1.1) Let A be as in (III.0.3). An A-Hodge complex K of weight n consists of:\par
\par\medskip
\(\alpha )\) a complex \(K_{A}\) \(\in \) Ob \(D^{+}(A)\), such that \(H^{k}(K_{A})\) is a finite-type A-module for every k;\par
\par\medskip
\(\beta )\) a filtration F on \(K_{A}\) \(\otimes \) \(\mathbb{C}\), i.e. a filtered complex \((K_{\mathbb{C}},F)\) \(\in \) Ob \(D^{+}F(\mathbb{C})\), and an isomorphism \(\alpha \) : \(K_{\mathbb{C}}\) \(\simeq \) \(K_{A}\) \(\otimes \) \(\mathbb{C}\) in \(D^{+}(\mathbb{C})\).\par
\par\medskip
The following axioms are required:\par
\par\medskip
(CH 1) the differential d of \(K_{\mathbb{C}}\) is strictly compatible with the filtration F; in other words, the spectral sequence defined by \((K_{\mathbb{C}},F)\) degenerates at \(E_{1}\) \((E_{1}\) \(=\) \(E_{\infty })\) (1.3.2);\par
\par\medskip
(CH 2) for every k, the filtration F on \(H^{k}(K_{\mathbb{C}})\) \(\simeq \) \(H^{k}(K_{A})\) \(\otimes \) \(\mathbb{C}\) defines an A-Hodge structure of weight n \(+\) k on \(H^{k}(K_{A})\): the filtration F is (n \(+\) k)-opposed to its complex conjugate (this makes sense because A \(\subset \) \(\mathbb{R})\).\par
\par\medskip
(8.1.2) Let X be a topological space. A cohomological A-Hodge complex K of weight n on X consists of:\par
\par\medskip
\DFormula{\alpha )\,a\,complex\,K_{A}\,\in \,Ob\,D^{+}(X,A),}
\par\medskip
\DFormula{\beta )\,a\,filtered\,complex\,(K_{\mathbb{C}},F)\,\in \,Ob\,D^{+}F(X,\mathbb{C}),}
\par\medskip
\DFormula{\gamma )\,\mathrm{an}\,isomorphism\,\alpha \,:\,K_{\mathbb{C}}\,\simeq \,K_{A}\,\otimes \,\mathbb{C}\,in\,D^{+}(X,\mathbb{C}).}
\par\medskip
The following axiom is required:\par
\par\medskip
(CHC) The triple \((R\Gamma (K_{A})\), \(R\Gamma ((K_{\mathbb{C}},F))\), \(R\Gamma (\alpha ))\) is a Hodge complex of weight n.\par
\DAsset{P0024-A01}{application of the lemma and Proposition (7.2.8)}
\DAsset{P0024-A02}{opening of §8, Hodge complexes (8.1.1)-(8.1.2), and axioms}

\DArticleAnchor{0025}{25}{26}{692}{29}{P05}

When A \(=\) \(\mathbb{Z}\), we shall simply speak of a Hodge complex and of a cohomological Hodge complex.\par
\par\medskip
The following statement reformulates part of classical Hodge theory (cf. (2.2.1)).\par
\par\medskip
Scholium (8.1.3). — Let X be a compact Kähler manifold. Let \(K_{\mathbb{Z}}\) be the complex \(\mathbb{Z}[0]\) concentrated in the constant sheaf \(\mathbb{Z}\), let \(K_{\mathbb{C}}\) be the holomorphic de Rham complex \(\Omega _{X}^{\bullet }\), and let F be its stupid filtration. There is\par
\par\medskip
\DFormula{\alpha \,:\,K_{\mathbb{Z}}\,\otimes \,\mathbb{C}\,\simeq \,\mathbb{C}[0]\,\widetilde{\to }\,\Omega _{X}^{\bullet }}
\par\medskip
(Poincaré lemma), and \((K_{\mathbb{Z}}\), \((K_{\mathbb{C}},F)\), \(\alpha )\) is a cohomological Hodge complex of weight 0.\par
\par\medskip
Remark (8.1.4). — If K is a Hodge complex (respectively a cohomological Hodge complex) of weight n, then K[m] is a Hodge complex (respectively a cohomological Hodge complex) of weight n \(+\) m.\par
\par\medskip
(8.1.5) A mixed A-Hodge complex K consists of:\par
\par\medskip
\(\alpha )\) a complex \(K_{A}\) \(\in \) Ob \(D^{+}(A)\), such that \(H^{k}(K_{A})\) is a finite-type A-module for every k;\par
\par\medskip
\(\beta )\) a filtered complex \((K_{A\otimes \mathbb{Q}},W)\) \(\in \) Ob \(D^{+}F(A\otimes \mathbb{Q})\) (W an increasing filtration) and an isomorphism \(K_{A\otimes \mathbb{Q}}\) \(\simeq \) \(K_{A}\) \(\otimes \) \(\mathbb{Q}\) in \(D^{+}(A\otimes \mathbb{Q})\);\par
\par\medskip
\(\gamma )\) a bifiltered complex \((K_{\mathbb{C}},W,F)\) \(\in \) Ob \(D^{+}F_{2}(\mathbb{C})\) (W an increasing filtration and F a decreasing filtration) and an isomorphism \(\alpha \) : \((K_{\mathbb{C}},W)\) \(\simeq \) \((K_{A\otimes \mathbb{Q}},W)\) \(\otimes \) \(\mathbb{C}\) in \(D^{+}F(\mathbb{C})\).\par
\par\medskip
The following axiom is required:\par
\par\medskip
(CHM) For every n, the system formed by the complex \(\operatorname{Gr}_{n}^{W}(K_{A\otimes \mathbb{Q}})\) \(\in \) Ob \(D^{+}(A\otimes \mathbb{Q})\), the complex \(\operatorname{Gr}_{n}^{W}(K_{\mathbb{C}})\) \(\in \) Ob \(D^{+}F(\mathbb{C})\) filtered by F, and the isomorphism\par
\par\medskip
\DFormula{\operatorname{Gr}_{n}^{W}(\alpha )\,:\,\operatorname{Gr}_{n}^{W}(K_{A\otimes \mathbb{Q}})\,\otimes \,\mathbb{C}\,\simeq \,\operatorname{Gr}_{n}^{W}(K_{\mathbb{C}}),}
\par\medskip
is an \(A\otimes \mathbb{Q}-Hodge\) complex of weight n.\par
\par\medskip
(8.1.6) A cohomological mixed A-Hodge complex K on a topological space X consists of:\par
\par\medskip
\(\alpha )\) a complex \(K_{A}\) \(\in \) Ob \(D^{+}(X,A)\), such that the \(H^{k}(X,K_{A})\) are finite-type A-modules and \(H^{\bullet }(X,K_{A})\) \(\otimes \) \(\mathbb{Q}\) \(\widetilde{\to }\) \(H^{\bullet }(X,K_{A}\) \(\otimes \) \(\mathbb{Q})\);\par
\par\medskip
\(\beta )\) a filtered complex \((K_{A\otimes \mathbb{Q}},W)\) \(\in \) Ob \(DF^{+}(X,A\otimes \mathbb{Q})\) (W an increasing filtration) and an isomorphism \(K_{A\otimes \mathbb{Q}}\) \(\simeq \) \(K_{A}\) \(\otimes \) \(\mathbb{Q}\) in \(D^{+}(X,A\otimes \mathbb{Q})\);\par
\par\medskip
\(\gamma )\) a bifiltered complex \((K_{\mathbb{C}},W,F)\) \(\in \) Ob \(D^{+}F_{2}(X,\mathbb{C})\) (W an increasing filtration and F a decreasing filtration) and an isomorphism \(\alpha \) : \((K_{\mathbb{C}},W)\) \(\simeq \) \((K_{A\otimes \mathbb{Q}},W)\) \(\otimes \) \(\mathbb{C}\) in \(D^{+}F(X,\mathbb{C})\).\par
\par\medskip
The following axiom is required:\par
\par\medskip
(CHMC) For every n, the system formed by the complex \(\operatorname{Gr}_{n}^{W}(K_{A\otimes \mathbb{Q}})\) \(\in \) Ob \(D^{+}(X,A\otimes \mathbb{C})\), the complex \(\operatorname{Gr}_{n}^{W}(K_{\mathbb{C}})\) \(\in \) Ob \(D^{+}F(X,\mathbb{C})\) filtered by F, and the isomorphism\par
\par\medskip
\DFormula{\operatorname{Gr}_{n}^{W}(\alpha )\,:\,\operatorname{Gr}_{n}^{W}(K_{\mathbb{C}})\,\simeq \,\operatorname{Gr}_{n}^{W}(K_{A\otimes \mathbb{Q}})\,\otimes \,\mathbb{C}}
\par\medskip
is a cohomological \(A\otimes \mathbb{Q}-Hodge\) complex of weight n.\par
\DAsset{P0025-A01}{Scholium (8.1.3), Remark (8.1.4), and definition (8.1.5)}
\DAsset{P0025-A02}{axiom (CHM), cohomological definition (8.1.6), and axiom (CHMC)}

\DArticleAnchor{0026}{26}{27}{693}{30}{P05}

One checks immediately:\par
\par\medskip
Proposition (8.1.7). — If K \(=\) \((K_{A}\), \((K_{A\otimes \mathbb{Q}},W)\), \((K_{\mathbb{C}},W,F))\) is a cohomological mixed A-Hodge complex, then \(R\Gamma K\) \(=\) \((R\Gamma K_{A}\), \(R\Gamma (K_{A\otimes \mathbb{Q}},W)\), \(R\Gamma (K_{\mathbb{C}},W,F))\) is a mixed A-Hodge complex.\par
\par\medskip
Scholium (8.1.8). — Let U be the complement, in a proper smooth scheme X, of a normal-crossings divisor D, and let j : U \(\hookrightarrow \) X be the inclusion morphism. Let W be the canonical filtration of \(Rj_{\ast }\mathbb{Q}\) \(=\) \(Rj_{\ast }\mathbb{Z}\) \(\otimes \) \(\mathbb{Q}\):\par
\par\medskip
\DFormula{W_{n}(Rj_{\ast }\mathbb{Q})\,=\,\tau _{\leq n}(Rj_{\ast }\mathbb{Q}).}
\par\medskip
One has (3.1.8)\par
\par\medskip
\DFormula{(Rj_{\ast }\mathbb{Q})\,\otimes \,\mathbb{C}\,\simeq \,Rj_{\ast }\mathbb{C}\,\simeq \,j_{\ast }\Omega _{U}^{\bullet }\,\simeq \,\Omega _{X}^{\bullet }(\operatorname{log}\,D).}
\par\medskip
Let W be the weight filtration on \(\Omega _{X}^{\bullet }(\operatorname{log}\) D), and F its Hodge filtration (by the \(\sigma _{\geq p})\). Then (3.1.8) gives, in \(D^{+}F(X,\mathbb{C})\), an isomorphism\par
\par\medskip
\DFormula{(Rj_{\ast }\mathbb{Q},W)\,\otimes \,\mathbb{C}\,\simeq \,(\Omega _{X}^{\bullet }(\operatorname{log}\,D),W),}
\par\medskip
and, by (3.1.5) and (3.1.9),\par
\par\medskip
\DFormula{(Rj_{\ast }\mathbb{Z},\,(Rj_{\ast }\mathbb{Q},W),\,(\Omega _{X}^{\bullet }(\operatorname{log}\,D),W,F))}
\par\medskip
is a cohomological mixed Hodge complex on X.\par
\par\medskip
The following result is an “abstract” version of subsection (3.2).\par
\par\medskip
Scholium (8.1.9). — Let K be a mixed A-Hodge complex.\par
\par\medskip
(i) On the terms \(_{W}\) \(E_{r}^{pq}\) of the spectral sequence of \((K_{\mathbb{C}},W)\), the recurrent filtration and the two direct filtrations defined by F coincide.\par
\par\medskip
(ii) The filtration W[n] (1.1.2) on \(H^{n}(K_{A\otimes \mathbb{Q}})\) \(\simeq \) \(H^{n}(K_{A})\) \(\otimes \) \(\mathbb{Q}\) and the filtration F on \(H^{n}(K_{\mathbb{C}})\) \(\simeq \) \(H^{n}(K_{A})\) \(\otimes _{A}\) \(\mathbb{C}\) define on \(H^{n}(K_{A})\) a mixed A-Hodge structure ((2.3.1), enlarged by (III.0.4)).\par
\par\medskip
(iii) The morphisms \(_{W}\) \(d_{1}\) : \(_{W}\) \(E_{1}^{pq}\) \(\to \) \(_{W}\) \(E_{1}^{p+1,q}\) are compatible with the Hodge bigrading; in particular, they are strictly compatible with the Hodge filtration.\par
\par\medskip
(iv) The spectral sequence of \((K_{A\otimes \mathbb{Q}},W)\) degenerates at \(E_{2}\): \(_{W}\) \(E_{2}\) \(=\) \(_{W}\) \(E_{\infty }\).\par
\par\medskip
(v) The spectral sequence of \((K_{\mathbb{C}},F)\) degenerates at \(E_{1}\): \(_{F}\) \(E_{1}\) \(=\) \(_{F}\) \(E_{\infty }\).\par
\par\medskip
(vi) The spectral sequence of the complex \(\operatorname{Gr}_{F}^{p}(K_{\mathbb{C}})\), equipped with the filtration induced by W, degenerates at \(E_{2}\).\par
\par\medskip
The proof of (i), (ii), (iii), and (iv) is parallel to that of (3.2.5). The analogues of Lemmas (3.2.6) and (3.2.7) were taken as axioms (CH 1 and CH 2). The proofs of (3.2.8), (3.2.9), and (3.2.10) then carry over verbatim, and, as in (3.2.10), yield (i), (ii), (iii), and (iv). Assertions (v) and (vi) then follow from (7.2.8).\par
\par\medskip
(8.1.10) We shall abbreviate “differential graded” to DG and “differential graded with degrees bounded below” to \(DG^{+}\). A complex of DG A-modules may be viewed as a double complex of A-modules; the first degree is the degree of the complex, and the second is the one defined by the grading of the DG modules. We denote by \(D^{+}((A-\operatorname{mod}\) \(DG)^{+})\)\par
\DAsset{P0026-A01}{Proposition (8.1.7), Scholium (8.1.8), and logarithmic comparison}
\DAsset{P0026-A02}{Scholium (8.1.9), spectral-sequence degenerations, and opening of (8.1.10)}

\DArticleAnchor{0027}{27}{28}{694}{31}{P05}

the derived category of the category of bounded-below complexes of DG A-modules, with degrees uniformly bounded below.\par
\par\medskip
A DG mixed A-Hodge complex consists of:\par
\par\medskip
\DFormula{a)\,a\,complex\,K_{A}\,\in \,Ob\,D^{+}((A-\operatorname{mod}\,DG)^{+});}
\par\medskip
b) a filtered complex \((K_{A\otimes \mathbb{Q}},W)\) \(\in \) Ob \(D^{+}F(((A\otimes \mathbb{Q})-\operatorname{mod}\) \(DG)^{+})\), and an isomorphism \(K_{A}\) \(\otimes \) \(\mathbb{Q}\) \(\simeq \) \(K_{A\otimes \mathbb{Q}}\) in \(D^{+}\);\par
\par\medskip
c) a bifiltered complex \((K_{\mathbb{C}},F,W)\) \(\in \) Ob \(D^{+}F_{2}((\mathbb{C}-\operatorname{mod}\) \(DG)^{+})\), and an isomorphism \(K_{A\otimes \mathbb{Q}}\) \(\otimes \) \(\mathbb{C}\) \(\simeq \) \(K_{\mathbb{C}}\) in \(D^{+}F\).\par
\par\medskip
It is required that, for every n, the component \((K^{\bullet n},F,W)\) of second degree n of K be a mixed A-Hodge complex.\par
\par\medskip
We denote by L the filtration of sK by the second degree. We denote by \(Dec_{1}W\) the filtration of \(K_{A\otimes \mathbb{Q}}\) obtained from W by décalage (1.3.3). The filtration that it induces on \(sK_{A\otimes \mathbb{Q}}\) (again denoted \(Dec_{1}W)\) is not to be confused with the décalage of the filtration induced on \(sK_{A\otimes \mathbb{Q}}\) by W.\par
\par\medskip
Variant (8.1.10.1). — Cosimplicial (respectively r-cosimplicial) mixed A-Hodge complexes are defined similarly by replacing DG everywhere by cosimplicial (respectively r-cosimplicial). The functor (cosimplicial A-module) \(\to \) (DG A-module) takes cosimplicial mixed A-Hodge complexes to DG mixed A-Hodge complexes. Likewise for r-cosimplicial.\par
\par\medskip
(8.1.11) Let \(X_{\bullet }\) be a simplicial topological space. A cohomological mixed A-Hodge complex K on \(X_{\bullet }\) consists of:\par
\par\medskip
\DFormula{\alpha )\,a\,complex\,K_{A}\,\in \,Ob\,D^{+}(X_{\bullet },A);}
\par\medskip
\(\beta )\) a filtered complex \((K_{A\otimes \mathbb{Q}},W)\) \(\in \) Ob \(D^{+}F(X_{\bullet },A\otimes \mathbb{Q})\), and an isomorphism \(K_{A\otimes \mathbb{Q}}\) \(\simeq \) \(K_{A}\) \(\otimes \) \(\mathbb{Q}\) in \(D^{+}(X_{\bullet },A\otimes \mathbb{Q})\);\par
\par\medskip
\(\gamma )\) a bifiltered complex \((K_{\mathbb{C}},W,F)\) \(\in \) Ob \(D^{+}F_{2}(X_{\bullet },\mathbb{C})\), and an isomorphism\par
\par\medskip
\DFormula{(K_{\mathbb{C}},W)\,\simeq \,(K_{A\otimes \mathbb{Q}},W)\,\otimes \,\mathbb{C}}
\par\medskip
\DFormula{in\,D^{+}F(X_{\bullet },\mathbb{C}).}
\par\medskip
The following axiom is required:\par
\par\medskip
\((CHM_{\bullet })\) The restriction of K to each \(X_{n}\) is a cohomological mixed A-Hodge complex.\par
\par\medskip
For A \(=\) \(\mathbb{Z}\), we shall simply speak of a cohomological mixed Hodge complex.\par
\par\medskip
Example (8.1.12). — Let \(X_{\bullet }\) be a proper smooth simplicial scheme (over \(\mathbb{C})\), and let \(Y_{\bullet }\) be a normal-crossings divisor in \(X_{\bullet }\). Set \(U_{\bullet }\) \(=\) \(X_{\bullet }\) \(−\) \(Y_{\bullet }\), and let j denote the inclusion of \(U_{\bullet }\) in \(X_{\bullet }\). The constructions of (8.1.8) can be carried out “uni-\par
\DAsset{P0027-A01}{definition of a DG mixed A-Hodge complex and the \(L/Dec_{1}W\) filtrations}
\DAsset{P0027-A02}{cosimplicial variant, definition (8.1.11), and opening of Example (8.1.12)}

\DArticleAnchor{0028}{28}{29}{695}{32}{P05}

formly in \(X_{n}\)” as (5.2.5) and (6.2.7) show. They therefore provide a cohomological mixed Hodge complex\par
\par\medskip
\DFormula{(Rj_{\ast }\mathbb{Z},\,(Rj_{\ast }\mathbb{Q},\tau _{\leq }),\,(\Omega _{X_{\bullet }}^{\bullet }(\operatorname{log}\,Y_{\bullet }),W,F))}
\par\medskip
\DFormula{on\,X_{\bullet }.}
\par\medskip
(8.1.13) Let K be a cohomological mixed A-Hodge complex on \(X_{\bullet }\). Apply the functor \(R\Gamma ^{\bullet }\) to K. One obtains:\par
\par\medskip
\DFormula{a)\,a\,cosimplicial\,complex\,R\Gamma ^{\bullet }K_{A}\,\in \,Ob\,D^{+}((cosimplicial\,A-modules));}
\par\medskip
b) a filtered cosimplicial complex \(R\Gamma ^{\bullet }(K_{A\otimes \mathbb{Q}},W)\) \(\in \) Ob \(D^{+}F((cosimplicial\) \(A\otimes \mathbb{Q}-modules))\);\par
\par\medskip
c) a bifiltered cosimplicial complex\par
\par\medskip
\DFormula{R\Gamma ^{\bullet }(K_{\mathbb{C}},W,F)\,\in \,Ob\,D^{+}F_{2}((cosimplicial\,\mathbb{C}-vector\,spaces));}
\par\medskip
d) the isomorphisms (8.1.10.b), c)) between these objects.\par
\par\medskip
It is clear that the \(R\Gamma ^{\bullet }(K,W,F)\) described above is a cosimplicial mixed A-Hodge complex. It defines a DG mixed A-Hodge complex (8.1.10.1).\par
\par\medskip
(8.1.14) Let K be a DG mixed A-Hodge complex, and consider the spectral sequence defined by the filtration L and abutting to the cohomology of sK. By hypothesis, the \(E_{1}\) terms of this spectral sequence carry mixed A-Hodge structures. The content of the theorem below is that the entire spectral sequence, including its abutment, carries a mixed A-Hodge structure. In other words, it is a spectral sequence in the abelian category of mixed A-Hodge structures. The structure on the abutment corresponds to a natural “structure” of mixed A-Hodge complex on sK.\par
\par\medskip
Theorem (8.1.15). — Let K be a DG mixed A-Hodge complex.\par
\par\medskip
(i) With the notation of (8.1.10), (8.1.10.1), and (7.1.6), \((sK,\delta (W,L),F)\) is a mixed A-Hodge complex.\par
\par\medskip
(ii) One has\par
\par\medskip
\DFormula{(8.1.15.1)\,\,_{\delta (W,L)}E_{1}^{ab}(sK\,\otimes \,\mathbb{Q})\,=\,\sum _{a=\gamma −\beta ,\,b=\alpha +\beta }\,H^{\alpha }(\operatorname{Gr}_{\beta }^{W}\,K^{\bullet ,\gamma });}
\par\medskip
the complex \((_{\delta (W,L)}E_{1}^{\bullet b},d_{1})\) is the simple complex associated with the following double complex of \(A\otimes \mathbb{Q}-Hodge\) structures of weight b:\par
\par\medskip
\DAsset{P0028-A02}{double complex (8.1.15.2), with horizontal differentials \(\partial \) and vertical differentials d″, and components \(H^{b−\beta }(\operatorname{Gr}_{\beta }^{W}\) \(K^{\bullet ,\gamma })\) and their neighbors}
\DAsset{P0028-A01}{end of Example (8.1.12), construction (8.1.13), and introduction to (8.1.14)}

\DArticleAnchor{0029}{29}{30}{696}{33}{P05}

(iii) On the terms \(_{L}\) \(E_{r}\) of the spectral sequence defined by \((sK_{A\otimes \mathbb{Q}},L)\), the recurrent filtration and the two direct filtrations defined by \(Dec_{1}(W)\) coincide; likewise, on the \(_{L}\) \(E_{r}(sK_{\mathbb{C}})\), the recurrent filtration and the two direct filtrations induced by F coincide. We shall denote these filtrations by W and F.\par
\par\medskip
(iv) For r \(\geq \) 1, \((_{L}\) \(E_{r},W,F)\) is a mixed A-Hodge structure (III.0.4). The differentials \(d_{r}\) are morphisms of mixed A-Hodge structures.\par
\par\medskip
(v) The filtration L on \(H^{\bullet }(sK)\) is a filtration in the category of mixed Hodge structures, and\par
\par\medskip
\DFormula{(8.1.15.3)\,\operatorname{Gr}_{L}^{p}(H^{p+q}(sK),\delta (W,L)[p+q],F)\,=\,(_{L}\,E_{\infty }^{pq},W,F).}
\par\medskip
Remark (8.1.16). — In (8.1.15), the filtration \(Dec_{1}(W)\) (8.1.10) plays a larger role than W. One has\par
\par\medskip
\DFormula{(8.1.16.1)\,Dec(\delta (W,L))\,=\,Dec_{1}(W)\,on\,sK,}
\par\medskip
and (8.1.15.3) may be rewritten\par
\par\medskip
\DFormula{(8.1.16.2)\,\operatorname{Gr}_{L}(H(K),Dec_{1}W,F)\,=\,(_{L}\,E_{\infty },W,F).}
\par\medskip
In the same spirit, note that, with the notation of (8.1.8),\par
\par\medskip
\DFormula{R\Gamma (Rj_{\ast }\mathbb{Z},\,(Rj_{\ast }\mathbb{Q},Dec\,W),\,(\Omega _{X}^{\bullet }(\operatorname{log}\,D),Dec\,W,F))}
\par\medskip
depends only on U and not on the chosen compactification X.\par
\par\medskip
(8.1.17) Let us prove (8.1.15). Assertions (i) and (ii) follow from\par
\par\medskip
\(\operatorname{Gr}_{n}^{\delta (W,L)}(sK\) \(\otimes \) \(\mathbb{Q})\) \(=\) \(\sum _{p−q=n}\) \(\operatorname{Gr}_{p}^{W}\) \(\operatorname{Gr}_{L}^{q}(sK\) \(\otimes \) \(\mathbb{Q})\) \(=\) \(\sum _{p−q=n}\) \((\operatorname{Gr}_{p}^{W}\) \(K^{\bullet q})[−q]\)\par
\par\medskip
(this isomorphism is compatible with F, and one applies (8.1.4)). Verification of (8.1.15.2) is left to the reader.\par
\par\medskip
Consider the assertions:\par
\par\medskip
\((a_{r})\) The morphisms \(d_{r}\) of \(_{L}\) \(E_{r}(sK_{A\otimes \mathbb{Q}})\) are strictly compatible with the recurrent filtration defined by \(Dec_{1}(W)\).\par
\par\medskip
\((b_{r})\) The morphisms \(d_{r}\) of \(_{L}\) \(E_{r}(sK_{\mathbb{C}})\) are strictly compatible with the recurrent filtration defined by F.\par
\par\medskip
\((c_{r})\) \(_{L}\) \(E_{r}\), equipped with the filtrations considered in \((a_{r})\) and \((b_{r})\), is a mixed A-Hodge structure.\par
\par\medskip
\((d_{r})\) The \(d_{r}\) are morphisms of mixed A-Hodge structures.\par
\par\medskip
We shall prove \((a_{r})\) (r \(\geq \) 0), \((b_{r})\) (r \(\geq \) 0), \((c_{r})\) (r \(\geq \) 1), and \((d_{r})\) (r \(\geq \) 1) by simultaneous induction.\par
\par\medskip
\DFormula{A)\,Proof\,of\,(a_{0}):\,one\,has}
\par\medskip
\DFormula{(_{L}\,E_{0}^{p\bullet }(sK_{A\otimes \mathbb{Q}}),Dec_{1}W)\,=\,(K_{A\otimes \mathbb{Q}}^{\bullet p},Dec\,W).}
\par\medskip
By (8.1.9)(iv), the spectral sequence of \((K_{A\otimes \mathbb{Q}}^{\bullet p},W)\) degenerates at \(E_{2}\). By (1.3.4), the spectral sequence of \((K_{A\otimes \mathbb{Q}}^{\bullet p},Dec\) W) degenerates at \(E_{1}\), which proves \((a_{0})\) (1.3.2).\par
\DAsset{P0029-A01}{Theorem (8.1.15)(iii)-(v) and formulas (8.1.15.3), (8.1.16.1)-(2)}
\DAsset{P0029-A02}{proof (8.1.17), assertions \((a_{r})-(d_{r})\), and proof of \((a_{0})\)}

\DArticleAnchor{0030}{30}{31}{697}{34}{P05}

\DFormula{B)\,Proof\,of\,(b_{0}):\,one\,has}
\par\medskip
\DFormula{(_{L}\,E_{0}^{p\bullet }(sK_{\mathbb{C}}),F)\,=\,(K_{\mathbb{C}}^{\bullet p},F).}
\par\medskip
Apply (8.1.9)(v) and (1.3.2).\par
\par\medskip
C) Proof of \((c_{1})\): one must show that \(H^{n}(K_{A}^{\bullet p})\), equipped with the filtrations induced by Dec W on \(H^{n}(K_{A\otimes \mathbb{Q}}^{\bullet p})\) \(=\) \(H^{n}(K_{A}^{\bullet p})\) \(\otimes \) \(\mathbb{Q}\) and by F on \(H^{n}(K_{\mathbb{C}}^{\bullet p})\) \(=\) \(H^{n}(K_{A}^{\bullet p})\) \(\otimes _{A}\) \(\mathbb{C}\), is a mixed A-Hodge structure. Dec W induces on \(H^{n}(K_{A\otimes \mathbb{Q}}^{\bullet p})\) the same filtration as W[n], and the conclusion follows from (8.1.9)(ii).\par
\par\medskip
D) Proof of \((c_{r})\) \(+\) \((d_{r})\) \(\Rightarrow \) \((a_{r})\) \(+\) \((b_{r})\) \(+\) \((c_{r+1})\): this is the key point of the proof, and a simple application of (2.3.5), enlarged by (III.0.4).\par
\par\medskip
E) Proof of \((a_{r})_{r\leq r_{0}−2}\) \(+\) \((b_{r})_{r\leq r_{0}−2}\) \(+\) \((c_{r_{0}})\) \(\Rightarrow \) \((d_{r_{0}})\): by (1.3.16), the recurrent filtration and the direct filtrations induced by \(Dec_{1}W\) (respectively F) on \(_{L}\) \(E_{r}\) coincide for r \(\leq \) \(r_{0}\), and one applies (1.3.13)(i).\par
\par\medskip
This completes the proof by induction of (a), (b), (c), and (d).\par
\par\medskip
Assertion (iii) was proved in part E of the induction, and the assertions of (iv) are \((c_{r})\) and \((d_{r})\). We shall prove (8.1.15.3) in the form (8.1.16.2).\par
\par\medskip
That\par
\par\medskip
\DFormula{\operatorname{Gr}_{L}(H(sK_{A\otimes \mathbb{Q}}),Dec_{1}W)\,=\,(_{L}\,E_{\infty }(sK_{A\otimes \mathbb{Q}}),W)}
\par\medskip
follows from \((a_{r})\) and (1.3.17). That\par
\par\medskip
\DFormula{\operatorname{Gr}_{L}(H(sK_{\mathbb{C}}),F)\,=\,(_{L}\,E_{\infty }(sK_{\mathbb{C}}),F)}
\par\medskip
follows from \((b_{r})\) and (1.3.17). We conclude with the following lemma, whose verification is left to the reader:\par
\par\medskip
Lemma (8.1.18). — Let H \(=\) \((H_{A},W,F)\) be a mixed A-Hodge structure. For a filtration L of \(H_{A}\) to arise from a filtration of H in the abelian category of mixed A-Hodge structures, it is necessary and sufficient that, for every n, \((\operatorname{Gr}_{L}^{n}(H_{A}),\operatorname{Gr}_{L}^{n}(W),\operatorname{Gr}_{L}^{n}(F))\) be a mixed A-Hodge structure.\par
\par\medskip
(8.1.19) Apply (8.1.13) and (8.1.15) to Example (8.1.12).\par
\par\medskip
One obtains a mixed Hodge complex consisting of a), b), and c) below.\par
\par\medskip
a) The complex \(R\Gamma Rj_{\ast }\mathbb{Z}\) \(\simeq \) \(R\Gamma (U_{\bullet },\mathbb{Z})\), whose cohomology groups are the \(H^{n}(U_{\bullet },\mathbb{Z})\).\par
\par\medskip
b) The filtration \(\delta (W,L)\) on \(R\Gamma j_{\ast }\mathbb{Q}\) \(\simeq \) \(R\Gamma (U_{\bullet },\mathbb{Q})\). With the notation of (3.1.4), one has\par
\par\medskip
\DAsset{P0030-A02}{three-line decomposition of \(\operatorname{Gr}_{n}^{\delta (W,L)}R\Gamma (U_{\bullet },\mathbb{Q})\), through \(\oplus _{m}\) \(\operatorname{Gr}_{n+m}^{W}\) \(R\Gamma (U_{m},\mathbb{Q})[−m]\), then \(R\Gamma (\widetilde{Y}_{m}^{n+m},\varepsilon _{\mathbb{Q}}^{n+m})\) with shifts \([−n−m][−m]\) and \([−n−2m]\)}
\DAsset{P0030-A01}{proofs of \((b_{0})\), \((c_{1})\), steps D-E, and passage to the abutment}

\DArticleAnchor{0031}{31}{32}{698}{35}{P06}

The term \(E_{1}\) of the corresponding spectral sequence is a sum of groups of the form \(H^{p}(\widetilde{Y}_{q}^{r},\varepsilon _{\mathbb{Q}}^{r})\); the latter contributes to \(H^{p+q+r}(U_{\bullet })\); it corresponds to \(\operatorname{Gr}_{a}^{\delta (W,F)}\) for \(p+2r=(p+q+r)+a\):\par
\par\medskip
\DFormula{(8.1.19.1)\,_{\delta (W,L)}E_{1}^{−a,b}\,=\,⨁_{p+2r=b,\,q−r=−a}\,H^{p}(\widetilde{Y}_{q}^{r},\varepsilon _{\mathbb{Q}}^{r})\,\Rightarrow \,H^{−a+b}(U_{\bullet },\mathbb{Q}).}
\par\medskip
c) The bifiltration \((\delta (W,L),F)\) on \(R\Gamma Rj_{\ast }\mathbb{C}\) \(\simeq \) \(sR\Gamma ^{\bullet }\Omega _{X_{\bullet }}^{\bullet }(\operatorname{log}\) D).\par
\par\medskip
The filtration F equips \(_{\delta (W,L)}E_{1}^{−a,b}\) with a Hodge structure of weight b, and the differentials \(d_{1}\) of \(_{\delta (W,L)}E\) are morphisms of rational Hodge structures. The complex \(_{\delta (W,L)}E_{1}^{\bullet b}\) is the simple complex associated with the following double complex, whose components are Hodge structures of weight b:\par
\par\medskip
\DAsset{P0031-A01}{double complex of (8.1.19.2), with terms \(H^{b−2(r+2)}(\widetilde{Y}_{q+1}^{r+1},\varepsilon _{\mathbb{Q}}^{r+1})(−r−1)\), \(H^{b−2r}(\widetilde{Y}_{q+1}^{r},\varepsilon _{\mathbb{Q}}^{r})(−r)\), \(H^{b−2(r−2)}(\widetilde{Y}_{q+1}^{r−1},\varepsilon _{\mathbb{Q}}^{r−1})(−r+1)\), corresponding lower terms, Gysin arrows, and vertical arrows \(\sum _{i}(−1)^{i}\delta _{i}\)}
\par\medskip
The indices of the components in \(E_{1}^{ab}\) satisfy \(−a=q−r\).\par
\par\medskip
In particular, the groups \(H^{n}(U_{\bullet },\mathbb{Z})\) are thereby equipped with mixed Hodge structures (8.1.9).\par
\par\medskip
Proposition (8.1.20). — Under the preceding hypotheses:\par
\par\medskip
(i) The mixed Hodge structure on \(H^{n}(U_{\bullet },\mathbb{Z})\) is functorial in the pair \((U_{\bullet },X_{\bullet })\).\par
\par\medskip
(ii) In rational cohomology, the spectral sequence (8.1.19.1) degenerates at \(E_{2}\) \((E_{2}=E_{\infty })\) and abuts to the weight filtration on \(H^{n}(U_{\bullet },\mathbb{Q})\). The Hodge structure on \(E_{2}\) induced by that on \(E_{1}\) is also that on \(\operatorname{Gr}_{W}\) \(H^{n}(U_{\bullet },\mathbb{Q})\).\par
\par\medskip
(iii) The Hodge numbers of \(H^{n}(U_{\bullet },\mathbb{Z})\) satisfy\par
\par\medskip
\(h^{pq}\neq 0\) \(\Rightarrow \) \(0\leq p\leq n\) and \(0\leq q\leq n\).\par
\par\medskip
(iv) For \(Y_{\bullet }=\varnothing \), the Hodge numbers of \(H^{n}(X_{\bullet },\mathbb{Z})\) satisfy\par
\par\medskip
\(h^{pq}\neq 0\) \(\Rightarrow \) \(0\leq p\leq n\), \(0\leq q\leq n\), and \(p+q\leq n\).\par
\par\medskip
(i) and (ii) follow from the general theory (8.1.9). To prove (iii), inspect (8.1.19.2) and observe that the Hodge structures \(H^{p}(\widetilde{Y}_{q}^{r},\varepsilon _{\mathbb{Q}}^{r})(−r)\) \((p+q+r=n)\) satisfy (iii). To prove (iv), observe that the Hodge structures \(H^{p}(X_{q})\) \((p+q=n)\) satisfy (iv).\par
\par\medskip
(8.1.21) A cohomological mixed Hodge A-complex on an r-simplicial topological space is defined by transposing (8.1.11). If K is such a complex, one sees as in (8.1.13) that \(R\Gamma ^{\bullet }K\) is an r-cosimplicial mixed Hodge A-complex (8.1.10.1).\par
\DAsset{P0031-A02}{Proposition (8.1.20) and r-simplicial definition (8.1.21)}

\DArticleAnchor{0032}{32}{33}{699}{36}{P06}

Let \(X_{\bullet }\) be a proper smooth r-simplicial scheme. A normal-crossings divisor \(D_{\bullet }\) in \(X_{\bullet }\) is a system of normal-crossings divisors in the \(X_{n}\) such that \(X_{\bullet }−D_{\bullet }=U_{\bullet }\) is an r-simplicial subscheme of \(X_{\bullet }\). Let \(j:U_{\bullet }\hookrightarrow X_{\bullet }\). As in (8.1.12), one defines a cohomological mixed Hodge complex \(Rj_{\ast }\mathbb{Z}\) on \(X_{\bullet }\).\par
\par\medskip
As in (8.1.20), one could deduce from it a mixed Hodge structure on \(H^{\bullet }(U_{\bullet },\mathbb{Z})\), but, since\par
\par\medskip
\DFormula{H^{\bullet }(U_{\bullet },\mathbb{Z})=H^{\bullet }(\delta U_{\bullet },\mathbb{Z}),}
\par\medskip
this would teach us nothing new (the preceding equality would be an identity between mixed Hodge structures).\par
\par\medskip
(8.1.22) Set \(r=2\). For every bounded-below bicosimplicial complex \(K=(K^{p\,n_{1}\,n_{2}})\) \((p=differential\) degree, \(n_{i}=i-th\) simplicial degree), let \(s_{1}K\) denote the cosimplicial complex of simplicial degree \(n_{2}\) given by\par
\par\medskip
\DFormula{(s_{1}K)^{\bullet \,n_{2}}=sK^{\bullet \bullet \,n_{2}}.}
\par\medskip
Let \(L_{1}\) denote the filtration on \(s_{1}K\) which, on each \(sK^{\bullet \bullet \,n_{2}}\), induces the filtration by the simplicial degree \(n_{1}\).\par
\par\medskip
Let K be a bicosimplicial mixed Hodge A-complex. It follows from (8.1.15) that \((s_{1}K,\delta (W,L_{1}),F)\) is a cosimplicial mixed Hodge complex. Let \(L_{2}\) be the filtration on \(sK=ss_{1}K\) by the simplicial degree \(n_{2}\). By (8.1.14) and (8.1.15), the spectral sequence defined by \(L_{2}\) comes from a spectral sequence in the abelian category of mixed Hodge A-structures:\par
\par\medskip
\DFormula{(8.1.22.1)\,_{L_{2}}E_{1}^{pq}=H^{q}(sK^{\bullet \bullet p})\,\Rightarrow \,H^{p+q}(sK).}
\par\medskip
In this spectral sequence, the mixed Hodge A-structure on \(E_{1}^{pq}\) is the mixed Hodge A-structure ((8.1.9)(ii)) on the cohomology of the mixed Hodge A-complex \((sK^{\bullet \bullet p},\delta (W,L_{1}),F)\). The mixed Hodge A-structure on the abutment is the mixed Hodge A-structure on the cohomology of the mixed Hodge A-complex \((sK,\delta (\delta (W,L_{1}),L_{2}),F)\). The filtration \(\delta (\delta (W,L_{1}),L_{2})=_{dfn}\delta (W,L_{1},L_{2})\) is given by\par
\par\medskip
\DFormula{\delta (W,L_{1},L_{2})_{k}(K^{p\,n_{1}\,n_{2}})=W_{k+n_{1}+n_{2}}K^{p\,n_{1}\,n_{2}}.}
\par\medskip
(8.1.23) In particular, a cohomological mixed Hodge A-complex on a bisimplicial topological space \(X_{\bullet \bullet }\) defines a spectral sequence of mixed Hodge A-structures\par
\par\medskip
\DFormula{(8.1.23.1)\,E_{1}^{pq}=H^{q}(X_{\bullet p},K)\,\Rightarrow \,H^{p+q}(X_{\bullet \bullet },K).}
\par\medskip
In the particular case considered in (8.1.21), for \(r=2\) one has \(A=\mathbb{Z}\), and (8.1.23.1) becomes\par
\par\medskip
\DFormula{(8.1.23.2)\,E_{1}^{pq}=H^{q}(U_{\bullet p},\mathbb{Z})\,\Rightarrow \,H^{p+q}(U_{\bullet \bullet },\mathbb{Z}).}
\DAsset{P0032-A01}{r-simplicial cohomology and bicosimplicial construction (8.1.22)}
\DAsset{P0032-A02}{spectral sequences (8.1.22.1), (8.1.23.1), and (8.1.23.2)}

\DArticleAnchor{0033}{33}{34}{700}{37}{P06}

(8.1.24) If K′ and K″ are two mixed Hodge A-complexes, their tensor product \(K′\otimes K″\), defined by\par
\par\medskip
\DFormula{(K′\otimes K″)_{A}=K′_{A}\otimes ^{L}\,K″_{A},}
\par\medskip
\DFormula{((K′\otimes K″)_{A\otimes \mathbb{Q}},W)=(K′_{A\otimes \mathbb{Q}},W′)\otimes (K″_{A\otimes \mathbb{Q}},W″),}
\par\medskip
\DFormula{((K′\otimes K″)_{\mathbb{C}},W,F)=(K′_{\mathbb{C}},W′,F′)\otimes (K″_{\mathbb{C}},W″,F″),}
\par\medskip
is again a mixed Hodge A-complex.\par
\par\medskip
Let U′ (respectively U″) be the complement of a normal-crossings divisor Y′ (respectively Y″) in a proper smooth scheme X′ (respectively X″). Then \(U=U′\times U″\) is the complement of a normal-crossings divisor Y in \(X=X′\times X″\). Let j (respectively j′,j″) be the inclusion of U (respectively U′,U″) in X (respectively X′,X″).\par
\par\medskip
There is a quasi-isomorphism\par
\par\medskip
\(Rj′_{\ast }\mathbb{Z}\otimes ^{L}\) \(Rj″_{\ast }\mathbb{Z}\) \(=_{dfn}\) \(\operatorname{pr}_{1}^{\ast }Rj′_{\ast }\mathbb{Z}\otimes ^{L}\) \(\operatorname{pr}_{2}^{\ast }Rj″_{\ast }\mathbb{Z}\) \(\simeq \) \(Rj_{\ast }\mathbb{Z}\),\par
\par\medskip
a filtered quasi-isomorphism\par
\par\medskip
\DFormula{(Rj′_{\ast }\mathbb{Q},\tau _{\leq })\otimes ^{L}(Rj″_{\ast }\mathbb{Q},\tau _{\leq })\,\simeq \,(Rj_{\ast }\mathbb{Q},\tau _{\leq }),}
\par\medskip
and a bifiltered morphism\par
\par\medskip
\DFormula{(\Omega _{X′}^{\bullet }(\operatorname{log}\,Y′),W,F)\otimes (\Omega _{X″}^{\bullet }(\operatorname{log}\,Y″),W,F)\,\to \,(\Omega _{X}^{\bullet }(\operatorname{log}\,Y),W,F).}
\par\medskip
Apply the functor \(R\Gamma \). By the Künneth formula and its analogue for the cohomology of coherent analytic sheaves, one obtains an isomorphism of mixed Hodge complexes\par
\par\medskip
\DFormula{R\Gamma (U′,\mathbb{Z})\otimes ^{L}\,R\Gamma (U″,\mathbb{Z})\,\simeq \,R\Gamma (U,\mathbb{Z}),}
\par\medskip
and hence an isomorphism of rational mixed Hodge structures\par
\par\medskip
\DFormula{H^{\bullet }(U′,\mathbb{Q})\otimes H^{\bullet }(U″,\mathbb{Q})\,\simeq \,H^{\bullet }(U,\mathbb{Q}).}
\par\medskip
(8.1.25) Standard arguments based on the Cartier-Eilenberg-Zilber theorem give the following simplicial variant of (8.1.24).\par
\par\medskip
Let \(U′_{\bullet }\) (respectively \(U″_{\bullet })\) be the complement of a normal-crossings divisor in a proper smooth simplicial scheme \(X′_{\bullet }\) (respectively \(X″_{\bullet })\). Let \(U_{\bullet }=U′_{\bullet }\times U″_{\bullet }\) and \(X_{\bullet }=X′_{\bullet }\times X″_{\bullet }\).\par
\par\medskip
There is an isomorphism of mixed Hodge complexes\par
\par\medskip
\DFormula{R\Gamma (U′_{\bullet },\mathbb{Z})\otimes ^{L}\,R\Gamma (U″_{\bullet },\mathbb{Z})\,\simeq \,R\Gamma (U_{\bullet },\mathbb{Z}),}
\par\medskip
and in particular an isomorphism of rational mixed Hodge structures\par
\par\medskip
\DFormula{H^{\bullet }(U′_{\bullet },\mathbb{Q})\otimes H^{\bullet }(U″_{\bullet },\mathbb{Q})\,\simeq \,H^{\bullet }(U_{\bullet },\mathbb{Q}).}
\par\medskip
(8.1.26) Let \(U′_{\bullet \bullet }\) (respectively \(U″_{\bullet \bullet })\) be the complement of a normal-crossings divisor in a proper smooth bisimplicial scheme \(X′_{\bullet \bullet }\) (respectively \(X″_{\bullet \bullet })\). Let \(U_{\bullet \bullet }=U′_{\bullet \bullet }\times U″_{\bullet \bullet }\) and\par
\DAsset{P0033-A01}{tensor product of mixed Hodge complexes and Künneth quasi-isomorphisms}
\DAsset{P0033-A02}{simplicial and bisimplicial variants (8.1.25)-(8.1.26)}

\DArticleAnchor{0034}{34}{35}{701}{38}{P06}

\(X_{\bullet \bullet }=X′_{\bullet \bullet }\times X″_{\bullet \bullet }\). Standard arguments show that, in rational cohomology, the spectral sequence (8.1.23.2) for \(U_{\bullet \bullet }\) is the tensor product of the spectral sequences (8.1.23.2) for \(U′_{\bullet \bullet }\) and \(U″_{\bullet \bullet }\).\par
\par\medskip
8.2. Separated algebraic spaces.\par
\par\medskip
(8.2.1) Let X be a scheme (or algebraic space) of finite type over \(\mathbb{C}\), assumed separated. There then exists a diagram\par
\par\medskip
\DAsset{P0034-A01}{diagram (8.2.1.1): open immersion \(Y_{\bullet }\hookrightarrow \overline{Y}_{\bullet }\) and proper augmentation \(\alpha :Y_{\bullet }\to X\)}
\par\medskip
(8.2.1.1)\par
\par\medskip
in which \(Y_{\bullet }\) is the complement of a normal-crossings divisor in a proper smooth simplicial scheme \(\overline{Y}_{\bullet }\), and is a proper hypercovering of X. One has (5.3.5(II))\par
\par\medskip
\DFormula{(8.2.1.2)\,H^{\bullet }(X)\,\simeq \,H^{\bullet }(Y_{\bullet }).}
\par\medskip
In (8.1.19), \(H^{\bullet }(Y_{\bullet })\) was equipped with a mixed Hodge structure. Let \(\mathcal{H}(\alpha )\) denote the mixed Hodge structure on \(H^{\bullet }(X)\) obtained from it through (8.2.1.2).\par
\par\medskip
Proposition (8.2.2). — With the preceding notation, the mixed Hodge structure \(\mathcal{H}(\alpha )\) on \(H^{\bullet }(X)\) is independent of the choice of \(\alpha \). It is called the mixed Hodge structure on \(H^{\bullet }(X,\mathbb{Z})\). For every morphism of schemes \(f:X_{1}\to X_{2}\), \(f^{\bullet }:H^{\bullet }(X_{2})\to H^{\bullet }(X_{1})\) is a morphism of mixed Hodge structures.\par
\par\medskip
The proof uses (6.2.8) and is parallel to that of (3.2.11 C)) (cf. also (8.3.3)).\par
\par\medskip
Remark (8.2.3). — When X is smooth, by taking \(\overline{Y}_{\bullet }\) to be a constant simplicial scheme one checks that the mixed Hodge structure (8.2.2) coincides with that defined in (3.2.12).\par
\par\medskip
Theorem (8.2.4). — The pairs (p,q) for which the Hodge number \(h^{pq}\) of \(H^{n}(X)\) is nonzero satisfy the following conditions:\par
\par\medskip
\DFormula{(i)\,(p,q)\in [0,n]\times [0,n].}
\par\medskip
(ii) If dim \(X=N\) and \(n\geq N\), then\par
\par\medskip
\DFormula{(p,q)\in [n−N,N]\times [n−N,N].}
\par\medskip
\DFormula{(iii)\,If\,X\,is\,proper,\,then\,p+q\leq n.}
\par\medskip
(iv) If X is smooth, then \(p+q\geq n\). The same conclusion holds when X is an equidimensional “rational homology manifold” of dimension N, i.e. when, for every \(x\in X\), \(H^{i}_{\lbrace x\rbrace }(X,\mathbb{Q})=0\) if \(i\neq 2N\) and \(=\mathbb{Q}\) if \(i=2N\).\par
\DAsset{P0034-A02}{Proposition (8.2.2), Remark (8.2.3), and Theorem (8.2.4)}

\DArticleAnchor{0035}{35}{36}{702}{39}{P06}

\DAsset{P0035-A01}{two diagrams of the admissible (p,q)-regions: for \(n\leq N\), the square \([0,n]^{2}\) divided by \(p+q=n\), with “proper” below the diagonal and “smooth” above; for \(n\geq N\), the square \([n−N,N]^{2}\) with the same regions, axes p,q, and marks N,n}
\par\medskip
Assertion (ii) will be proved in (8.3.10). Assertion (i) follows from (8.1.20(iii)). For X proper, one can find a diagram (8.2.1.1) in which \(Y_{\bullet }=\overline{Y}_{\bullet }\), and (iii) follows from (8.1.20(iv)). For X smooth, (iv) follows from (8.2.3) and (3.2.15). The general case follows: if \(p:\widetilde{X}\to X\) is a resolution of the singularities of X, \(p^{\ast }:H^{\bullet }(X,\mathbb{Q})\to H^{\bullet }(\widetilde{X},\mathbb{Q})\) is injective, because the transpose under Poincaré duality, \(p_{!}\), of \(p^{\ast }:H_{c}^{\bullet }(X,\mathbb{Q})\to H_{c}^{\bullet }(\widetilde{X},\mathbb{Q})\) is a retraction of it.\par
\par\medskip
Proposition (8.2.5). — Suppose X is proper. If \(u:Y\to X\) is a proper surjective morphism, with Y smooth, then the quotient of weight n of \(H^{n}(X,\mathbb{Q})\) is the image of \(H^{n}(X,\mathbb{Q})\) in \(H^{n}(Y,\mathbb{Q})\). The sequence\par
\par\medskip
\DFormula{(8.2.5.1)\,H^{n}(X,\mathbb{Q})\,\to \,H^{n}(Y,\mathbb{Q})\,\rightrightarrows _{\operatorname{pr}_{1}^{\ast }−\operatorname{pr}_{2}^{\ast }}\,H^{n}(Y\times _{X}\,Y,\mathbb{Q})}
\par\medskip
is exact.\par
\par\medskip
Let \(Y_{\bullet }\) be as in (8.2.1.1), with \(Y_{0}=Y\) and \(Y_{\bullet }=\overline{Y}_{\bullet }\). By ((8.1.20(ii))), there is an exact sequence\par
\par\medskip
\DFormula{(8.2.5.2)\,0\,\to \,H^{n}(X,\mathbb{Q})/W^{n−1}H^{n}(X,\mathbb{Q})\,\to \,H^{n}(Y_{0},\mathbb{Q})\,\rightrightarrows _{\delta _{0}^{\ast },\delta _{1}^{\ast }}\,H^{n}(Y_{1},\mathbb{Q})}
\par\medskip
which implies (8.2.5) (conversely, this exact sequence is readily deduced from (8.2.5) and (8.2.7) below).\par
\par\medskip
Proposition (8.2.6). — Let there be morphisms of schemes\par
\par\medskip
\DFormula{Y\,\xrightarrow{f}\,X\,\overset{j}{\hookrightarrow}\,\overline{X}.}
\par\medskip
Assume that Y is proper, X is smooth, and \(\overline{X}\) is a smooth compactification of X. Then \(H^{\bullet }(\overline{X},\mathbb{Q})\) and \(H^{\bullet }(X,\mathbb{Q})\) have the same image in \(H^{\bullet }(Y,\mathbb{Q})\).\par
\par\medskip
The proof is parallel to that of (3.2.18). It is enough to prove that, for all i and n, \(\operatorname{Gr}_{i}^{W}(f^{\ast })\) and \(\operatorname{Gr}_{i}^{W}((jf)^{\ast })\) have the same image in \(\operatorname{Gr}_{i}^{W}\) \(H^{n}(Y,\mathbb{Q})\).\par
\par\medskip
a) For \(i>n\), these images are zero because \(\operatorname{Gr}_{i}^{W}\) \(H^{n}(Y,\mathbb{Q})=0\) (8.2.4(iii)).\par
\par\medskip
b) For \(i<n\), these images are zero because \(\operatorname{Gr}_{i}^{W}\) \(H^{n}(X,\mathbb{Q})=0\) (8.2.4(iv)).\par
\par\medskip
c) For \(i=n\), these images are equal because \(\operatorname{Gr}_{n}^{W}(j^{\ast })\) is surjective (3.2.17).\par
\DAsset{P0035-A02}{Propositions (8.2.5)-(8.2.6) and exact sequences}

\DArticleAnchor{0036}{36}{37}{703}{40}{P06}

Proposition (8.2.7). — Let there be morphisms of schemes\par
\par\medskip
\DFormula{\widetilde{X}\,\xrightarrow{\pi }\,X\,\xrightarrow{f}\,Y.}
\par\medskip
Assume that Y is smooth, X is proper, \(\widetilde{X}\) is proper and smooth, and \(\pi \) is surjective. Then the kernels of \(f^{\ast }\) and \((f\pi )^{\ast }\) in \(H^{\bullet }(Y,\mathbb{Q})\) are equal.\par
\par\medskip
Let us prove that Ker \(\operatorname{Gr}_{i}^{W}(f^{\ast })=\operatorname{Ker}(\operatorname{Gr}_{i}^{W}((f\pi )^{\ast }))\) in \(\operatorname{Gr}_{i}^{W}(H^{n}(Y,\mathbb{Q}))\).\par
\par\medskip
a) For \(i<n\), these kernels are zero because \(\operatorname{Gr}_{i}^{W}\) \(H^{n}(Y,\mathbb{Q})=0\) ((8.2.4(iv))).\par
\par\medskip
b) For \(i>n\), these kernels are all of \(\operatorname{Gr}_{i}^{W}(H^{n}(Y,\mathbb{Q}))\) because \(\operatorname{Gr}_{i}^{W}\) \(H^{n}(X,\mathbb{Q})=0\) ((8.2.4(iii))).\par
\par\medskip
c) For \(i=n\), these kernels are equal because \(\operatorname{Gr}_{n}^{W}(\pi ^{\ast })\) is injective (8.2.5.2).\par
\par\medskip
Corollary (8.2.8). — Let there be morphisms of schemes\par
\par\medskip
\DFormula{\widetilde{X}\,\xrightarrow{\pi }\,X\,\overset{f}{\hookrightarrow}\,Y.}
\par\medskip
Assume that Y is proper and smooth, X is a closed subscheme of Y, and \(\widetilde{X}\) is a resolution of the singularities of X. Let \(q=f\pi \) and \(U=Y−X\). Then the sequence\par
\par\medskip
\DFormula{(8.2.8.1)\,H^{\bullet }(\widetilde{X},\mathbb{Q})\,\xrightarrow{q_{!}}\,H^{\bullet }(Y,\mathbb{Q})\,\to \,H^{\bullet }(U,\mathbb{Q})}
\par\medskip
is exact.\par
\par\medskip
By (8.2.7), one has\par
\par\medskip
\DFormula{\operatorname{Ker}(f^{\ast }:H^{\bullet }(Y,\mathbb{Q})\to H^{\bullet }(X,\mathbb{Q}))=\operatorname{Ker}(q^{\ast }:H^{\bullet }(Y,\mathbb{Q})\to H^{\bullet }(\widetilde{X},\mathbb{Q})).}
\par\medskip
The long exact cohomology sequence of the pair (Y,X) therefore gives an exact sequence\par
\par\medskip
\DFormula{(8.2.8.2)\,H_{c}^{\bullet }(U,\mathbb{Q})\,\to \,H^{\bullet }(Y,\mathbb{Q})\,\to \,H^{\bullet }(\widetilde{X},\mathbb{Q}),}
\par\medskip
and (8.2.8.1) is the transpose of (8.2.8.2) under Poincaré duality.\par
\par\medskip
Remark (8.2.9). — Corollary (8.2.8) is used under more general hypotheses in ([8], (9.7), p. 161). As Lieberman pointed out to the author, the justification given in loc. cit. is insufficient.\par
\par\medskip
It follows from (8.1.25) that:\par
\par\medskip
Proposition (8.2.10). — The Künneth isomorphisms\par
\par\medskip
\DFormula{H^{\bullet }(X\times Y,\mathbb{Q})\,\simeq \,H^{\bullet }(X,\mathbb{Q})\otimes H^{\bullet }(Y,\mathbb{Q})}
\par\medskip
are isomorphisms of mixed Hodge structures.\par
\par\medskip
Corollary (8.2.11). — The cup product\par
\par\medskip
\DFormula{H^{\bullet }(X)\otimes H^{\bullet }(X)\,\to \,H^{\bullet }(X)}
\par\medskip
is a morphism of mixed Hodge structures.\par
\par\medskip
This follows from (8.2.10) and from (8.2.2) applied to the diagonal map \(\Delta :X\to X\times X\).\par
\DAsset{P0036-A01}{Proposition (8.2.7), Corollary (8.2.8), and morphism diagrams}
\DAsset{P0036-A02}{sequences (8.2.8.1)-(2), Remark (8.2.9), Künneth, and cup product}

\DArticleAnchor{0037}{37}{38}{704}{41}{P07}

8.3. Hodge theory of simplicial schemes.\par
\par\medskip
(8.3.1) For \(X_{\bullet }\) a separated simplicial scheme (or algebraic space), we shall have to consider diagrams of simplicial schemes\par
\par\medskip
\DAsset{P0037-A01}{diagram (8.3.1.1): \(Z_{\bullet }\) \(\overset{j}{\hookrightarrow}\) \(\overline{Z}_{\bullet }\) and \(p:Z_{\bullet }\to X_{\bullet }\)}
\par\medskip
satisfying:\par
\par\medskip
a) j is an open immersion with dense image; \(\overline{Z}_{\bullet }\) is proper.\par
\par\medskip
b) The evident morphisms\par
\par\medskip
\DFormula{Z_{n}\,\to \,(\operatorname{cosq}_{n−1}^{X_{\bullet }}\,\operatorname{sq}_{n−1}\,Z_{\bullet })_{n}}
\par\medskip
are proper and surjective.\par
\par\medskip
For \(\overline{Z}_{\bullet }\) reduced, one may regard j as a simplicial object of the category \(\mathcal{C}\) of pairs \((Y,\overline{Y})\) formed by a reduced proper scheme \(\overline{Y}\) and a dense (Zariski) open subset Y of \(\overline{Y}\). The category \(\mathcal{C}\) satisfies (6.2.4.1), which allows the analogue of (6.2.4) to be applied to it.\par
\par\medskip
Proposition (8.3.2). — (i) For every separated simplicial scheme \(X_{\bullet }\), there exists a diagram (8.3.1.1) satisfying a), b), such that \(\overline{Z}_{\bullet }\) is proper and smooth and \(Z_{\bullet }\) is the complement of a normal-crossings divisor.\par
\par\medskip
(ii) Let \(F:X_{\bullet 1}\to X_{\bullet 2}\) be a morphism of separated simplicial schemes. Let \((X_{\bullet i},Z_{\bullet i},\overline{Z}_{\bullet i})_{i=1,2}\) be two diagrams (8.3.1.1) satisfying a) and b). There then exist \((X_{\bullet 1},Z_{\bullet },\overline{Z}_{\bullet })\) as in (i) and a commutative diagram\par
\par\medskip
\DAsset{P0037-A02}{diagram of (8.3.2)(ii), with \(Z_{\bullet }\hookrightarrow \overline{Z}_{\bullet }\) dominating the two systems and \(f:X_{\bullet 1}\to X_{\bullet 2}\)}
\par\medskip
The proof, analogous to but more complicated than (6.2.5), is left to the reader.\par
\par\medskip
(8.3.3) Let there be a diagram (8.3.1.1) as in (8.3.2(i)). By (5.3.5(V)), the map \(H^{\bullet }(X_{\bullet },\mathbb{Z})\to H^{\bullet }(Z_{\bullet },\mathbb{Z})\) is an isomorphism. In (8.1.19) and (8.1.20), we endowed \(H^{\bullet }(Z_{\bullet },\mathbb{Z})\) with a mixed Hodge structure, hence a structure of\par

\DArticleAnchor{0038}{38}{39}{705}{42}{P07}

mixed Hodge structure on \(H^{\bullet }(X_{\bullet },\mathbb{Z})\). It is independent of the choice of diagram (8.3.1.1): by (8.3.2(ii)) with \(f=Id\), two systems \((Z_{\bullet i},\overline{Z}_{\bullet i},p_{i})_{i=1,2}\) are always dominated by a third \((Z_{\bullet },\overline{Z}_{\bullet },p)\), whence a commutative diagram of isomorphisms\par
\par\medskip
\DAsset{P0038-A01}{independence diagram: \(H^{\bullet }(X_{\bullet },\mathbb{Z})\) mapping to \(H^{\bullet }(Z_{\bullet 1},\mathbb{Z})\), \(H^{\bullet }(Z_{\bullet },\mathbb{Z})\), and \(H^{\bullet }(Z_{\bullet 2},\mathbb{Z})\), with horizontal isomorphisms toward \(H^{\bullet }(Z_{\bullet },\mathbb{Z})\)}
\par\medskip
Definition (8.3.4). — The mixed Hodge structure on \(H^{\bullet }(X_{\bullet },\mathbb{Z})\) is the one defined above.\par
\par\medskip
It follows from (8.3.2)(ii) that this mixed Hodge structure is functorial in \(X_{\bullet }\).\par
\par\medskip
Proposition (8.3.5). — Let \(X_{\bullet }\) be a separated simplicial scheme. The spectral sequence\par
\par\medskip
\DFormula{E_{1}^{pq}=H^{q}(X_{p},\mathbb{Z})\,\Rightarrow \,H^{p+q}(X_{\bullet },\mathbb{Z})}
\par\medskip
is a spectral sequence of mixed Hodge structures.\par
\par\medskip
For a diagram as in (8.3.1), the maps \(p_{n}:Z_{n}\to X_{n}\) are proper and surjective. One may therefore apply (6.2.4) to the category of simplicial objects of the category \(\mathcal{C}\) (8.3.1), obtaining as in (6.2.5) the following lemma:\par
\par\medskip
Lemma (8.3.6). — Let \(X_{\bullet }\) be a separated simplicial scheme. There exist a bisimplicial scheme \(Z_{\bullet \bullet }\) augmented to \(X_{\bullet }\),\par
\par\medskip
\DFormula{\varepsilon _{n,m}:Z_{n,m}\to X_{n},}
\par\medskip
and \(i:Z_{\bullet \bullet }\hookrightarrow \overline{Z}_{\bullet \bullet }\), such that:\par
\par\medskip
a) \(\overline{Z}_{\bullet \bullet }\) is a proper smooth bisimplicial scheme, and \(Z_{\bullet \bullet }\) is the complement of a normal-crossings divisor in \(\overline{Z}_{\bullet \bullet }\);\par
\par\medskip
b) for every n, the augmented simplicial scheme \(Z_{n\bullet }\to X_{n}\) is a proper hypercover of \(X_{n}\).\par
\par\medskip
In (8.1.2.1), we endowed \(H^{\bullet }(Z_{\bullet \bullet },\mathbb{Z})\) with a mixed Hodge structure. If \(\delta Z_{\bullet \bullet }\) is the diagonal simplicial scheme of \(Z_{\bullet \bullet }\), one has a commutative diagram of morphisms\par
\par\medskip
\DAsset{P0038-A02}{triangular diagram: \(H^{\bullet }(X_{\bullet },\mathbb{Z})\) \(\xrightarrow{\gamma }\) \(H^{\bullet }(Z_{\bullet \bullet },\mathbb{Z})\), \(\alpha \) toward \(H^{\bullet }(\delta Z_{\bullet \bullet },\mathbb{Z})\), and \(\beta \) between \(H^{\bullet }(Z_{\bullet \bullet },\mathbb{Z})\) and \(H^{\bullet }(\delta Z_{\bullet \bullet },\mathbb{Z})\)}

\DArticleAnchor{0039}{39}{40}{706}{43}{P07}

in which \(\beta \) and \(\gamma \) are isomorphisms, and \(\alpha \) and \(\beta \) are morphisms of mixed Hodge structures; \(\gamma \) is therefore an isomorphism of mixed Hodge structures.\par
\par\medskip
There is a morphism of spectral sequences\par
\par\medskip
\DAsset{P0039-A01}{morphism from \(H^{q}(X_{p},\mathbb{Z})\Rightarrow H^{p+q}(X_{\bullet },\mathbb{Z})\) (5.2.3) to \(H^{q}(Z_{p\bullet },\mathbb{Z})\Rightarrow H^{p+q}(Z_{\bullet \bullet },\mathbb{Z})\) (8.1.23.1)}
\par\medskip
and this morphism is an isomorphism compatible with the mixed Hodge structures, so that (8.3.5) follows from (8.1.23).\par
\par\medskip
Example (8.3.7). — Let \(X_{\bullet }\) be a simplicial scheme that is not necessarily separated. There then exist a separated simplicial scheme \(Y_{\bullet }\) and \(u:Y_{\bullet }\to X_{\bullet }\) such that, for \(n\geq 0\), the map\par
\par\medskip
\DFormula{u′_{n}:Y_{n}\to (\operatorname{cosq}_{n−1}^{X_{\bullet }}\,\operatorname{sq}_{n−1}Y_{\bullet })_{n}}
\par\medskip
is étale and surjective (for example, the disjoint sum of the open subsets in an open covering).\par
\par\medskip
The map \(u^{\ast }:H^{\bullet }(X_{\bullet },\mathbb{Z})\to H^{\bullet }(Y_{\bullet },\mathbb{Z})\) is an isomorphism ((5.3.5)(V) and (III)). By definition, the mixed Hodge structure on \(H^{\bullet }(X_{\bullet },\mathbb{Z})\) is deduced from that on \(H^{\bullet }(Y_{\bullet },\mathbb{Z})\) via \(u^{\ast }\). It does not depend on the choice of \(Y_{\bullet }\).\par
\par\medskip
For X a nonseparated scheme, the mixed Hodge structure on \(H^{\bullet }(X,\mathbb{Z})\) is that on \(H^{\bullet }(X_{\bullet },\mathbb{Z})\), where \(X_{\bullet }\) is the constant simplicial scheme defined by X. It is likely, but not proved, that (8.2.4)(iv) remains valid for X smooth and nonseparated.\par
\par\medskip
Example (8.3.8). — Let \(u:X_{\bullet }\to Y_{\bullet }\) be a morphism of simplicial schemes. The relative cohomology groups \(H^{n}(Y_{\bullet }\) mod \(X_{\bullet },\mathbb{Z})\) are, by definition (6.3.2), the cohomology groups of the simplicial scheme C(u). They therefore carry a mixed Hodge structure.\par
\par\medskip
Proposition (8.3.9). — The long exact sequence of relative cohomology\par
\par\medskip
\DFormula{H^{n}(Y_{\bullet },\mathbb{Z})\,\to \,H^{n}(X_{\bullet },\mathbb{Z})\,\xrightarrow{j}\,H^{n+1}(Y_{\bullet }\,\operatorname{mod}\,X_{\bullet },\mathbb{Z})\,\to \,H^{n+1}(Y,\mathbb{Z})}
\par\medskip
is a long exact sequence of mixed Hodge structures.\par
\par\medskip
We reduce to the case in which \(X_{\bullet }\) and \(Y_{\bullet }\) are smooth and, moreover, there exists a commutative diagram\par
\par\medskip
\DAsset{P0039-A02}{square \(X_{\bullet }\) \(\xrightarrow{u}\) \(Y_{\bullet }\), with immersions into \(\overline{X}_{\bullet }\) and \(\overline{Y}_{\bullet }\) and lower map \(\overline{u}:\overline{X}_{\bullet }\to \overline{Y}_{\bullet }\)}
\par\medskip
with \(\overline{X}_{\bullet }\) and \(\overline{Y}_{\bullet }\) proper and smooth, and \(X_{\bullet }\) and \(Y_{\bullet }\) the complements of normal-crossings divisors \(D_{\bullet }\) and \(E_{\bullet }\) in \(\overline{X}_{\bullet }\) and \(\overline{Y}_{\bullet }\).\par

\DArticleAnchor{0040}{40}{41}{707}{44}{P07}

There then exist bifiltered resolutions (whose associated graded has acyclic components)\par
\par\medskip
\(a:\Omega _{\overline{X}_{\bullet }}^{\bullet }(\operatorname{log}\) \(D_{\bullet })\to K\)   and   \(b:\Omega _{\overline{Y}_{\bullet }}^{\bullet }(\operatorname{log}\) \(E_{\bullet })\to L\)\par
\par\medskip
which fit into a commutative diagram of bifiltered complexes\par
\par\medskip
\DAsset{P0040-A01}{bifiltered square \(\Omega _{\overline{X}_{\bullet }}^{\bullet }(\operatorname{log}\) \(D_{\bullet })\) \(\xleftarrow{\overline{u}^{\ast }}\) \(\Omega _{\overline{Y}_{\bullet }}^{\bullet }(\operatorname{log}\) \(E_{\bullet })\), vertically to K \(\xleftarrow{v}\) L, followed by the distinguished triangle \(C(\Gamma v)\), \(\Gamma L\), and \(\Gamma K\)}
\par\medskip
\((\overline{u}^{\ast }\) and v are \(\overline{u}-morphisms)\). The long exact sequence (8.3.9) is then identified with the exact sequence deduced from the distinguished triangle represented in P0040-A01.\par
\par\medskip
The complex C(v) on \(C(\overline{u})\) is a bifiltered resolution of \(C(\overline{u}^{\ast }:\Omega _{\overline{Y}_{\bullet }}^{\bullet }(\operatorname{log}\) \(E_{\bullet })\to \Omega _{\overline{X}_{\bullet }}^{\bullet }(\operatorname{log}\) \(D_{\bullet }))\). It follows that the Hodge and weight filtrations on \(H^{\bullet }(Y_{\bullet }\) mod \(X_{\bullet },\mathbb{C})\) are deduced from the filtrations on the cone \(C(\Gamma v)\), and (8.3.9) follows formally.\par
\par\medskip
(8.3.10) Proof of (8.2.4)(ii).\par
\par\medskip
The proof proceeds by induction on \(N=\operatorname{dim}(X)\). We may assume X reduced. Let U be a dense smooth Zariski-open subset of X and \(Y=X−U\). By resolution of singularities, there exists a proper birational morphism \(p:X′\to X\) with smooth source that induces an isomorphism of \(p^{−1}(U)\) with U. Let \(Y′=p^{−1}(Y)\). For \(n\geq N\), consider the diagram of mixed Hodge structures\par
\par\medskip
\DAsset{P0040-A02}{exact diagram: \(H^{n}(X\) mod \(Y,\mathbb{Z})\to H^{n}(X,\mathbb{Z})\to H^{n}(Y,\mathbb{Z})\), vertically to \(H^{n−1}(Y′,\mathbb{Z})\to H^{n}(X′\) mod \(Y′,\mathbb{Z})\to H^{n}(X′,\mathbb{Z})\to H^{n}(Y′,\mathbb{Z})\), with \(p^{\ast }\) marked}
\par\medskip
Let \(\Phi \) be the family of supports on U formed by the \(A\subset U\) such that, in X (or X′, which is equivalent), \(\overline{A}\subset U\). One has\par
\par\medskip
\DFormula{H^{n}(X\,\operatorname{mod}\,Y,\mathbb{Z})=H^{n}_{\Phi }(U,\mathbb{Z})=H^{n}(X′\,\operatorname{mod}\,Y′,\mathbb{Z}),}
\par\medskip
so the arrow marked \(p^{\ast }\) is an isomorphism.\par
\par\medskip
For \(h^{pq}\) a nonzero Hodge number of one of the mixed Hodge structures under consideration, one has:\par

\DArticleAnchor{0041}{41}{42}{708}{45}{P07}

\DFormula{a)\,for\,H^{n}(X′,\mathbb{Z}):\,(p,q)\in [n−N,N]\times [n−N,N]}
\par\medskip
\DFormula{b)\,for\,H^{n−1}(Y′,\mathbb{Z}):\,(p,q)\in [n−N,N−1]\times [n−N,N−1]}
\par\medskip
\DFormula{c)\,for\,H^{n}(X′\,\operatorname{mod}\,Y′,\mathbb{Z})=H^{n}(X\,\operatorname{mod}\,Y,\mathbb{Z}):\,(p,q)\in [n−N,N]\times [n−N,N]}
\par\medskip
\DFormula{d)\,for\,H^{n}(Y,\mathbb{Z}):\,(p,q)\in [n−N+1,N−1]\times [n−N+1,N−1]}
\par\medskip
\DFormula{e)\,for\,H^{n}(X,\mathbb{Z}):\,(p,q)\in [n−N,N]\times [n−N,N].}
\par\medskip
a) follows because X′ is smooth and the assertion is true for the terms \(_{W}\) \(E_{1}\) of the spectral sequence (3.2.4.1). b) and d) follow from the induction hypothesis: dim Y, dim \(Y′\leq N−1\). c) follows from a) and b), and e) follows from c) and d). This completes the proof.\par
\par\medskip
9. EXAMPLES AND APPLICATIONS\par
\par\medskip
9.1. Cohomology of groups and classifying spaces.\par
\par\medskip
Let G be a connected complex linear algebraic group. We propose to calculate the mixed Hodge structure on \(H^{\bullet }(G)\). We shall first calculate that on \(H^{\bullet }(B_{\bullet G})\), by reduction to the case of tori; we shall then calculate that on \(H^{\bullet }(G)\) via the Eilenberg-Moore spectral sequence (6.1.5). As a bonus, we obtain the curious Theorem (9.1.7).\par
\par\medskip
Theorem (9.1.1). — Let G be a linear algebraic group. One has \(H^{2n+1}(B_{\bullet G},\mathbb{Q})=0\), and \(H^{2n}(B_{\bullet G},\mathbb{Q})\otimes \mathbb{C}\) is purely of type (n,n) (III.0.5).\par
\par\medskip
Let \(G^{0}\) be the identity component of G and let \(T\subset G^{0}\) be a maximal torus. Since\par
\par\medskip
\DFormula{H^{\bullet }(B_{\bullet G},\mathbb{Q})=H^{\bullet }(B_{\bullet G^{0}},\mathbb{Q})^{G/G^{0}},}
\par\medskip
by the splitting principle (see (6.1.6)) one has\par
\par\medskip
\DFormula{H^{\bullet }(B_{\bullet G},\mathbb{Q})\,\hookrightarrow \,H^{\bullet }(B_{\bullet T},\mathbb{Q}).}
\par\medskip
This inclusion is a morphism of mixed Hodge structures: it suffices to treat the case of a torus, indeed, by the Künneth formula, the case \(G=G_{m}\). Here are two proofs of (9.1.1) in this case.\par
\par\medskip
First proof. — The cohomology of \(G_{m}\) \((=the\) projective line minus two points) is \(H^{0}=\mathbb{Z}\), \(H^{1}=\mathbb{Z}\). The only Hodge numbers \(h^{pq}\) of \(H^{1}(G_{m})\) that could be nonzero are \(h^{01}=h^{10}\) and \(h^{11}\) ((8.2.4)(i)(iv)). Since \(h^{01}+h^{10}+h^{11}=2h^{01}+h^{11}=1\), one has \(h^{01}=0\) and \(H^{1}(G_{m})\) is purely of type (1,1) (for another proof, see (10.3.8)). The term \(E_{2}^{p,q}\) of the spectral sequence (6.1.5), zero for \(p+q\) odd, is therefore purely of type \(((p+q)/2,(p+q)/2)\). Its abutment has the same homogeneity.\par
\par\medskip
Second proof. — One knows that \(ℙ^{r}(\mathbb{C})\), equipped with \(\mathcal{O}(1)\), is an “approximation” to \(B_{G_{m}}\): the homomorphism \([\mathcal{O}(1)]:H^{i}(B_{\bullet G_{m}})\to H^{i}(ℙ^{r}(\mathbb{C}))\) is bijective for \(i\leq 2r\). The cohomology group \(H^{2i}(ℙ^{r}(\mathbb{C}))\) is generated by the cohomology class of an algebraic cycle (namely, a linear subspace), and is therefore of type (i,i). The following proposition allows one to deduce that \(H^{2i}(B_{G_{m}})\) is of type (i,i).\par
\DAsset{P0041-A01}{Hodge bounds, end of (8.3.10), section 9, and Theorem (9.1.1)}
\DAsset{P0041-A02}{two proofs of Theorem (9.1.1) and projective approximation}

\DArticleAnchor{0042}{42}{43}{709}{46}{P07}

Proposition (9.1.2). — Let G be an algebraic group (over \(\mathbb{C})\), and let P be an algebraic principal homogeneous space under G over a scheme X. The morphism\par
\par\medskip
\DFormula{[P]:H^{\bullet }(B_{G})=H^{\bullet }(B_{\bullet G})\to H^{\bullet }(X)}
\par\medskip
is a morphism of mixed Hodge structures.\par
\par\medskip
This follows from the definition (6.1.3.1) of [P] and the functoriality (8.3.4).\par
\par\medskip
Corollary (9.1.3). — The characteristic classes in \(H^{2n}(S)\) of an algebraic principal homogeneous space P over S, whose structure group is a linear group G, are purely of type (n,n): they come from morphisms \(\mathbb{Q}(−n)\to H^{2n}(S,\mathbb{Q})\).\par
\par\medskip
In the particular case where S is quasi-projective and smooth, it is known that the characteristic classes of P are even algebraic, that is, they lie in the image of the Chow ring \((\otimes \mathbb{Q})\).\par
\par\medskip
(9.1.4) Let G be a linear algebraic group over a finite field \(𝔽_{q}\), and let \(\overline{G}/\overline{𝔽}_{q}\) be the group obtained by extension of scalars to an algebraic closure of \(𝔽_{q}\). A proof parallel to (9.1.1) shows that, in \(\ell -adic\) cohomology, the eigenvalues of (geometric) Frobenius acting on \(H^{2n}(B_{\bullet \overline{G}},\mathbb{Q}_{\ell })\) are of the form \(q^{n}·\zeta \) \((\zeta \) a root of unity). This agrees with the philosophy of (I.3). Likewise, statement (9.1.5) below is inspired by the formulas giving the numbers of points of reductive groups over finite fields.\par
\par\medskip
Theorem (9.1.5). — Let G be a connected linear algebraic group. Let \(P^{\bullet }(G)\subset H^{\bullet }(G,\mathbb{Q})\) be the primitive part of the Hopf algebra \(H^{\bullet }(G,\mathbb{Q})\). It is a mixed Hodge substructure. One has \(P^{2i}(G)=0\), and \(P^{2i−1}(G)\) is purely of type (i,i). Finally, the isomorphism\par
\par\medskip
\DFormula{(9.1.5.1)\,\Lambda P^{\bullet }(G)\to H^{\bullet }(G,\mathbb{Q})}
\par\medskip
is an isomorphism of mixed Hodge structures.\par
\par\medskip
In the spectral sequence (6.1.5), \(E^{0,q}=0\) for \(q>0\). For \(n>0\), the edge homomorphism\par
\par\medskip
\DFormula{H^{2n}(B_{G},\mathbb{Q})\to E_{1}^{1,2n}=H^{2n−1}(G,\mathbb{Q})}
\par\medskip
is a morphism of mixed Hodge structures (8.3.5); its image \(P^{2n−1}(G)\) is therefore a mixed Hodge substructure, and it is purely of type (n,n) by (9.1.1). The isomorphism (9.1.5.1) is given by the cup product, and one concludes by (8.2.11).\par
\par\medskip
Corollary (9.1.6). — Let G be a linear algebraic group. For \(n>0\), one has\par
\par\medskip
\DFormula{W_{n}(H^{n}(G,\mathbb{Q}))=0.}
\par\medskip
It suffices to reduce to the case where G is connected, and then apply (9.1.5).\par
\par\medskip
Theorem (9.1.7). — Let X be a proper smooth algebraic space on which a connected linear algebraic group G acts. Let \(\rho :G\times X\to X\). The diagram\par
\DAsset{P0042-A01}{Proposition (9.1.2), Corollary (9.1.3), Frobenius, and beginning of (9.1.5)}
\DAsset{P0042-A02}{Theorem (9.1.5), Corollary (9.1.6), and beginning of (9.1.7)}

\DArticleAnchor{0043}{43}{44}{710}{47}{P08}

\DAsset{P0043-A01}{continuation of the diagram in Theorem (9.1.7): \(H^{\bullet }(X,\mathbb{Q})\) \(\xrightarrow{\rho ^{\ast }}\) \(H^{\bullet }(X,\mathbb{Q})\otimes H^{\bullet }(G,\mathbb{Q})\), with lower row \(H^{\bullet }(X,\mathbb{Q})\otimes _{\mathbb{Q}}\mathbb{Q}\) \(=\) \(H^{\bullet }(X,\mathbb{Q})\otimes H^{0}(G,\mathbb{Q})\)}
\par\medskip
is commutative.\par
\par\medskip
Since the composite map\par
\par\medskip
\DFormula{H^{\bullet }(X,\mathbb{Q})\,\xrightarrow{\rho ^{\ast }}\,H^{\bullet }(X,\mathbb{Q})\otimes H^{\bullet }(G,\mathbb{Q})\,\to \,H^{\bullet }(X,\mathbb{Q})\otimes H^{0}(G,\mathbb{Q})=H^{\bullet }(X,\mathbb{Q})}
\par\medskip
is the identity (the composite \(X=\lbrace e\rbrace \times X\to G\times X\to X\) is the identity), it is enough to prove that, for \(i>0\), the i-th Künneth component of \(\rho ^{\ast }\),\par
\par\medskip
\DFormula{\rho _{i}^{\ast }\,:\,H^{n}(X,\mathbb{Q})\,\to \,H^{n−i}(X,\mathbb{Q})\otimes H^{i}(G,\mathbb{Q}),}
\par\medskip
is zero. Since \(H^{n}(X,\mathbb{Q})\) is pure of weight n and \(H^{n−i}(X,\mathbb{Q})\) is pure of weight \(n−i\), one has\par
\par\medskip
\DFormula{\operatorname{Im}(H^{n}(X,\mathbb{Q}))\,\subset \,W_{n}(H^{n−i}(X,\mathbb{Q})\otimes H^{i}(G,\mathbb{Q}))\,=\,H^{n−i}(X,\mathbb{Q})\otimes W_{i}\,H^{i}(G,\mathbb{Q}).}
\par\medskip
By (9.1.6), this image is therefore zero.\par
\par\medskip
Variant (9.1.8). — Let X be a separated algebraic space on which a connected linear algebraic group G acts.\par
\par\medskip
(i) If X is smooth and \(\overline{X}\) is a compactification of X, the diagram\par
\par\medskip
\DAsset{P0043-A02}{diagram (9.1.8)(i): \(H^{\bullet }(\overline{X},\mathbb{Q})\) \(\to \) \(H^{\bullet }(X,\mathbb{Q})\) \(\xrightarrow{\rho ^{\ast }}\) \(H^{\bullet }(X,\mathbb{Q})\otimes H^{\bullet }(G,\mathbb{Q})\), with lower row \(H^{\bullet }(X,\mathbb{Q})=H^{\bullet }(X,\mathbb{Q})\otimes H^{0}(G,\mathbb{Q})\)}
\par\medskip
is commutative.\par
\par\medskip
(ii) If X is proper and \(\widetilde{X}\) is a resolution of the singularities of X, the diagram\par
\par\medskip
\DAssetReference{P0043-A02}{diagram (9.1.8)(ii): \(H^{\bullet }(X,\mathbb{Q})\) \(\xrightarrow{\rho ^{\ast }}\) \(H^{\bullet }(X,\mathbb{Q})\otimes H^{\bullet }(G,\mathbb{Q})\) \(\to \) \(H^{\bullet }(\widetilde{X},\mathbb{Q})\otimes H^{\bullet }(G,\mathbb{Q})\), with lower row \(H^{\bullet }(\widetilde{X},\mathbb{Q})=H^{\bullet }(\widetilde{X},\mathbb{Q})\otimes H^{0}(G,\mathbb{Q})\)}
\par\medskip
is commutative.\par

\DArticleAnchor{0044}{44}{45}{711}{48}{P08}

It is enough to prove that, for \(i>0\), the Künneth components\par
\par\medskip
\DFormula{H^{n}(\overline{X},\mathbb{Q})\,\to \,H^{n}(X,\mathbb{Q})\,\to \,H^{n−i}(X,\mathbb{Q})\otimes H^{i}(G,\mathbb{Q})}
\par\medskip
\DFormula{H^{n}(X,\mathbb{Q})\,\to \,H^{n−i}(X,\mathbb{Q})\otimes H^{i}(G,\mathbb{Q})\,\to \,H^{n−i}(\widetilde{X},\mathbb{Q})\otimes H^{i}(G,\mathbb{Q})}
\par\medskip
are respectively zero. In case (i) (respectively (ii)), one has, by (8.2.4),\par
\par\medskip
\DFormula{W_{n}\,H^{n}(\overline{X},\mathbb{Q})=H^{n}(\overline{X},\mathbb{Q})\,\,\,\,(respectively\,W_{n}\,H^{n}(X,\mathbb{Q})=H^{n}(X,\mathbb{Q})),}
\par\medskip
whereas, by (8.2.4) and (9.1.6),\par
\par\medskip
\DFormula{W_{n}(H^{n−i}(X,\mathbb{Q})\otimes H^{i}(G,\mathbb{Q}))=0\,\,\,\,(respectively\,W_{n}(H^{n−i}(\widetilde{X},\mathbb{Q})\otimes H^{i}(G,\mathbb{Q}))=0),}
\par\medskip
which proves the assertion.\par
\par\medskip
9.2. Hodge theory of smooth hypersurfaces, following Griffiths.\par
\par\medskip
This subsection depends only on §§ 1 to 3. We explain in it Theorem (8.6) of [7].\par
\par\medskip
(9.2.1) Let X be a proper smooth algebraic variety, and let Y be a smooth divisor in X. Let j be the inclusion of \(U=X−Y\) into X. One has \(R^{i}\) \(j_{\ast }\mathbb{Z}=0\) for \(i\neq 0,1\), and the Leray spectral sequence for j is identified with the long exact cohomology sequence\par
\par\medskip
\DFormula{(9.2.1.1)\,\,…\,\to \,H_{Y}^{n}(X)\,\to \,H^{n}(X)\,\to \,H^{n}(U)\,\xrightarrow{\partial }\,…}
\par\medskip
\DFormula{H_{Y}^{n}(X)=H^{n−2}(Y,\mathcal{H}_{Y}^{2}).}
\par\medskip
For a reason made clear in (3.1.9), we shall identify \(\mathcal{H}_{Y}^{2}\) with the constant sheaf \(\mathbb{Z}(−1)_{Z}\). Since the Leray spectral sequence of j converges, by definition, to the mixed Hodge structure on \(H^{n}(U)\), it follows that the exact sequence\par
\par\medskip
\DFormula{(9.2.1.2)\,\,…\,\to \,H^{n−2}(Y,\mathbb{Z})(−1)\,\to \,H^{n}(X,\mathbb{Z})\,\to \,H^{n}(U,\mathbb{Z})\,\xrightarrow{\partial }\,…}
\par\medskip
is an exact sequence of mixed Hodge structures.\par
\par\medskip
The Leray spectral sequence of j is also the spectral sequence of the filtered complex \((\Omega _{X}^{\bullet }(\operatorname{log}\) Y),W); the preceding exact sequence is therefore also identified with the one deduced from the short exact sequence\par
\par\medskip
\DFormula{(9.2.1.3)\,\,0\,\to \,\Omega _{X}^{\bullet }\,\to \,\Omega _{X}^{\bullet }(\operatorname{log}\,Y)\,\xrightarrow{Res}\,\Omega _{Y}^{\bullet −1}\,\to \,0.}
\par\medskip
For this reason, the homomorphism \(\partial \) in (9.2.1.2) is also called the residue operator.\par
\par\medskip
(9.2.2) The residue operator\par
\par\medskip
\DFormula{(9.2.2.1)\,\,Res\,:\,H^{n+1}(U,\mathbb{Z})\,\to \,H^{n}(Y,\mathbb{Z})(−1)}
\par\medskip
is strictly compatible with the Hodge filtration. The Hodge filtration on \(H^{n}(U,\mathbb{C})\) therefore determines that on \(\operatorname{Ker}(H^{n}(Y,\mathbb{C})\) \(\to \) \(H^{n+2}(X,\mathbb{C}))\).\par
\DAsset{P0044-A01}{end of Variant (9.1.8), section 9.2, and exact sequence (9.2.1.1)}
\DAsset{P0044-A02}{residue exact sequences (9.2.1.2)-(9.2.1.3) and operator (9.2.2.1)}

\DArticleAnchor{0045}{45}{46}{712}{49}{P08}

Assume Y ample. If X is pure of dimension \(N+1\), then\par
\par\medskip
\DFormula{H^{i}(X,\mathbb{Z})\,\xrightarrow{\sim }\,H^{i}(Y,\mathbb{Z})\,for\,i<N\,\,(Lefschetz).}
\par\medskip
By Poincaré duality, for \(i\neq N\), \(H^{i}(Y,\mathbb{C})\) can therefore be calculated from \(H^{\bullet }(X,\mathbb{C})\). Put\par
\par\medskip
\DFormula{H_{ev}^{N}(Y,\mathbb{C})=\operatorname{Ker}(H^{N}(Y,\mathbb{C})\,\to \,H^{N+2}(X,\mathbb{C}))}
\par\medskip
\(=\) the orthogonal complement of \(H^{N}(X,\mathbb{C})\) \(\hookrightarrow \) \(H^{N}(Y,\mathbb{C})\)\par
\par\medskip
(the vanishing part of the cohomology). One has\par
\par\medskip
\DFormula{H^{N}(Y,\mathbb{C})=H_{ev}^{N}(Y,\mathbb{C})\oplus H^{N}(X,\mathbb{C}),}
\par\medskip
and one therefore obtains much information about the Hodge structure on \(H^{\bullet }(Y,\mathbb{C})\) when one knows that on \(H^{\bullet }(X,\mathbb{C})\) and the Hodge filtration on \(H^{N+1}(U,\mathbb{C})\).\par
\par\medskip
(9.2.3) By definition, the Hodge filtration on \(H^{\bullet }(U,\mathbb{C})\) is the abutment of the spectral sequence, degenerate at \(E_{1}\),\par
\par\medskip
\DFormula{E_{1}^{pq}=H^{q}(X,\Omega _{X}^{p}(\operatorname{log}\,Y))\,\Rightarrow \,H^{p+q}(U,\mathbb{C}).}
\par\medskip
By (3.1.11), this spectral sequence coincides with the spectral sequence of the complex \(j_{\ast }\Omega _{U}^{\bullet }\), equipped with the filtration P by order of pole.\par
\par\medskip
In particular, one has\par
\par\medskip
\DFormula{F^{p}(H^{\bullet }(U,\mathbb{C}))=H^{\bullet }(X,P^{p}(j_{\ast }\Omega _{U}^{\bullet })).}
\par\medskip
Recall that \(P^{p}(j_{\ast }\Omega _{U}^{\bullet })\) is the following complex, whose first nonzero component is in degree p:\par
\par\medskip
\DFormula{\Omega _{X}^{p}(Y)\,\to \,\Omega _{X}^{p+1}(2Y)\,\to \,…\,\to \,\Omega _{X}^{i}((i−p+1)Y)\,\to \,…}
\par\medskip
If \(H^{i}(X,\Omega _{X}^{p}(nY))=0\) for \(i>0\), \(n>0\), and \(p\geq 0\), one simply has\par
\par\medskip
\DFormula{H^{\bullet }(X,P^{p}(j_{\ast }\Omega _{U}^{\bullet }))=H^{\bullet }(\Gamma (X,P^{p}(j_{\ast }\Omega _{U}^{\bullet }))),}
\par\medskip
whence the following result:\par
\par\medskip
Proposition (9.2.4). — Assume that \(H^{i}(X,\Omega _{X}^{p}(nY))=0\) for \(i>0\), \(n>0\), and \(p\geq 0\) (this is the case if Y is sufficiently ample). Then a cohomology class \(c\in H^{d}(U,\mathbb{C})\) has Hodge filtration \(\geq p\) if and only if it can be represented by a closed d-form \(\alpha \) on U having a pole of order \(\leq d−p+1\) along Y. If \(\alpha \) is a closed form on U having a pole of order k along Y, and if the class of \(\alpha \) in \(H^{d}(U,\mathbb{C})\) is zero, then \(\alpha =d\beta \), where \(\beta \) has a pole of order \(\leq k−1\) along Y; for \(k\leq 1\), one even has \(\alpha =0\).\par
\par\medskip
(9.2.5) It is well known (Bott's theorem) that the hypotheses of (9.2.4) are satisfied when Y is a hypersurface in projective space \(ℙ^{n}(\mathbb{C})\) \((n\geq 1)\), that is,\par
\par\medskip
\DFormula{H^{i}(ℙ^{n}(\mathbb{C}),\Omega ^{j}(m))=0\,\,\,\,for\,i>0,\,m>0}
\par\medskip
(a proof may be found, for example, in SGA 7 XI). Moreover, for \(n>1\), the residue map identifies \(H^{n}(U,\mathbb{C})\) with the primitive part of the cohomology \(H^{n−1}(Y,\mathbb{C})\). Hence ([7], (8.6)):\par
\DAsset{P0045-A01}{vanishing cohomology and Hodge filtration from the logarithmic spectral sequence}
\DAsset{P0045-A02}{pole-order complex, Proposition (9.2.4), and Bott theorem (9.2.5)}

\DArticleAnchor{0046}{46}{47}{713}{50}{P08}

Proposition (9.2.6). — Let Y be a smooth hypersurface in \(ℙ^{n}(\mathbb{C})\) \((n>1)\), and let \(U=ℙ^{n}(\mathbb{C})−Y\).\par
\par\medskip
(i) For a primitive cohomology class \(c\in H^{n−1}(Y,\mathbb{C})\) to have Hodge filtration \(\leq p\), it is necessary and sufficient that it be the residue of a regular n-form on U having a pole of order \(\leq n−p\) along Y.\par
\par\medskip
(ii) Let \(\alpha \) be a regular n-form on U having a pole of order \(k\geq 1\) along Y. For \(Res(\alpha )\in H^{n−1}(Y,\mathbb{C})\) to be zero, it is necessary and sufficient that \(\alpha =d\beta \) for an \((n−1)-form\) on U having a pole of order \(\leq k−1\) along Y.\par
\par\medskip
9.3. Construction of complexes of first-order differential operators.\par
\par\medskip
This subsection depends only on §§ 5 and 6. We prove the following theorem in it.\par
\par\medskip
Proposition (9.3.1). — Let X be a quasi-projective scheme (over \(\mathbb{C})\). There exists on X a complex of sheaves K having the following properties.\par
\par\medskip
a) The \(K^{n}\) are coherent algebraic sheaves, and \(K^{n}=0\) for \(n<0\).\par
\par\medskip
b) The differentials \(d:K^{n}\to K^{n+1}\) are first-order algebraic differential operators.\par
\par\medskip
c) The complex \(K^{\mathrm{an}}\) of coherent analytic sheaves on \(X^{\mathrm{an}}\) is a resolution of the constant sheaf \(\mathbb{C}\). More precisely, \(\mathcal{H}^{i}(K)=0\) for \(i>0\), and there exists\par
\par\medskip
\DFormula{f:\Omega _{X}^{\bullet }\to K}
\par\medskip
with \(f^{n}\) algebraically linear, such that the composite map \(\mathbb{C}\to \mathcal{O}_{X}\) \(\xrightarrow{f^{0}}\) \(K^{0}\) identifies \(\mathbb{C}\) with \(\mathcal{H}^{0}(K)\).\par
\par\medskip
d) The natural map\par
\par\medskip
\DFormula{H^{i}(X,K)\,\to \,H^{i}(X^{\mathrm{an}},K^{\mathrm{an}})\,\xleftarrow[(c)]{\sim }\,H^{i}(X^{\mathrm{an}},\mathbb{C})}
\par\medskip
is an isomorphism.\par
\par\medskip
As a corollary, one recovers a special case of a theorem of Bloom and Herrera [2].\par
\par\medskip
Corollary (9.3.2). — The map deduced from \(\mathbb{C}\to \Omega _{X}^{\bullet \mathrm{an}}\),\par
\par\medskip
\DFormula{H^{\bullet }(X^{\mathrm{an}},\mathbb{C})\,\to \,H^{\bullet }(X^{\mathrm{an}},\Omega _{X}^{\bullet \mathrm{an}}),}
\par\medskip
identifies \(H^{\bullet }(X^{\mathrm{an}},\mathbb{C})\) with a direct factor of \(H^{\bullet }(X^{\mathrm{an}},\Omega _{X}^{\bullet \mathrm{an}})\).\par
\par\medskip
We shall first prove the corollary. The proposition will be obtained by repeating the proof at the level of complexes. Let \(\varepsilon :X_{\bullet }\to X\) be a simplicial scheme augmented to X. We assume that \(\varepsilon \) is a proper hypercover and that the \(X_{n}\) are smooth. There is a commutative diagram\par
\par\medskip
\DAsset{P0046-A02}{diagram (9.3.2.1): upper row \(H^{\bullet }(X^{\mathrm{an}},\mathbb{C})\) \(\xrightarrow[1]{\sim }\) \(H^{\bullet }(X_{\bullet }^{\mathrm{an}},\mathbb{C})\); lower row \(H^{\bullet }(X^{\mathrm{an}},\Omega _{X}^{\bullet })\) \(\to \) \(H^{\bullet }(X_{\bullet }^{\mathrm{an}},\Omega _{X_{\bullet }}^{\bullet })\); both vertical arrows point downward, with the right vertical arrow labelled \(\sim \) and 2}
\DAsset{P0046-A01}{Proposition (9.2.6), section 9.3, and Proposition (9.3.1)}

\DArticleAnchor{0047}{47}{48}{714}{51}{P08}

Arrow 1 is bijective by (5.3.5)(II)(IV). Arrow 2 is bijective because \(\Omega _{X_{\bullet }}^{\bullet }\) is a resolution of \(\mathbb{C}\) (the Poincaré lemma on the \(X_{n})\). Diagram (9.3.2.1) supplies a retraction, compatible with the cup product, of \(H^{\bullet }(X^{\mathrm{an}},\Omega _{X}^{\bullet })\) onto \(H^{\bullet }(X^{\mathrm{an}},\mathbb{C})\), and this proves (9.3.2).\par
\par\medskip
(9.3.3) The complex K that we wish to construct will be none other, in the derived category, than \(sR\varepsilon _{\bullet \ast }\Omega _{X_{\bullet }}^{\bullet }\). We shall treat only the case in which the augmented simplicial space \(X_{\bullet }\) is s-split, defined as in (6.2.5) by proper surjective maps with quasi-projective smooth sources\par
\par\medskip
\DFormula{f′_{k}:N_{k}\to (\operatorname{cosq}\,\operatorname{sq}_{k−1}(X_{\bullet }))_{k}\,\,\,\,(k\geq 0).}
\par\medskip
We shall construct a “compatible system” of affine open coverings of the \(X_{n}\). For this construction it is useful to regard an open covering of a space Y as a diagram\par
\par\medskip
\DFormula{Y\,\xleftarrow{u}\,U\,\xrightarrow{\varphi }\,I,}
\par\medskip
with I discrete, u surjective, and \(u_{i}:U_{i}=\varphi ^{−1}(i)\hookrightarrow Y\) an open immersion. Proceeding as in (6.2.5), one constructs a diagram\par
\par\medskip
\DFormula{(9.3.3.1)\,\,X_{\bullet }\,\xleftarrow{u}\,U_{\bullet }\,\xrightarrow{\varphi }\,I_{\bullet }}
\par\medskip
of the following type.\par
\par\medskip
a) \(I_{\bullet }\) is a discrete simplicial set. The \(I_{n}\) are finite.\par
\par\medskip
\DFormula{b)\,For\,every\,k\geq 0,\,(9.3.3.1)\,induces\,a\,diagram}
\par\medskip
\DAsset{P0047-A01}{diagram b): (cosq \(\operatorname{sq}_{k−1}(X_{\bullet }))_{k}\) \(\leftarrow \) (cosq \(\operatorname{sq}_{k−1}(U_{\bullet }))_{k}\) \(\xrightarrow{\varphi }\) (cosq \(\operatorname{sq}_{k−1}(I_{\bullet }))_{k}\), above \(N_{k}\) \(\leftarrow \) \(N_{k}(U_{\bullet })\) \(\xrightarrow{\varphi }\) \(N_{k}(I_{\bullet })\), with vertical arrows \(f′_{k}\) and \(a′_{k}\)}
\par\medskip
c) In diagram b), for \(i\in (\operatorname{cosq}\) \(\operatorname{sq}_{k−1}(I_{\bullet }))_{k}\), the \(\varphi ^{−1}(j)\), for \(a′_{k}(j)=i\), form an affine open covering of\par
\par\medskip
\DFormula{N_{k}\,\times _{(\operatorname{cosq}\,\operatorname{sq}_{k−1}(X_{\bullet }))_{k}}\,\varphi ^{−1}(i).}
\par\medskip
d) (follows inductively from c)). The first (and hence the second) row of diagram b) defines an open covering. The \(X_{k}\leftarrow U_{k}\to I_{k}\) therefore also form open coverings.\par
\par\medskip
e) (follows from a suitable choice of the affine coverings in c)). \(U_{n_{i}}=\varphi _{n}^{−1}(i)\) is the complement of a closed subscheme \(D_{n,i}\) of \(X_{n}\), which is a union of some connected components of \(X_{n}\) and a divisor. For every coherent sheaf \(\mathcal{F}\) on \(X_{n}\) and integers \((a_{i})_{i\in I_{n}}\) \((a_{i}\geq 0)\), put\par
\par\medskip
\DAsset{P0047-A02}{piecewise definition: \(\mathcal{F}(\sum _{i}\) \(a_{i}D_{n,i})\) equals 0 on the components of \(X_{n}\) contained in a \(D_{n,i}\) with \(a_{i}>0\), and \(\mathcal{F}\otimes \mathcal{O}(\sum _{a_{i}>0}a_{i}D_{n,i})\) on the other components}

\DArticleAnchor{0048}{48}{49}{715}{52}{P08}

In the inductive construction of (9.3.3.1), the open coverings in c) are chosen arbitrarily. This allows us to assume:\par
\par\medskip
f) Whatever the \(a_{i}>0\) \((i\in I_{k})\),\par
\par\medskip
\DFormula{R^{\ell }\,\varepsilon _{k\ast }(\Omega _{X_{k}}^{m}(\sum _{i}\,a_{i}D_{k,i}))=0\,\,\,\,for\,\ell >0.}
\par\medskip
g) For every \(\alpha :\Delta _{m}\to \Delta _{n}\) and \(i\in I_{n}\), one has\par
\par\medskip
\DFormula{X_{\bullet }(\alpha )^{−1}(D_{m,\alpha (i)})\,\subset \,D_{n,i}\,\,\,\,(as\,schemes).}
\par\medskip
For each k, the complex of sheaves on X\par
\par\medskip
\DFormula{R\varepsilon _{k\ast }(\Omega _{X_{k}}^{\bullet })}
\par\medskip
can be calculated by a Čech procedure. If, for \(P\subset I_{k}\), \(j_{P}\) is the inclusion of \(U_{P}=⋂_{i\in P}U_{k_{i}}\) into \(X_{k}\), it is the simple complex associated with the double complex whose components are\par
\par\medskip
\DFormula{\widetilde{K}^{k,p,q}=\sum _{\#P=p+1}(\varepsilon _{k}\,j_{P})_{\ast }\,j_{P}^{\ast }\Omega _{X_{k}}^{q}.}
\par\medskip
This double complex contains as a double subcomplex the complex whose components are the coherent sheaves\par
\par\medskip
\DFormula{K^{k,p,q}=\sum _{\#P=p+1}\,\varepsilon _{k\ast }\Omega _{X_{k}}^{q}((q+1)\sum _{i\in P}D_{k,i}).}
\par\medskip
For fixed k and q, it follows from f) that the inclusion of \(K^{k,\bullet ,q}\) into \(\widetilde{K}^{k,\bullet ,q}\) is a quasi-isomorphism. For fixed k, the inclusion of \(K^{k,\bullet ,\bullet }\) into \(\widetilde{K}^{k,\bullet ,\bullet }\) is likewise a quasi-isomorphism. Finally, as k varies, the \(K^{k,\bullet ,\bullet }\) form, by g), a simplicial double complex. The associated simple complex K coincides with \(R\varepsilon _{\bullet \ast }\Omega _{X_{\bullet }}^{\bullet }\) in the derived category.\par
\par\medskip
Let us prove that K satisfies (9.3.1). Properties a) and b) of K are clear; we define \(f:\Omega _{X}^{\bullet }\to K\) by the evident maps from \(\Omega _{X}^{q}\) into \(K^{0,0,q}\).\par
\par\medskip
In the analytic category, the complex \(K^{\mathrm{an}}\) of coherent analytic sheaves is given by the same formulas as K. There is a commutative diagram\par
\par\medskip
\DAsset{P0048-A02}{diagram: a headless horizontal line labelled \(\sim \) joins \(H^{i}(X,K)\) and \(H^{i}(X_{\bullet },\Omega _{X_{\bullet }}^{\bullet })\); a headless horizontal line labelled \(\sim \) joins \(H^{i}(X^{\mathrm{an}},K^{\mathrm{an}})\) and \(H^{i}(X_{\bullet }^{\mathrm{an}},\Omega _{X_{\bullet }^{\mathrm{an}}}^{\bullet })\); two vertical arrows point downward from each upper term to its corresponding lower term, with the right vertical arrow labelled 1}
\par\medskip
and 1 is the abutment of a morphism of spectral sequences\par
\par\medskip
\DAssetReference{P0048-A02}{morphism of spectral sequences: upper row \(H^{a}(X_{b},\Omega _{X_{\bullet }}^{\bullet })\) \(\Rightarrow \) \(H^{a+b}(X_{\bullet },\Omega _{X_{\bullet }}^{\bullet })\); lower row \(H^{a}(X_{b}^{\mathrm{an}},\Omega _{X_{\bullet }^{\mathrm{an}}}^{\bullet })\) \(\Rightarrow \) \(H^{a+b}(X_{\bullet }^{\mathrm{an}},\Omega _{X_{\bullet }^{\mathrm{an}}}^{\bullet })\); two vertical arrows point downward from each upper term to its corresponding lower term}
\par\medskip
This morphism is an isomorphism by [9], which proves d).\par
\DAsset{P0048-A01}{conditions f)-g), Čech calculation, and double complexes K and \(\widetilde{K}\)}

\DArticleAnchor{0049}{49}{50}{716}{53}{P09}

To prove c), consider the commutative diagram\par
\par\medskip
\DAsset{P0049-A01}{diagram: a headless horizontal link labelled \(\approx \) joins \(\mathbb{C}\) and \(R\varepsilon _{\bullet \ast }\mathbb{C}\); the lower row is \(\Omega _{X_{\bullet }}^{\bullet \mathrm{an}}\) \(\xrightarrow{f}\) \(K^{\mathrm{an}}\), followed by a headless horizontal link labelled \(\approx \) joining \(K^{\mathrm{an}}\) and \(R\varepsilon _{\bullet \ast }\Omega _{X_{\bullet }}^{\bullet \mathrm{an}}\); the left vertical arrow points down from \(\mathbb{C}\) to \(\Omega _{X_{\bullet }}^{\bullet \mathrm{an}}\) and is labelled \(\text{\textcircled{1}}\); the right vertical arrow points down from \(R\varepsilon _{\bullet \ast }\mathbb{C}\) to \(R\varepsilon _{\bullet \ast }\Omega _{X_{\bullet }}^{\bullet \mathrm{an}}\) and is labelled \(\approx \) and \(\text{\textcircled{2}}\)}
\par\medskip
That 1 is a quasi-isomorphism follows from (5.3.5).\par
That 2 is one follows from the Poincaré lemma on the \(X_{n}\). Assertion c) follows.\par
\par\medskip
(9.3.4) The same construction furnishes resolutions satisfying a), b) for sheaves \(\mathcal{F}\) on \(X^{\mathrm{an}}\) other than the constant sheaf \(\mathbb{C}\) (one can doubtless take for \(\mathcal{F}\) any “algebraically constructible” sheaf). Unfortunately, I cannot prove that the resolutions obtained are essentially unique (in a suitable derived category). In the absence of such uniqueness, (9.3.1) is of little interest.\par
\par\medskip
\DFormula{10.\,HODGE\,THEORY\,IN\,LEVEL\,\leq \,1}
\par\medskip
10.1. 1-motives\par
\par\medskip
Definition (10.1.1). — (i) A mixed Hodge structure H is torsion-free if \(H_{\mathbb{Z}}\) is torsion-free.\par
\par\medskip
(ii) If H is torsion-free, W also denotes the filtration of \(H_{\mathbb{Z}}\subset H_{\mathbb{Q}}\) induced by W. We denote by \(\operatorname{Gr}_{n}^{W}(H)\) the torsion-free \(\mathbb{Z}-module\) \(\operatorname{Gr}_{n}^{W}(H_{\mathbb{Z}})\), endowed with its Hodge structure of weight n.\par
\par\medskip
Definition (10.1.2). — A 1-motive over an algebraically closed field k consists of:\par
\par\medskip
a) a finite-type free \(\mathbb{Z}-module\) X, an abelian variety A, and a torus T;\par
b) an extension G of A by T;\par
\DFormula{c)\,a\,homomorphism\,u:X\to G(k).}
\par\medskip
It is often useful to regard a 1-motive M as a complex of group schemes concentrated in degrees 0 and 1,\par
\par\medskip
\DFormula{M=[X\,\xrightarrow{u}\,G].}
\par\medskip
Construction (10.1.3). — We shall construct an equivalence \(M↦T(M)\) from the category of 1-motives over \(\mathbb{C}\) to the category of torsion-free mixed Hodge structures H of type\par
\par\medskip
\DFormula{\lbrace (0,0),\,(0,−1),\,(−1,0),\,(−1,−1)\rbrace }
\DAsset{P0049-A02}{Section 10, definitions of 1-motives, and Construction (10.1.3)}

\DArticleAnchor{0050}{50}{51}{717}{54}{P09}

such that \(\operatorname{Gr}_{−1}^{W}(H)\) is polarizable. Moreover, we shall construct isomorphisms natural in \(M=(X,A,T,G,u)\):\par
\par\medskip
\DFormula{H_{1}(T,\mathbb{Z})\,\simeq \,\operatorname{Gr}_{−2}^{W}\,T(M)_{\mathbb{Z}}}
\DFormula{H_{1}(A,\mathbb{Z})\,\simeq \,\operatorname{Gr}_{−1}^{W}\,T(M)\,\,\,\,\,(isomorphism\,of\,Hodge\,structures)}
\DFormula{X\,\simeq \,\operatorname{Gr}_{0}^{W}\,T(M)_{\mathbb{Z}}.}
\par\medskip
Let M be a 1-motive. The exponential map \(\operatorname{exp}:\operatorname{Lie}(G)\to G\) has kernel \(H_{1}(G)\). We define \(T_{\mathbb{Z}}(M)\) (which we shall rather denote \(T_{Z}(M))\) to be the fiber product of Lie(G) and X over G:\par
\par\medskip
\DAsset{P0050-A01}{diagram (10.1.3.1): \(0\to H_{1}(G)\to \operatorname{Lie}(G)\xrightarrow{\operatorname{exp}}G\to 0\) above \(0\to H_{1}(G)\to T_{Z}(M)\xrightarrow{\beta }X\to 0\); the vertical arrows are the identity, \(\alpha \), and u}
\par\medskip
Set\par
\par\medskip
\DFormula{W_{−1}(T_{Z}(M))=\operatorname{Ker}(\beta )=H_{1}(G),}
\DFormula{W_{−2}(T_{Z}(M))=H_{1}(T)=\operatorname{Ker}(H_{1}(G)\to H_{1}(A)).}
\par\medskip
This defines the weight filtration. The morphism \(\alpha \) extends to\par
\par\medskip
\DFormula{\alpha _{\mathbb{C}}:T_{Z}(M)\otimes \mathbb{C}\to \operatorname{Lie}(G).}
\par\medskip
Set \(F^{0}(T_{Z}(M)\otimes \mathbb{C})=\operatorname{Ker}(\alpha _{\mathbb{C}})\); this defines the Hodge filtration of\par
\par\medskip
\DFormula{T_{\mathbb{C}}(M)=T_{Z}(M)\otimes \mathbb{C}.}
\par\medskip
Lemma (10.1.3.2). — The triple \(T(M)=(T_{Z}(M),W,F)\) is a torsion-free mixed Hodge structure of type \(\lbrace (0,0),(0,−1),(−1,0),(−1,−1)\rbrace \), and \(\operatorname{Gr}_{−1}^{W}(T(M))\) is polarizable.\par
\par\medskip
\DAsset{P0050-A02}{large proof diagram: exact rows \(0\to H_{1}(T,\mathbb{Z})\otimes \mathbb{C}\to H_{1}(G,\mathbb{Z})\otimes \mathbb{C}\to H_{1}(A,\mathbb{Z})\otimes \mathbb{C}\to 0\) and \(0\to \operatorname{Lie}(T)\to \operatorname{Lie}(G)\to \operatorname{Lie}(A)\to 0\); above them, \(W_{−1}(T_{\mathbb{C}}(M))\cap F^{0}\) \(\xrightarrow[\text{\textcircled{2}}]{\sim }\) \(\operatorname{Ker}(\alpha _{A,\mathbb{C}})\); exact columns read downward, \(0\to W_{−1}(T_{\mathbb{C}}(M))\cap F^{0}\to H_{1}(G,\mathbb{Z})\otimes \mathbb{C}\xrightarrow{\alpha _{G,\mathbb{C}}}\operatorname{Lie}(G)\to 0\) and \(0\to \operatorname{Ker}(\alpha _{A,\mathbb{C}})\to H_{1}(A,\mathbb{Z})\otimes \mathbb{C}\xrightarrow{\alpha _{A,\mathbb{C}}}\operatorname{Lie}(A)\to 0\); the left vertical arrow points down from \(H_{1}(T,\mathbb{Z})\otimes \mathbb{C}\) to Lie(T) and is labelled \(\approx \) and \(\text{\textcircled{1}}\)}

\DArticleAnchor{0051}{51}{52}{718}{55}{P09}

a) One has \(W_{−2}(T_{\mathbb{C}}(M))\cap F^{0}=0\), because 1 is injective; hence \((W_{−2}(T_{\mathbb{C}}(M)),F)\) is a Hodge structure of type \((−1,−1)\).\par
\par\medskip
b) Since 2 is surjective, F induces on \(\operatorname{Gr}_{−1}^{W}(T_{Z}(M))\otimes \mathbb{C}\simeq H_{1}(A,\mathbb{Z})\otimes \mathbb{C}\) the Hodge filtration of \(H_{1}(A,\mathbb{Z})\otimes \mathbb{C}\): it is a polarizable Hodge structure of type \((−1,0)+(0,−1)\) (cf. (4.4.3)).\par
\par\medskip
c)\par
\par\medskip
\DAsset{P0051-A01}{exact diagram: \(0\to W_{−1}(T_{\mathbb{C}}(M))\cap F^{0}\to F^{0}\to V\to 0\) above \(0\to H_{1}(G,\mathbb{Z})\otimes \mathbb{C}\to T_{Z}(M)\otimes \mathbb{C}\to X\otimes \mathbb{C}\to 0\); the first two columns descend to Lie(G), and \(V\to X\otimes \mathbb{C}\) is the arrow printed in the authority}
\par\medskip
\(F^{0}\) maps onto \(\operatorname{Gr}_{0}^{W}(T_{\mathbb{C}}(M))\), which is therefore of type (0,0).\par
This completes the construction of T(M). It is clear that T(M) is functorial in M.\par
\par\medskip
Let H be a torsion-free mixed Hodge structure of type\par
\par\medskip
\DFormula{\lbrace (0,0),\,(0,−1),\,(−1,0),\,(−1,−1)\rbrace .}
\par\medskip
Suppose that \(\operatorname{Gr}_{−1}^{W}(H)\) is polarizable; the complex torus\par
\par\medskip
\DFormula{A=H_{\mathbb{Z}}\mathbin{\backslash}\,\operatorname{Gr}_{−1}^{W}(H_{\mathbb{C}})/F^{0}\operatorname{Gr}_{−1}^{W}(H_{\mathbb{C}})}
\par\medskip
is then an abelian variety (cf. (4.4.3)). Let T be the torus whose character group is the dual of \(\operatorname{Gr}_{−2}^{W}(H_{\mathbb{Z}})\):\par
\par\medskip
\DFormula{H_{1}(T)=\operatorname{Gr}_{−2}^{W}(H_{\mathbb{Z}}).}
\par\medskip
The complex analytic group\par
\par\medskip
\DFormula{G=W_{−1}(H_{\mathbb{Z}})\mathbin{\backslash}\,W_{−1}(H_{\mathbb{C}})/(F^{0}\cap W_{−1}(H_{\mathbb{C}}))}
\par\medskip
is an extension of A by T.\par
\par\medskip
Lemma (10.1.3.3). — The functor \(E↦E^{\mathrm{an}}\) is an equivalence from the category of algebraic groups that are extensions of A by T to the category of analytic groups that are extensions of A by T.\par
\DAsset{P0051-A02}{inverse construction from H and Lemma (10.1.3.3)}

\DArticleAnchor{0052}{52}{53}{719}{56}{P09}

We reduce to the case \(T=G_{m}\). By G.A.G.A., \(E↦E^{\mathrm{an}}\) is then, more generally, an equivalence from the category of \(G_{m}-torsors\) over A to the category of \(G_{m}-torsors\) over \(A^{\mathrm{an}}\).\par
\par\medskip
We therefore have G, an extension of A by T. Let \(X=\operatorname{Gr}_{0}^{W}(H_{\mathbb{Z}})\). We define \(u:X\to G\) as follows (cf. (2.2.1)):\par
\par\medskip
\DAsset{P0052-A01}{diagram: \(0\to W_{−1}(H_{\mathbb{Z}})\to W_{−1}(H_{\mathbb{C}})/(W_{−1}(H_{\mathbb{C}})\cap F^{0})\to G\to 0\) above \(0\to W_{−1}(H_{\mathbb{Z}})\to H_{\mathbb{Z}}\to X\to 0\); the middle quotient descends to \(H_{\mathbb{C}}/F^{0}\) and then to \(H_{\mathbb{Z}}\), while X maps upward to G}
\par\medskip
It is clear that the preceding construction, which associates a torsion-free 1-motive with H, is inverse to T and is functorial.\par
\par\medskip
If \(G_{i}\) \((i=1,2)\) is an extension of an abelian variety \(A_{i}\) by a torus \(T_{i}\), then for every morphism \(u:G_{1}\to G_{2}\), \(u(T_{1})\subset T_{2}\) and \(\operatorname{Ker}(u:A_{1}\to A_{2})^{0}=\operatorname{Im}(\operatorname{Ker}(u:G_{1}\to G_{2})\to A_{1})^{0}\). This corresponds to the fact that a morphism of \(\mathbb{Q}-mixed\) Hodge structures is strictly compatible with the weight filtration.\par
\par\medskip
(10.1.4) Let H be a torsion-free mixed Hodge structure of type\par
\par\medskip
\DFormula{\lbrace (0,0),\,(0,−1),\,(−1,0),\,(−1,−1)\rbrace .}
\par\medskip
The filtration W of \(H_{\mathbb{Z}}\) then defines a filtration of H by mixed-Hodge substructures. In parallel, if \(M=(X,A,T,G,u)\) is a 1-motive, W denotes the following increasing filtration of M:\par
\par\medskip
\(W_{i}(M)=0\)   for \(i<−2\)   and   \(W_{i}(M)=M\)   for \(i\geq 0\);\par
\DFormula{W_{−1}(M)=G\,\,\,(i.e.\,(\lbrace 0\rbrace ,A,T,G,0));}
\DFormula{W_{−2}(M)=T\,\,\,(i.e.\,(\lbrace 0\rbrace ,0,T,T,0)).}
\par\medskip
In the evident sense, the successive \(\operatorname{Gr}_{i}^{W}(M)\), for \(i=0,−1,−2\), are X, A, and T.\par
\par\medskip
(10.1.5) Construction (10.1.3), which associates \(T_{Z}(M)\) with a 1-motive M over \(\mathbb{C}\), is transcendental. We shall show that\par
\par\medskip
\DFormula{\widehat{T}(M)=T_{Z}(M)\otimes \widehat{\mathbb{Z}}=\prod _{\ell }\,T_{Z}(M)\otimes \mathbb{Z}_{\ell }}
\par\medskip
can be defined purely algebraically.\par
\par\medskip
Let M be a 1-motive over an algebraically closed field k of characteristic 0. Identify M with a complex \([X\xrightarrow{u}G]\) (degrees 0 and 1). For every integer \(n>0\), let \([\mathbb{Z}\xrightarrow{n}\mathbb{Z}]\) be the complex in degrees \(−1\) and 0, and let \(T_{\mathbb{Z}/n\mathbb{Z}}(M)\) be \(H^{0}\) of the complex\par
\par\medskip
\DFormula{[X\xrightarrow{u}G]\otimes [\mathbb{Z}\xrightarrow{n}\mathbb{Z}]:}
\DAsset{P0052-A02}{filtration of 1-motives and algebraic definition of T-hat}

\DArticleAnchor{0053}{53}{54}{720}{57}{P09}

\DAsset{P0053-A01}{diagram of the complex: \(X\xrightarrow{u}G\) above \(X\xrightarrow{u}G\), with vertical arrows n and \(−n\)}
\par\medskip
One has\par
\par\medskip
\DFormula{T_{\mathbb{Z}/n\mathbb{Z}}(M)=\lbrace (x,g)\,\mid \,u(x)=ng\rbrace /\lbrace (nx,u(x))\,\mid \,x\in X\rbrace .}
\par\medskip
In the derived category,\par
\par\medskip
\DFormula{T_{\mathbb{Z}/n\mathbb{Z}}(M)=H^{0}(M\otimes ^{L}\mathbb{Z}/n\mathbb{Z}).}
\par\medskip
\DFormula{For\,n=md,\,define\,transition\,morphisms}
\par\medskip
\DFormula{\varphi _{m,n}:T_{\mathbb{Z}/n\mathbb{Z}}(M)\to T_{\mathbb{Z}/m\mathbb{Z}}(M)\,\,\,by\,\,\,\varphi _{m,n}((x,g))=(x,dg).}
\par\medskip
Set\par
\par\medskip
\DFormula{\widehat{T}(M)=\varprojlim _{n}\,T_{\mathbb{Z}/n\mathbb{Z}}(M).}
\par\medskip
The filtration W of M induces a filtration W on the \(T_{\mathbb{Z}/n\mathbb{Z}}(M)\) and on \(\widehat{T}(M)\). One has\par
\par\medskip
\DFormula{\operatorname{Gr}_{0}^{W}(\widehat{T}(M))=X\otimes \widehat{\mathbb{Z}},}
\DFormula{\operatorname{Gr}_{−1}^{W}(\widehat{T}(M))=\varprojlim _{n}\,A_{n}=\widehat{T}(A),}
\DFormula{\operatorname{Gr}_{−2}^{W}(\widehat{T}(M))=\varprojlim _{n}\,T_{n}=Y\otimes \widehat{\mathbb{Z}}(1),}
\par\medskip
where Y is the dual of the character group of \(T=W_{−2}(M)\).\par
\par\medskip
Construction (10.1.6). — Let \(M=[X\xrightarrow{u}G]\) be a 1-motive over \(\mathbb{C}\). One has\par
\par\medskip
\DFormula{(10.1.6.1)\,\widehat{T}(M)\simeq T_{Z}(M)\otimes \widehat{\mathbb{Z}}.}
\par\medskip
The natural morphism \([T_{Z}(M)\to \operatorname{Lie}(G)]\to [X\to G]\) is a quasi-isomorphism. The quasi-isomorphisms\par
\par\medskip
\DAsset{P0053-A02}{three bracketed complexes: \([X\to G]\) above itself with n and \(−n\); \([T_{Z}\to \operatorname{Lie}(G)]\) above itself; \([T_{Z}\to 0]\) above \([T_{Z}\to 0]\), connected by quasi-isomorphisms and bearing the labels u, n, \(−n\), and c}
\par\medskip
furnish isomorphisms\par
\par\medskip
\DFormula{(10.1.6.2)\,T_{\mathbb{Z}/n\mathbb{Z}}(M)\simeq T_{Z}/nT_{Z}.}
\par\medskip
Under this isomorphism, \(t\in W_{−1}(T_{Z})\) is sent to \(\operatorname{exp}(t/n)\in G_{n}\). Passing to the limit, one deduces (10.1.6.1) from (10.1.6.2).\par
\par\medskip
Grothendieck’s techniques make it possible to transpose to de Rham cohomology what we have just done in \(\ell -adic\) cohomology. Let M be a 1-motive over an algebraically closed field k.\par

\DArticleAnchor{0054}{54}{55}{721}{58}{P09}

Construction (10.1.7). — Let us construct a vector space \(T_{DR}(M)\) over k, equipped with filtrations W and F.\par
\par\medskip
Let \(M=(X,A,T,G,u)\), and regard M as a complex concentrated in degrees \(−1\) and 0. An extension of M by \(G_{a}\) is an extension of M by the complex reduced to \(G_{a}\) placed in degree 0. One has \(\operatorname{Hom}(G,G_{a})=\operatorname{Ext}^{1}(X,G_{a})=0\). Hence:\par
\par\medskip
a) extensions of M by \(G_{a}\) have no automorphisms;\par
\par\medskip
b) the sequence\par
\par\medskip
\DFormula{0\to \operatorname{Hom}(X,G_{a})\to \operatorname{Ext}^{1}(M,G_{a})\to \operatorname{Ext}^{1}(G,G_{a})\to 0}
\par\medskip
is exact.\par
\par\medskip
One has \(\operatorname{Hom}(T,G_{a})=\operatorname{Ext}^{1}(T,G_{a})=0\), so that\par
\par\medskip
\DFormula{c)\,\operatorname{Ext}^{1}(A,G_{a})\xrightarrow{\sim}\operatorname{Ext}^{1}(G,G_{a}).}
\par\medskip
The vector space \(\operatorname{Ext}^{1}(M,G_{a})\) is therefore finite-dimensional and, by a), there is consequently a universal extension of M by a vector group; it is denoted \(M^{\natural }:X\to G^{\natural }\):\par
\par\medskip
\DAsset{P0054-A01}{universal-extension diagram: \(0\to X=X\) above \(0\to \operatorname{Ext}^{1}(M,G_{a})^{\ast }\to G^{\natural }\to G\to 0\), with the vertical arrows printed in the authority}
\par\medskip
Set\par
\par\medskip
\DFormula{T_{DR}(M)=\operatorname{Lie}(G^{\natural }),}
\par\medskip
\DFormula{F^{0}T_{DR}(M)=\operatorname{Ker}(\operatorname{Lie}\,G^{\natural }\to \operatorname{Lie}\,G)\simeq (\operatorname{Ext}^{1}(M,G_{a}))^{\ast }.}
\par\medskip
\(T_{DR}\) is functorial in M (as is \(M^{\natural })\), and the filtration W of \(T_{DR}(M)\) is defined as coming from the filtration W of M.\par
\par\medskip
Construction (10.1.8). — Let M be a 1-motive over \(\mathbb{C}\). One has\par
\par\medskip
\DFormula{(T_{DR}(M),F,W)\simeq (T_{\mathbb{C}}(M),F,W).}
\par\medskip
The \(\operatorname{Ext}^{i}(X,G_{a})\), \(\operatorname{Ext}^{i}(A,G_{a})\), and \(\operatorname{Ext}^{i}(T,G_{a})\) \((i=0,1)\) are the same in the algebraic and analytic categories. Consider the map \(T_{\mathbb{C}}(M)=T_{Z}(M)\otimes \mathbb{C}\to \operatorname{Lie}(G)\to G\). The diagram\par
\par\medskip
\DAsset{P0054-A02}{terminal diagram: \(X\to T_{\mathbb{C}}(M)/H_{1}(G)\) above \(X\to G\); the left vertical arrow is the identity and the right one is the natural map}

\DArticleAnchor{0055}{55}{56}{722}{59}{P10}

defines an extension of \(M^{\mathrm{an}}\) by the vector group \(F^{0}T_{\mathbb{C}}M\), hence an extension of M by this group. We must prove that it is the universal extension. For this, it is enough to verify that the category of extensions of \([X\to T_{\mathbb{C}}(M)/H_{1}(G)]\) by \(G_{a}\) is trivial (one object, no automorphism). It is also the category of extensions of \(T_{\mathbb{C}}(M)/T_{Z}(M)\) by \(G_{a}\), and indeed one has\par
\par\medskip
\DFormula{\operatorname{Ext}^{i}(T_{\mathbb{C}}(M)/T_{Z}(M),G_{a})=0\,\,\,\,\,for\,i=0,1.}
\par\medskip
Remark (10.1.9). — \(M^{\natural }\) is characterized by the commutative diagram\par
\par\medskip
\DAsset{P0055-A01}{diagram: \(X\to G^{\natural }\to G\) above \(T_{Z}(M)\to \operatorname{Lie}(G^{\natural })\to \operatorname{Lie}(G)\), with three vertical arrows directed upward}
\par\medskip
whose outer square is (10.1.3.1), and which induces an isomorphism \(T_{\mathbb{C}}(M)\xrightarrow{\sim }\operatorname{Lie}(G^{\natural })\).\par
\par\medskip
Variant (10.1.10). — A smooth 1-motive M over a scheme S consists of:\par
\par\medskip
a) a group scheme X over S which, locally for the étale topology, is a constant group scheme defined by a finitely generated free \(\mathbb{Z}-module\); an abelian scheme A over S, and a torus T over S;\par
b) an extension G of A by T;\par
\DFormula{c)\,a\,morphism\,u:X\to G.}
\par\medskip
For every integer n, a construction analogous to (10.1.5) associates with a smooth 1-motive M over S a group scheme \(T_{\mathbb{Z}/n\mathbb{Z}}(M)\), finite and flat over S. For \(\ell \) invertible on S, the projective system\par
\par\medskip
\DFormula{T_{\ell }(M)=\varprojlim \,T_{\mathbb{Z}/\ell ^{n}\mathbb{Z}}(M)}
\par\medskip
is a \(\mathbb{Z}_{\ell }-sheaf\) over S. In general, \(T_{\ell }(M)\) corresponds to an \(\ell -divisible\) \((=\) Barsotti–Tate) group over S.\par
\par\medskip
The filtration W of M, defined as in (10.1.4), defines a filtration W of \(T_{\ell }(M)\). One still has\par
\par\medskip
\DFormula{\operatorname{Gr}_{0}^{W}\,T_{\ell }(M)=X\otimes \widehat{\mathbb{Z}}}
\DFormula{\operatorname{Gr}_{−1}^{W}\,T_{\ell }(M)=T_{\ell }(A)}
\DFormula{\operatorname{Gr}_{−2}^{W}\,T_{\ell }(M)=Y\otimes \mathbb{Z}_{\ell }(1)}
\par\medskip
\DFormula{for\,Y=\operatorname{Hom}(T,G_{m})^{\vee }.}
\par\medskip
Likewise, a construction analogous to (10.1.7) associates with M a vector bundle \(T_{DR}(M)\).\par
\par\medskip
Terminology (10.1.11). — Let M be a 1-motive. T(M), \(T_{\ell }(M)\), and \(T_{DR}(M)\) are called the Hodge, \(\ell -adic\), and de Rham realizations of the 1-motive M.\par
\DAsset{P0055-A02}{\(\ell -adic\) realization, source-printed weight-graded pieces, and Terminology (10.1.11)}

\DArticleAnchor{0056}{56}{57}{723}{60}{P10}

10.2. 1-motives and biextensions\par
\par\medskip
The results of this subsection will not be used in the sequel of this section. We shall use the notion of biextension (S.G.A. 7, VII (2.1)) and the following generalization.\par
\par\medskip
(10.2.1) Let \(K_{i}:[A_{i}\to B_{i}]\) \((i=1,2)\) be two complexes of sheaves, concentrated in degree 0 and \(−1\) (on an arbitrary topos). A biextension of \(K_{1}\) and \(K_{2}\) by an abelian sheaf H consists of:\par
\par\medskip
a) a biextension P of \(B_{1}\) and \(B_{2}\) by H, i.e. an H-torsor P over \(B_{1}\times B_{2}\), with \(P_{b_{1},b_{2}}\) depending biadditively on \(b_{1}\) and \(b_{2}\);\par
b) a trivialization \((=\) biadditive section) of the biextension of \(B_{1}\) and \(A_{2}\) by H obtained by inverse image from P;\par
c) a trivialization of the biextension of \(A_{1}\) and \(B_{2}\) by H obtained by inverse image from P; the trivializations b) and c) are required to coincide over \(A_{1}\times A_{2}\).\par
\par\medskip
The category of biextensions of \(K_{1}\) and \(K_{2}\) by H is denoted \(\operatorname{Biext}(K_{1},K_{2};H)\). It is a Picard category (its objects may be added). The group \(\operatorname{Biext}^{0}(K_{1},K_{2};H)\) of automorphisms of any biextension of \(K_{1}\) and \(K_{2}\) by H is\par
\par\medskip
\DFormula{\operatorname{Biext}^{0}(K_{1},K_{2};H)=\operatorname{Hom}(H^{0}(K_{1})\otimes H^{0}(K_{2}),H).}
\par\medskip
The group of isomorphism classes of biextensions is denoted \(\operatorname{Biext}^{1}(K_{1},K_{2};H)\). As in (S.G.A. 7, VII (3.6.5)), one verifies that\par
\par\medskip
\DFormula{\operatorname{Biext}^{i}(K_{1},K_{2};H)=\operatorname{Ext}^{i}(K_{1}\otimes ^{L}\,K_{2};H)\,\,\,\,\,(i=0,1).}
\par\medskip
From this formula, or elementarily, one verifies that if \(K_{i}′\to K_{i}\) is a quasi-isomorphism \((i=1,2)\), giving a biextension of \(K_{1}\) and \(K_{2}\) by H is equivalent to giving a biextension of \(K_{1}′\) and \(K_{2}′\) by H.\par
\par\medskip
This still applies when “sheaves” is replaced by “algebraic groups” (respectively “complex analytic groups”): these may be interpreted as sheaves on a suitable site.\par
\par\medskip
(10.2.2) In this subsection, we shall identify a 1-motive \(M=(X,T,A,G,u)\) with a complex concentrated in degrees 0 and \(−1\)\par
\par\medskip
\DFormula{M:[X\to G].}
\par\medskip
Over \(\mathbb{C}\), \(M^{\mathrm{an}}\) will denote the complex of complex analytic groups \([X\to G^{\mathrm{an}}]\). One has the following rigidity result:\par
\par\medskip
Lemma (10.2.2.1). — Let \(M_{i}=(X_{i},A_{i},T_{i},G_{i},u_{i})\) \((i=1,2)\) be two 1-motives. One has\par
\par\medskip
\DFormula{\operatorname{Biext}^{0}(M_{1},M_{2};G_{m})=0.}
\par\medskip
Over \(\mathbb{C}\), one likewise has \(\operatorname{Biext}^{0}(M_{1}^{\mathrm{an}},M_{2}^{\mathrm{an}};G_{m})=0\).\par
\par\medskip
Indeed, every biadditive morphism \(G_{1}\times G_{2}\to G_{m}\) is trivial.\par
\DAsset{P0056-A01}{Section 10.2 and Definition (10.2.1), biextension data and \(\operatorname{Biext}^{0}\)}
\DAsset{P0056-A02}{\(\operatorname{Biext}^{1}\) and derived Ext, degree convention, and rigidity lemma}

\DArticleAnchor{0057}{57}{58}{724}{61}{P10}

Construction (10.2.3). — Let \(M_{i}=(X_{i},A_{i},T_{i},G_{i},u_{i})\) \((i=1,2)\) be two 1-motives over \(\mathbb{C}\). One has\par
\par\medskip
\DFormula{\operatorname{Biext}^{1}(M_{1},M_{2};G_{m})=\operatorname{Hom}(T(M_{1})\otimes T(M_{2}),\mathbb{Z}(1)).}
\par\medskip
We shall first compute \(\operatorname{Biext}^{1}(M_{1}^{\mathrm{an}},M_{2}^{\mathrm{an}};G_{m})\).\par
If V and W are two vector spaces, then in the analytic category one has\par
\par\medskip
\DFormula{(10.2.3.1)_{\mathrm{an}}\,\,\,\,\,\operatorname{Biext}^{0}(V,W;G_{m})\simeq _{\operatorname{exp}}\operatorname{Hom}(V\otimes W,\mathbb{C})}
\DFormula{(10.2.3.2)_{\mathrm{an}}\,\,\,\,\,\operatorname{Biext}^{1}(V,W;G_{m})=0.}
\par\medskip
Since an extension of W by \(G_{m}\), over an arbitrary base, is always locally trivial, it indeed follows formally that\par
\par\medskip
\DFormula{\operatorname{Biext}^{i}(V,W;G_{m})=\operatorname{Ext}^{i}(V,\operatorname{Hom}(W,G_{m}))}
\DFormula{=\operatorname{Ext}^{i}(V,W\ast )\,\,\,\,\,(i=1,2).}
\par\medskip
Let \(V_{Z}\) and \(W_{Z}\) be two free \(\mathbb{Z}-modules\) of finite rank. Since\par
\par\medskip
\DFormula{[V_{Z}\to V_{\mathbb{C}}]\to [0\to V_{\mathbb{C}}/V_{Z}]}
\par\medskip
is a quasi-isomorphism, the “inverse-image” functor is an equivalence.\par
\par\medskip
\DFormula{(10.2.3.3)_{\mathrm{an}}\,\,\,\,\,\operatorname{Biext}(V_{\mathbb{C}}/V_{Z},W_{\mathbb{C}}/W_{Z};G_{m})\simeq \operatorname{Biext}([V_{Z}\to V_{\mathbb{C}}],[W_{Z}\to W_{\mathbb{C}}];G_{m}).}
\par\medskip
Every biadditive morphism from \(V_{\mathbb{C}}/V_{Z}\times W_{\mathbb{C}}/W_{Z}\) to \(G_{m}\) is trivial, whence a) below:\par
\par\medskip
Lemma (10.2.3.4)\par
\par\medskip
\DFormula{a)\,\operatorname{Biext}^{0}(V_{\mathbb{C}}/V_{Z},W_{\mathbb{C}}/W_{Z};G_{m})=0.}
\DFormula{b)\,\operatorname{Biext}^{1}(V_{\mathbb{C}}/V_{Z},W_{\mathbb{C}}/W_{Z};G_{m})=\operatorname{Hom}(V_{Z}\otimes W_{Z},\mathbb{Z}(1)).}
\par\medskip
Let \(\psi _{i}:V_{\mathbb{C}}\otimes W_{\mathbb{C}}\to \mathbb{C}\) be two bilinear maps, with \(\psi =\psi _{1}−\psi _{2}\) taking values in \(\mathbb{Z}(1)=2\pi i\mathbb{Z}\) on \(V_{Z}\otimes W_{Z}\). Let \(P(\psi _{1},\psi _{2})\) be the following biextension of \([V_{Z}\to V_{\mathbb{C}}]\) and \([W_{Z}\to W_{\mathbb{C}}]\) by \(G_{m}\): the trivial biextension of \(V_{\mathbb{C}}\) and \(W_{\mathbb{C}}\) by \(G_{m}\), equipped with the trivializations \(\varphi _{1}=\operatorname{exp}\) \(\psi _{1}:V_{Z}\otimes W_{\mathbb{C}}\to G_{m}\) and \(\varphi _{2}=\operatorname{exp}\) \(\psi _{2}:V_{\mathbb{C}}\otimes W_{Z}\to G_{m}\). By (10.2.3.2), every biextension is of this form; by (10.2.3.1), \(P(\psi _{1},\psi _{2})\simeq P(\psi _{1}′,\psi _{2}′)\) if and only if \(\psi _{1}′−\psi _{2}′=\psi _{1}−\psi _{2}\). Thus the \(P(\psi ,0)\) represent every biextension once and only once.\par
\par\medskip
(10.2.3.5) Under the general hypotheses of (10.2.1), let \(K_{i}=[A_{i}+F_{i}\to B_{i}]\) \((i=1,2)\) be complexes. Let \(K_{i}′\) be the subcomplex \([A_{i}\to B_{i}]\). One verifies that giving a biextension of \(K_{1}\) and \(K_{2}\) by H is equivalent to giving:\par
\par\medskip
a) a biextension P of \(K_{1}′\) and \(K_{2}′\) by H; and\par
b) trivializations of the biextensions of \([0\to F_{1}]\) and \(K_{2}′\) (respectively \(K_{1}′\) and \([0\to F_{2}])\) by H obtained by inverse image from P; these trivializations must coincide over \(F_{1}\times F_{2}\).\par
\par\medskip
In the special case where\par
\par\medskip
\DFormula{\operatorname{Biext}^{i}(F_{1},K_{2}′;H)=0\,\,\,\,\,(i=1,2)}
and\par
\DFormula{\operatorname{Biext}^{i}(K_{1}′,F_{2};H)=0\,\,\,\,\,(i=1,2),}
\DAsset{P0057-A01}{Construction (10.2.3), analytic biextension computation, and pullback equivalence}
\DAsset{P0057-A02}{Lemma (10.2.3.4), bilinear forms, and opening of (10.2.3.5)}

\DArticleAnchor{0058}{58}{59}{725}{62}{P10}

there exists one and only one trivialization \(\varphi _{1}\), a section of P over \(F_{1}\times B_{2}\) (respectively \(\varphi _{2}\), a section of P over \(B_{1}\times F_{2})\), as in b). The category \(\operatorname{Biext}(K_{1},K_{2};H)\) is then identified with the full subcategory of \(\operatorname{Biext}(K_{1}′,K_{2}′;H)\) formed by the P for which \(\varphi _{1}=\varphi _{2}\) over \(F_{1}\times F_{2}\).\par
\par\medskip
(10.2.3.6) Let \(V_{Z}\) and \(W_{Z}\) be two free \(\mathbb{Z}-modules\) of finite type, and let F,G be vector subspaces of \(V_{\mathbb{C}}\) and \(W_{\mathbb{C}}\). One has\par
\par\medskip
\DFormula{\operatorname{Biext}^{i}(F,[W_{Z}\to W_{\mathbb{C}}];G_{m})=0.}
\par\medskip
For \(i=0\), this follows from (10.2.3.1) and from the fact that a bilinear map\par
\par\medskip
\(B:F\times W_{\mathbb{C}}\to \mathbb{C}\)     with     \(B(F,W_{Z})\subset \mathbb{Z}(1)\)\par
\par\medskip
is zero. For \(i=1\), it follows from the fact that every biadditive morphism \(F\times W_{Z}\to G_{m}\) comes from a biadditive morphism \(F\times W_{\mathbb{C}}\to G_{m}\). In other words, one uses the long exact sequence\par
\par\medskip
\DFormula{0\to \operatorname{Biext}^{0}(F,[W_{Z}\to W_{\mathbb{C}}];G_{m})\to \operatorname{Biext}^{0}(F,W_{\mathbb{C}};G_{m})\xrightarrow{\sim}\operatorname{Biext}^{0}(F,W_{Z};G_{m})}
\DFormula{\to \operatorname{Biext}^{1}(F,[W_{Z}\to W_{\mathbb{C}}];G_{m})\to \operatorname{Biext}^{1}(F,W_{\mathbb{C}};G_{m}).}
\par\medskip
We can therefore apply (10.2.3.5) to the complexes \(K_{1}=[V_{Z}\oplus F\to V_{\mathbb{C}}]\) and \(K_{2}=[W_{Z}\oplus G\to W_{\mathbb{C}}]\). We find that biextensions of \(K_{1}\) and \(K_{2}\) by \(G_{m}\) are identified with certain biextensions of \(V_{\mathbb{C}}/V_{Z}\) and \(W_{\mathbb{C}}/W_{Z}\) by \(G_{m}\), i.e. with certain forms \(\psi :V_{Z}\otimes W_{Z}\to \mathbb{Z}(1)\) (10.2.3.4). For \(P(\psi _{1},\psi _{2})\) (10.2.3.4) to come from a biextension of \(K_{1}\) and \(K_{2}\) by \(G_{m}\), it is necessary and sufficient that exp \(\psi _{2}:F\times W_{\mathbb{C}}\to G_{m}\) and exp \(\psi _{1}:V_{\mathbb{C}}\times G\to G_{m}\) coincide over \(F\times G\), that is, that \(\psi (F,G)=0\).\par
\par\medskip
Lemma (10.2.3.7). — With the preceding notation, one has \(\operatorname{Biext}^{0}(K_{1},K_{2};G_{m})=0\), and \(\operatorname{Biext}^{1}(K_{1},K_{2};G_{m})\) is identified with the set of forms \(\psi :V_{Z}\otimes W_{Z}\to \mathbb{Z}(1)\) whose complexifications satisfy \(\psi (F,G)=0\).\par
\par\medskip
For \(\psi \) as in (10.2.3.7), with \(\psi =\psi _{1}−\psi _{2}\), the biextension corresponding to \(\psi \) is given by the following trivializations of the trivial biextension of \(V_{\mathbb{C}}\) and \(W_{\mathbb{C}}\) by \(G_{m}\):\par
\par\medskip
\DFormula{over\,(V_{Z}\oplus F)\times W_{\mathbb{C}}:\,\operatorname{exp}\,\psi _{1}(v_{Z},w_{\mathbb{C}})·\operatorname{exp}\,\psi _{2}(f,w_{\mathbb{C}})}
\DFormula{over\,V_{\mathbb{C}}\times (W_{Z}\oplus G):\,\operatorname{exp}\,\psi _{2}(v_{\mathbb{C}},w_{Z})·\operatorname{exp}\,\psi _{1}(v_{\mathbb{C}},g).}
\par\medskip
(10.2.3.8) For the 1-motives \(M_{i}\), one has quasi-isomorphisms\par
\par\medskip
\DAsset{P0058-A02}{diagram: \(T_{Z}(M_{i})\oplus F^{0}T_{\mathbb{C}}(M_{i})\to T_{Z}(M_{i})\to X_{i}\) above \(T_{\mathbb{C}}(M_{i})\to \operatorname{Lie}(G_{i})\to G_{i}\), with three downward vertical arrows}
\DAsset{P0058-A01}{completion of (10.2.3.5), Lemma (10.2.3.6), and long exact sequence}

\DArticleAnchor{0059}{59}{60}{726}{63}{P10}

The group \(\operatorname{Biext}^{1}(M_{1}^{\mathrm{an}},M_{2}^{\mathrm{an}};G_{m})\) is therefore identified with the group of maps\par
\par\medskip
\DFormula{\psi :T_{Z}(M_{1})\otimes T_{Z}(M_{2})\to \mathbb{Z}(1)}
\par\medskip
whose complexification is compatible with the Hodge filtration.\par
\par\medskip
Lemma (10.2.3.9). — The map \(\operatorname{Biext}^{1}(M_{1},M_{2};G_{m})\to \operatorname{Biext}^{1}(M_{1}^{\mathrm{an}},M_{2}^{\mathrm{an}};G_{m})\) is injective. Its image is formed by the biextensions P whose restriction to \(G_{1}\times G_{2}\) is the inverse image of a biextension of \(A_{1}\) and \(A_{2}\) by \(G_{m}\).\par
\par\medskip
One has\par
\par\medskip
\DFormula{\operatorname{Biext}^{1}(A_{1},A_{2};G_{m})\simeq \operatorname{Biext}^{1}(G_{1},G_{2};G_{m})\,\,\,\,\,(S.G.A.\,7,\,VIII\,(3.5))}
\par\medskip
and\par
\par\medskip
\DFormula{\operatorname{Biext}^{1}(A_{1},A_{2};G_{m})\simeq \operatorname{Biext}^{1}(A_{1}^{\mathrm{an}},A_{2}^{\mathrm{an}};G_{m})\,\,\,\,\,(G.A.G.A.).}
\par\medskip
Let P be a biextension of \(M_{1}\) and \(M_{2}\) by \(G_{m}\). If \(P^{\mathrm{an}}\) is trivial, then:\par
\par\medskip
a) the restriction of P to \(G_{1}\times G_{2}\) is the image of a biextension \(P_{0}\) of \(A_{1}\) and \(A_{2}\) by \(G_{m}\); the latter is trivial. Otherwise it would be analytically nontrivial, would give rise to a nonzero bilinear form \(\psi _{0}\), and the form \(\psi \) corresponding to P would a fortiori be nonzero. The restriction of P to \(G_{1}\times G_{2}\) is therefore trivial (and uniquely trivializable, even analytically).\par
\par\medskip
b) P is therefore defined by biadditive maps \(X_{1}\times G_{2}\to G_{m}\) and \(G_{1}\times X_{2}\to G_{m}\). These are zero in the analytic category, hence are zero, and P is trivial.\par
\par\medskip
Let P be a biextension of \(M_{1}^{\mathrm{an}}\) and \(M_{2}^{\mathrm{an}}\) by \(G_{m}\) satisfying condition (10.2.3.9). It remains to prove that P is algebraizable. By hypothesis and G.A.G.A., its restriction \(P_{1}\) to \(G_{1}\times G_{2}\) is algebraizable. We conclude by showing that a trivialization of \(P_{1}\) over \(X_{1}\times G_{2}\) (respectively \(G_{1}\times X_{2})\) is automatically algebraic: \(P_{1}\) may be interpreted as an extension of \(X_{1}\otimes G_{2}\) by \(G_{m}\), which reduces the matter to showing that\par
\par\medskip
\DFormula{\operatorname{Ext}^{i}(G_{1},G_{m})\simeq \operatorname{Ext}^{i}(G_{1}^{\mathrm{an}},G_{m})\,\,\,\,\,(i=0,1).}
\par\medskip
This follows from the same assertion for an abelian variety (G.A.G.A.) and for \(G_{m}\).\par
\par\medskip
By (10.2.3.9), \(\operatorname{Biext}^{1}(M_{1},M_{2};G_{m})\) is identified with the set of\par
\par\medskip
\DFormula{\psi :T_{Z}(M_{1})\otimes T_{Z}(M_{2})\to \mathbb{Z}(1)}
\par\medskip
such that:\par
\par\medskip
a) \(\psi \) is compatible with the Hodge filtration (10.2.3.8);\par
b) the restriction of \(\psi \) to \(W_{−1}T_{Z}(M_{1})\otimes W_{−1}T_{Z}(M_{2})\) comes from a form on\par
\par\medskip
\DFormula{\operatorname{Gr}_{−1}^{W}\,T_{Z}(M_{1})\otimes \operatorname{Gr}_{−1}^{W}\,T_{Z}(M_{2}).}
\par\medskip
Condition b) means that \(\psi \) is compatible with W, and this completes Construction (10.2.3).\par
\par\medskip
Remark (10.2.4). — A biextension P of \(M_{1}\) and \(M_{2}\) by \(G_{m}\) defines in the evident way a biextension \(P^{0}\) of \(M_{2}\) and \(M_{1}\) by \(G_{m}\). If \(\psi \) (respectively \(\psi ^{0})\) is the bilinear form\par
\DAsset{P0059-A01}{analytic biextensions, Lemma (10.2.3.9), and injectivity proof}
\DAsset{P0059-A02}{algebraization proof, Hodge/weight conditions, and opening of Remark (10.2.4)}

\DArticleAnchor{0060}{60}{61}{727}{64}{P10}

defined by P (respectively \(P^{0})\), one has \(\psi (x,y)=−\psi ^{0}(y,x)\). This is clear from the formulas of (10.2.3.7).\par
\par\medskip
Construction (10.2.3) has purely algebraic analogues in \(\ell -adic\) or de Rham cohomology.\par
\par\medskip
Construction (10.2.5). — Let \(M_{i}=(X_{i},A_{i},T_{i},G_{i},u_{i})\) \((i=1,2)\) be two 1-motives over an algebraically closed field k of characteristic 0. Let P be a biextension of \(M_{1}\) and \(M_{2}\) by \(G_{m}\). With it one associates a morphism\par
\par\medskip
\DFormula{T_{\mathbb{Z}/n\mathbb{Z}}(M_{1})\otimes T_{\mathbb{Z}/n\mathbb{Z}}(M_{2})\to \mathbb{Z}/n\mathbb{Z}(1).}
\par\medskip
Let \(m_{i}\in T_{\mathbb{Z}/n\mathbb{Z}}(M_{i})\), represented by \((x_{i},g_{i})\), with \(u_{i}(x_{i})=ng_{i}\). We construct two isomorphisms \(a_{1}\) and \(a_{2}\) from \(P_{g_{1},g_{2}}^{\otimes n}\) to the trivial torsor \(G_{m}\):\par
\par\medskip
\DFormula{a_{1}:P_{g_{1},g_{2}}^{\otimes n}\simeq P_{ng_{1},g_{2}}=P_{u(x_{1}),g_{2}}\simeq G_{m}}
\DFormula{a_{2}:P_{g_{1},g_{2}}^{\otimes n}\simeq P_{g_{1},ng_{2}}=P_{g_{1},u(x_{2})}\simeq G_{m}.}
\par\medskip
Set \(a_{2}=\Phi (m_{1},m_{2})a_{1}\). One verifies that \(\Phi (m_{1},m_{2})\) does not depend on the choice of the \((x_{i},g_{i})\) and is an n-th root of unity. This is the desired form.\par
\par\medskip
Proposition (10.2.6). — Over \(\mathbb{C}\), the form (10.2.5) is obtained from the one constructed in (10.2.3) by reduction modulo n.\par
\par\medskip
This is verified using the formulas given in (10.2.3.7).\par
\par\medskip
Construction (10.2.7). — Let \(M_{i}=(X_{i},A_{i},T_{i},G_{i},u_{i})\) be two 1-motives over an algebraically closed field k. Let P be a biextension of \(M_{1}\) and \(M_{2}\) by \(G_{m}\). With it one associates a morphism\par
\par\medskip
\DFormula{T_{DR}(M_{1})\otimes _{k}T_{DR}(M_{2})\to k.}
\par\medskip
Recall the definitions (due to Grothendieck) of \(\natural -extensions\) and biextensions.\par
\par\medskip
(10.2.7.1) A \(\natural -extension\) of a smooth commutative group G by a smooth commutative group H consists of:\par
\par\medskip
a) an extension E of G by H; if \(\mu :G\times G\to G\) is addition, E is regarded as an H-torsor over G equipped with a morphism of H-torsors over \(G\times G\) (addition)\par
\par\medskip
\DFormula{\nu :\operatorname{pr}_{1}^{\ast }E+\operatorname{pr}_{2}^{\ast }E\to \mu ^{\ast }E;}
\par\medskip
b) a connection on the H-torsor E such that the map \(\nu \) is horizontal. This connection automatically has zero curvature.\par
\par\medskip
If G is an extension of an abelian variety by a torus, every extension of G by \(G_{m}\) admits a \(\natural -structure\); two \(\natural -structures\) differ by an invariant form on G (with values in the Lie algebra of \(G_{m})\). The Baer sum of two \(\natural -extensions\) is defined in the evident way.\par
\DAsset{P0060-A01}{completion of Remark (10.2.4), Construction (10.2.5), and modulo-n form}
\DAsset{P0060-A02}{Proposition (10.2.6), de Rham construction, and \(\natural -extension\) definition}

\DArticleAnchor{0061}{61}{62}{728}{65}{P11}

(10.2.7.2) Let P be a biextension of \(G_{1}\) and \(G_{2}\) by H. Regard P as an H-torsor P over \(G_{1}\times G_{2}\), equipped with morphisms\par
\par\medskip
\DFormula{\nu _{1}\,:\,\operatorname{pr}_{13}^{\ast }P+\operatorname{pr}_{23}^{\ast }P\to (\mu _{1}\times Id)^{\ast }P\,\,\,\,\,over\,\,\,\,\,G_{1}\times G_{1}\times G_{2}}
\DFormula{\nu _{2}\,:\,\operatorname{pr}_{12}^{\ast }P+\operatorname{pr}_{13}^{\ast }P\to (Id\times \mu _{2})^{\ast }P\,\,\,\,\,over\,\,\,\,\,G_{1}\times G_{2}\times G_{2}.}
\par\medskip
A \(\natural -1-structure\) on P is a connection on the H-torsor P, relative to \(G_{1}\times G_{2}\to G_{2}\) (the covariant derivative is defined only for vector fields parallel to \(G_{1})\), such that \(\nu _{1}\) and \(\nu _{2}\) are horizontal. \(\natural -2-structures\) are defined in the same way. A \(\natural -structure\) is the datum of a \(\natural -1-structure\) and a \(\natural -2-structure\). It is also a connection on the H-torsor P such that \(\nu _{1}\) and \(\nu _{2}\) are horizontal.\par
\par\medskip
The curvature R of a \(\natural -biextension\) P is the curvature of the connection on P. It is a translation-invariant 2-form on \(G_{1}\times G_{2}\) with values in Lie(H), i.e. an alternating form on \(\operatorname{Lie}(G_{1})\times \operatorname{Lie}(G_{2})\) with values in Lie(H). Its restriction to \(\operatorname{Lie}(G_{1})\) and to \(\operatorname{Lie}(G_{2})\) is zero, so that R defines a pairing (also called the curvature of P)\par
\par\medskip
\DFormula{(10.2.7.3)\,\,\,\,\,\Phi \,:\,\operatorname{Lie}(G_{1})\otimes \operatorname{Lie}(G_{2})\to \operatorname{Lie}(H),}
\par\medskip
with\par
\par\medskip
\DFormula{R(g_{1}+g_{2},g_{1}′+g_{2}′)=\Phi (g_{1},g_{2}′)−\Phi (g_{1}′,g_{2}).}
\par\medskip
For complexes of smooth groups \(K_{i}:[A_{i}\to B_{i}]\), a \(\natural -biextension\) (respectively a \(\natural -i-biextension)\) of \(K_{1}\) and \(K_{2}\) by H is a \(\natural -biextension\) (respectively a \(\natural -i-biextension)\) of \(B_{1}\) and \(B_{2}\) by H, equipped with trivializations (as \(\natural -biextensions\), respectively as \(\natural -i-biextensions)\) as in (10.2.1).\par
\par\medskip
Proposition (10.2.7.4). — Let \(P^{\natural }\) be the biextension of \(M_{1}^{\natural }\) and \(M_{2}^{\natural }\) by \(G_{m}\) obtained by pulling back P. There exists on \(P^{\natural }\) one and only one \(\natural -structure\).\par
\par\medskip
Let us prove uniqueness: it must be checked that on the trivial biextension of \(M_{1}^{\natural }\) and \(M_{2}^{\natural }\) by \(G_{m}\) every \(\natural -structure\) is trivial. A \(\natural -structure\) is first a connection on the trivial \(G_{m}-torsor\) over \(G_{1}^{\natural }\times G_{2}^{\natural }\), i.e. a field of forms on \(G_{1}^{\natural }\times G_{2}^{\natural }\), say \(\Gamma _{g_{1},g_{2}}(t_{1}+t_{2})\). The axiom for \(\natural -structures\) gives\par
\par\medskip
\DFormula{\Gamma _{g_{1},g_{2}}(t_{1}+t_{2})=\Phi _{1}(g_{1},t_{2})+\Phi _{2}(t_{1},g_{2})}
\par\medskip
with \(\Phi _{1}\) and \(\Phi _{2}\) biadditive. \(G_{i}^{\natural }\) is an extension of \(X_{i}\otimes k\) by the universal extension \(G_{i}′\) of \(G_{i}\) by an additive group. Every morphism \(G_{i}′\to G_{a}\) is trivial. Hence \(\Phi _{1}\) factors through \((X_{1}\otimes k)\otimes \operatorname{Lie}(G_{2}^{\natural })\). Finally, the restriction of \(\Gamma \) to \(X_{1}\times G_{2}^{\natural }\) must be trivial, and therefore so must the restriction of \(\Phi _{1}\) to \(X_{1}\times \operatorname{Lie}(G_{2}^{\natural })\). Thus \(\Phi _{1}=0\). Likewise \(\Phi _{2}=0\), and \(\Gamma =0\).\par
\par\medskip
To prove existence, it is enough to construct a \(\natural -2-structure\) on the biextension of \(M_{1}^{\natural }\) and \(M_{2}\) by \(G_{m}\) obtained by pulling back P: by pullback, this will give a \(\natural -2-structure\) on \(P^{\natural }\) and, symmetrically, a \(\natural -1-structure\).\par
\par\medskip
For \(g\in G_{1}\), \(P_{g}\) is an extension of \(G_{2}\) by \(G_{m}\). Let \(C_{g}\) be the set of \(\natural -structures\) on \(P_{g}\). It is a torsor under \(\operatorname{Lie}(G_{2})^{\ast }\). For \(g\in G_{1}\), the \(C_{g}\) are the fibers of a torsor C over \(G_{1}\); the Baer sum of \(\natural -extensions\) even makes C an extension of \(G_{1}\) by\par
\DAsset{P0061-A01}{definition of natural biextensions, curvature, and pairing (10.2.7.2-3)}
\DAsset{P0061-A02}{Proposition (10.2.7.4), uniqueness, and the opening of the existence construction}

\DArticleAnchor{0062}{62}{63}{729}{66}{P11}

\(\operatorname{Lie}(G_{2})^{\ast }\). Lift \(u_{1}:X\to G_{1}\) to \(u_{1}′:X\to C\) by associating with \(x\in X\) the trivial connection on the trivialized extension \(P_{u_{1}(x)}\). In this way one obtains a \(\natural -2-structure\) on the biextension of \([X\to C]\) and \(M_{2}\) by \(G_{m}\) obtained by pulling back P.\par
\par\medskip
By the universal property of \(M_{1}^{\natural }\), there exists a unique commutative diagram\par
\par\medskip
\DAsset{P0062-A01}{triangular diagram: \([X\to G]\) at the top; \(M_{1}^{\natural }\) at lower left; \([X\to C]\) at lower right; an arrow v from \(M_{1}^{\natural }\) to \([X\to C]\) and two upward arrows to \([X\to G]\)}
\par\medskip
and, on pulling back by v, one finds the desired \(\natural -2-structure\).\par
\par\medskip
By definition, the pairing (10.2.7) is the negative of the curvature of the \(\natural -biextension\) (10.2.7.4).\par
\par\medskip
Proposition (10.2.8). — Over \(\mathbb{C}\), the pairing (10.2.7) is the complexification of the pairing (10.2.3).\par
\par\medskip
The \(\natural -biextension\) \(P^{\natural }\) defines, in the analytic category, a \(\natural -biextension\) of \([T_{Z}(M_{1})\to T_{\mathbb{C}}(M_{1})]\) and \([T_{Z}(M_{2})\to T_{\mathbb{C}}(M_{2})]\) by \(G_{m}\). Proposition (10.2.8) then follows from the\par
\par\medskip
Lemma (10.2.9). — Let \(V_{Z}\) and \(W_{Z}\) be two finite-rank free \(\mathbb{Z}-modules\), and let P be a \(\natural -biextension\) of \([V_{Z}\to V_{\mathbb{C}}]\) and \([W_{Z}\to W_{\mathbb{C}}]\) by \(G_{m}\), with underlying biextension P.\par
\par\medskip
Then the curvature \(\psi :V_{\mathbb{C}}\otimes W_{\mathbb{C}}\to \mathbb{C}\) of \(P^{\natural }\) coincides with the negative of the complexification of the form corresponding to P under (10.2.3.4).\par
\par\medskip
Take \(P=P(\psi _{1},\psi _{2})\) (10.2.3.4). A \(\natural -structure\) on P is first a connection on the trivial \(G_{m}-torsor\) over \(V_{\mathbb{C}}\times W_{\mathbb{C}}\), i.e. a field of forms \(\Gamma _{v,w}(v′+w′)\) on \(V_{\mathbb{C}}\times W_{\mathbb{C}}\). This connection defines a \(\natural -structure\) on this torsor if and only if\par
\par\medskip
\DFormula{\Gamma _{v,w}(v′+w′)=\Phi _{1}(v,w′)+\Phi _{2}(v′,w)}
\par\medskip
with \(\Phi _{1}\) and \(\Phi _{2}\) bilinear. The trivializations exp \(\psi _{1}(v_{Z},w_{\mathbb{C}})\) and exp \(\psi _{2}(v_{\mathbb{C}},w_{Z})\) are horizontal if and only if \(\Phi _{i}=−\psi _{i}\). The curvature of the connection is the field of 2-forms \(d\Gamma \):\par
\par\medskip
\DFormula{R_{v,w}(v′+w′,v″+w″)=\Phi _{1}(v′,w″)+\Phi _{2}(v″,w′)−\Phi _{1}(v″,w′)−\Phi _{2}(v′,w″).}
\par\medskip
The curvature of P is therefore the form\par
\par\medskip
\DFormula{\Phi _{1}(v′,w″)−\Phi _{2}(v′,w″)=−(\psi _{1}(v′,w″)−\psi _{2}(v′,w″))=−\psi (v′,w″).}
\par\medskip
(10.2.10) Let M be a 1-motive over \(\mathbb{C}\). To M corresponds a mixed Hodge structure T(M) of type \(\lbrace (0,0)\,(0,−1)\,(−1,0)\,(−1,−1)\rbrace \). The mixed Hodge structure \(\operatorname{Hom}(T(M),\mathbb{Z}(1))\) is again of this type, and therefore defines a 1-motive M*. The evident map \(T(M)\otimes T(M\ast )\to \mathbb{Z}(1)\) defines a biextension P of M and M* by \(G_{m}\). The 1-motive M*, equipped with the biextension P, is the Cartier dual of M. P is called the Poincaré biextension.\par
\DAsset{P0062-A02}{Lemma (10.2.9), curvature calculation, and Cartier duality (10.2.10)}

\DArticleAnchor{0063}{63}{64}{730}{67}{P11}

(10.2.11) Construction (10.2.10) can be made purely algebraic. To a 1-motive \(M=(X,A,T,G,u)\) over an algebraically closed field k, we shall associate a Cartier dual \(M\ast =(X′,A′,T′,G′,u′)\):\par
\par\medskip
a) X′ is the character group of T, A′ is the abelian variety dual to A, and T′ is the torus whose character group is X.\par
\par\medskip
b) First consider the case \(T=0\). An extension of M by \(G_{m}\) then consists of an extension E of A by \(G_{m}\), together with a trivialization \(\widetilde{u}\) of the extension \(u^{\ast }E\) of X by \(G_{m}\)\par
\par\medskip
\DAsset{P0063-A01}{diagram: exact sequence \(0\to G_{m}\to E\to A\to 0\); X above A, with arrow u to A and arrow \(\widetilde{u}\) to E}
\par\medskip
There is an exact sequence\par
\par\medskip
\DFormula{0\to \operatorname{Hom}(X,G_{m})\to \operatorname{Ext}^{1}(M,G_{m})\to \operatorname{Ext}^{1}(A,G_{m})\to 0.}
\par\medskip
Extensions of M by \(G_{m}\) have no nontrivial automorphism. The functor of isomorphism classes of extensions is representable; define G′ to be the algebraic group representing this functor. It is an extension of A′ by T′. There is a biextension P of M and G′ by \(G_{m}\).\par
\par\medskip
c) In the general case, set \(G′=\operatorname{Ext}^{1}(M/W_{−2}M,G_{m})\). The natural biextension P″ of \(M/W_{−2}M\) and G′ by \(G_{m}\) induces a biextension P′ of M and G′ by \(G_{m}\). For each \(\chi \in X′\), the extension M of \(M/W_{−2}M\) by T defines an extension of \(M/W_{−2}M\) by \(G_{m}\), and hence an element \(u′(\chi )\in G′\).\par
\par\medskip
\DAsset{P0063-A02}{first diagram: exact rows \(0\to T\to M\to M/W_{−2}M\to 0\) and \(0\to G_{m}\to P″_{u′(\chi )}\to M/W_{−2}M\to 0\); vertical arrows \(T\to G_{m}\), \(M\to P″_{u′(\chi )}\) labeled v, and the identity on \(M/W_{−2}M\). Second diagram: extension \(0\to G_{m}\to P′_{u′(\chi )}\to M\to 0\) above \(G_{m}\to P″_{u′(\chi )}\to M/W_{−2}M\), with equality at left, vertical arrows, and a diagonal v from M to \(P″_{u′(\chi )}\)}
\par\medskip
The map v trivializes the extension \(P′_{u′(\chi )}\) of M by \(G_{m}\).\par
\par\medskip
These trivializations make P′ into a biextension of M and \(M\ast =[X′\xrightarrow{u′}G′]\) by \(G_{m}\). The 1-motive M*, equipped with P′, is the desired Cartier dual.\par

\DArticleAnchor{0064}{64}{65}{731}{68}{P11}

(10.2.12) This construction gives a more symmetric description of 1-motives.\par
\par\medskip
Let \(M=(X,A,T,G,u)\) be a 1-motive, with Cartier dual \(M\ast =(X′,A′,T′,G′,u′)\), and let \(P_{0}\) be the Poincaré biextension. Passing to the quotient, the morphisms u and u′ define morphisms\par
\par\medskip
(10.2.12.1)     \(\overline{u}:X\to A\)     and     \(\overline{u}′:X′\to A′\).\par
\par\medskip
The biextension of G and G′ by \(G_{m}\) deduced from \(P_{0}\) is the pullback of the Poincaré biextension P of A and A′ by \(G_{m}\). By hypothesis, the trivializations (of the pullbacks) of P over \(X\times G′\) and \(G\times X′\) agree over \(X\times X′\), whence\par
\par\medskip
\DFormula{(10.2.12.2)\,\,\,\,\,a\,trivialization\,\psi \,of\,(\overline{u}\times \overline{u}′)^{\ast }P.}
\par\medskip
A morphism \(f:M_{1}\to M_{2}\), with transpose (in an evident sense) \(M_{2}\ast \to M_{1}\ast \), gives rise to commutative diagrams\par
\par\medskip
\DAsset{P0064-A01}{(10.2.12.3): square \(X_{1}\to A_{1}\) above \(X_{2}\to A_{2}\), with downward arrows \(f_{X}\) and \(f_{A}\); square \(X_{1}′\to A_{1}′\) above \(X_{2}′\to A_{2}′\), with upward arrows \(f_{X}′\) and \(f_{A}′\)}
\par\medskip
Moreover, \(f_{A}\) and \(f_{A}′\) are transposes and, via the isomorphism \((1,f_{A}′)^{\ast }P_{1}\simeq (f_{A},1)^{\ast }P_{2}\), one has\par
\par\medskip
\DFormula{(10.2.12.4)\,\,\,\,\,\psi _{1}(x_{1},f_{X}′(x_{2}′))=\psi _{2}(f_{X}(x_{1}),x_{2}′).}
\par\medskip
\DAsset{P0064-A02}{large diagram: \(X_{1}\times X_{1}′\to (P_{1}\) over \(A_{1}\times A_{1}′)\) by \(\psi _{1}\); \(X_{1}\times X_{2}′\) to \((1,f_{A}′)^{\ast }P_{1}=(f_{A},1)^{\ast }P_{2}\) over \(A_{1}\times A_{2}′\); \(X_{2}\times X_{2}′\to (P_{2}\) over \(A_{2}\times A_{2}′)\) by \(\psi _{2}\); upward vertical arrows above the middle row and downward vertical arrows below it}
\par\medskip
(10.2.13) One thus obtains a functor from the category of 1-motives to the following category. Its objects consist of:\par
\par\medskip
a) two abelian varieties A and A′ in duality;\par
b) two finite-rank free \(\mathbb{Z}-modules\) X and X′;\par
c) morphisms \(v:X\to A\) and \(v′:X′\to A′\);\par
d) a trivialization \(\psi \) of the pullback by (v,v′) of the Poincaré biextension of A and A′ by \(G_{m}\).\par

\DArticleAnchor{0065}{65}{66}{732}{69}{P11}

The arrows consist of the diagrams (10.2.12.3) satisfying the conditions of (10.2.12).\par
\par\medskip
Proposition (10.2.14). — The functor (10.2.13) is an equivalence of categories.\par
\par\medskip
The verification is left to the reader as an exercise.\par
\par\medskip
10.3. Algebraic interpretation of mixed \(H^{1}\): the case of curves\par
\par\medskip
(10.3.1) Let \(X_{0}\) be an algebraic curve over an algebraically closed field k. Let \(p:X′\to X_{0}\) be the normalization of \(X_{0}\), and let \(\overline{X}′\) be the nonsingular projective curve in which X′ is a dense open set. Denote by X (respectively \(\overline{X})\) the curve obtained from X′ (respectively from \(\overline{X}′)\) by contracting to a point each finite set \(p^{−1}(s)\), for \(s\in X_{0}\).\par
\par\medskip
\DAsset{P0065-A01}{normalization and compactification diagram: \(X′\hookrightarrow \overline{X}′\) by j′; vertical arrows \(r:X′\to X\) and \(\overline{r}:\overline{X}′\to \overline{X}\); inclusion \(j:X\hookrightarrow \overline{X}\); arrow \(p_{0}:X\to X_{0}\)}
\par\medskip
X and \(\overline{X}\) may be characterized by the following properties.\par
\par\medskip
a) \(p_{0}\) is finite and radicial, and therefore induces isomorphisms\par
\par\medskip
\(H^{\bullet }(X_{0},\mathbb{Z}_{\ell })\to H^{\bullet }(X,\mathbb{Z}_{\ell })\)     \((\ell \) prime to the characteristic).\par
\par\medskip
b) The singularities of X are analytically isomorphic to those presented by a union of coordinate axes in affine space.\par
\par\medskip
c) \(\overline{X}\) is projective, and \(\overline{X}−X\) is a finite set of points at which \(\overline{X}\) is smooth.\par
\par\medskip
(10.3.2) The Picard group \(\operatorname{Pic}(\overline{X})\) is described as follows.\par
\par\medskip
a) Let I be the set of irreducible components of \(\overline{X}\) \((=\) the set of irreducible components of \(X_{0})\). For each irreducible component Y of \(\overline{X}\), the degree of the restriction to Y of an invertible sheaf is a homomorphism\par
\par\medskip
\DFormula{\operatorname{deg}_{Y}:\operatorname{Pic}(\overline{X})\to \mathbb{Z}\,:\,\mathcal{L}\to \operatorname{deg}_{Y}(\mathcal{L}).}
\par\medskip
This gives\par
\par\medskip
\DFormula{\operatorname{deg}:\operatorname{Pic}(\overline{X})\to \mathbb{Z}^{I}.}
\par\medskip
This homomorphism is surjective, with kernel \(\operatorname{Pic}^{0}(\overline{X})\).\par
\par\medskip
b) The morphism from \(\operatorname{Pic}^{0}(\overline{X})\) to the abelian variety \(\operatorname{Pic}^{0}(\overline{X}′)\) is surjective. Its kernel is a torus.\par
\DAsset{P0065-A02}{properties of X and X-bar and the Picard-group decomposition}

\DArticleAnchor{0066}{66}{67}{733}{70}{P11}

(10.3.3) Let S be the finite set \(\overline{X}−X\), let \(\alpha _{0}:S\to I\) be the map associating with \(s\in S\) the unique irreducible component of \(\overline{X}\) containing s, and let \(\alpha :\mathbb{Z}^{S}\to \mathbb{Z}^{I}\) be induced by \(\alpha _{0}\).\par
\par\medskip
Each \(s\in S\) defines an invertible sheaf \(\mathcal{O}(s)\), whence \(u_{0}:\mathbb{Z}^{S}\to \operatorname{Pic}(\overline{X})\). One has deg \(u_{0}=\alpha \).\par
\par\medskip
Definition (10.3.4). — The motivic \(H^{1}\) of \(X_{0}\), denoted \(H_{m}^{1}(X_{0})(1)\), is the 1-motive\par
\par\medskip
\DFormula{\operatorname{Ker}(\alpha )\to \operatorname{Pic}^{0}(\overline{X}).}
\par\medskip
\DFormula{By\,definition,\,one\,therefore\,has\,H_{m}^{1}(X_{0})(1)=H_{m}^{1}(X)(1).}
\par\medskip
Remarks (10.3.5). — (i) One has\par
\par\medskip
\DFormula{\operatorname{Gr}_{0}^{W}(H_{m}^{1}(X)(1))=\operatorname{Ker}(\alpha :\mathbb{Z}^{S}\to \mathbb{Z}^{I})}
\DFormula{\operatorname{Gr}_{−1}^{W}(H_{m}^{1}(X)(1))=\operatorname{Pic}^{0}(\overline{X}′).}
\par\medskip
(ii) The cokernel of \(\alpha \) is torsion-free: it identifies with the free \(\mathbb{Z}-module\) with basis \(I−\alpha _{0}(S)\). There is a short exact sequence of complexes and a quasi-isomorphism\par
\par\medskip
\DAsset{P0066-A02}{source-printed diagram: \(0\to [\operatorname{Ker}(\alpha )\to \operatorname{Pic}^{0}(\overline{X}_{0})]\to [\mathbb{Z}^{S}\to \operatorname{Pic}(\overline{X}_{0})]\to [\operatorname{Im}(\alpha )\to \mathbb{Z}^{I}]\to 0\), with a vertical arrow marked \(\simeq \) from \([\operatorname{Im}(\alpha )\to \mathbb{Z}^{I}]\) to \([0\to \operatorname{coker}(\alpha )]\)}
\par\medskip
Construction (10.3.6). — For an integer \(n>0\) prime to the characteristic of k, we construct an isomorphism\par
\par\medskip
\DFormula{H^{1}(X,\mathbb{Z}/n\mathbb{Z})(1)=H^{1}(X,\mu _{n})\simeq T_{\mathbb{Z}/n\mathbb{Z}}(H_{m}^{1}(X)(1)).}
\par\medskip
It is known that \(H^{1}(X,\mu _{n})\) identifies with the set of isomorphism classes of invertible sheaves \(\mathcal{L}\) on X equipped with an isomorphism \(\mathcal{L}^{\otimes n}\simeq \mathcal{O}\). Since \(p_{0}:X_{0}\to X\) is finite and radicial, one has \(H^{1}(X,\mathbb{Z}/n)\simeq H^{1}(X_{0},\mathbb{Z}/n)\), and \(H^{1}(X,\mu _{n})\) is again the set of isomorphism classes of pairs \((\mathcal{L},\alpha )\), where \(\mathcal{L}\) is an invertible sheaf on \(X_{0}\) and \(\alpha \) is an isomorphism \(\alpha :\mathcal{L}^{\otimes n}\simeq \mathcal{O}\).\par
\par\medskip
Let \((\mathcal{L},\alpha )\) be such a pair. The invertible sheaf \(\mathcal{L}\) extends to an invertible sheaf \(\overline{\mathcal{L}}\) on \(\overline{X}\), and there exists a divisor D supported on S such that \(\alpha \) extends to \(\overline{\alpha }:\overline{\mathcal{L}}^{\otimes n}\simeq \mathcal{O}(D)\). If there exists an isomorphism \(\beta \) of \(\overline{\mathcal{L}}^{\otimes n}\) with \(\mathcal{O}(E)\), this isomorphism is uniquely determined up to multiplication by an element of the divisible group \(H^{0}(\overline{X},\mathcal{O}^{\ast })\). It follows that \((\overline{\mathcal{L}},D)\) determines \((\mathcal{L},\alpha )\) up to isomorphism. For a pair \((\overline{\mathcal{L}},D)\) to come from a suitable pair \((\mathcal{L},\alpha )\), it is necessary and sufficient that \(n[\overline{\mathcal{L}}]=[\mathcal{O}(D)]\) in \(\operatorname{Pic}(\overline{X})\). It comes from \((\mathcal{O}_{X_{0}},0)\) if and only if it is of the form \((\mathcal{O}_{X_{0}}(D),nD)\).\par
\DAsset{P0066-A01}{motivic \(H^{1}\), weight-graded pieces, and the exact-complex diagram}

\DArticleAnchor{0067}{67}{68}{734}{71}{P12}

This identifies \(H^{1}(X,\mu _{n})\) with the \(H^{0}\) of the tensor-product complex of \([\mathbb{Z}^{S}\to \operatorname{Pic}(\overline{X})]\) (degrees 0 and 1) and \([\mathbb{Z}\xrightarrow{n}\mathbb{Z}]\) (degrees \(−1\) and 0):\par
\par\medskip
\DAsset{P0067-A01}{source-printed diagram: square of complexes \(\mathbb{Z}^{S}\xrightarrow{u_{0}}\operatorname{Pic}(\overline{X})\), with upward vertical arrows labelled n and \(−n\) between two copies of the complex}
\par\medskip
\DFormula{H^{1}(X,\mu _{n})\simeq H^{0}([\mathbb{Z}^{S}\to \operatorname{Pic}(\overline{X})]\otimes ^{L}\mathbb{Z}/n\mathbb{Z}).}
\par\medskip
Since \(\operatorname{coker}(\alpha )\) is torsion-free, it follows from (10.3.5) (ii) that\par
\par\medskip
\DFormula{H^{1}(X,\mu _{n})\simeq H^{0}([\mathbb{Z}^{S}\to \operatorname{Pic}(\overline{X})]\otimes ^{L}\mathbb{Z}/n\mathbb{Z})\simeq H^{0}(H_{m}^{1}(X)\otimes ^{L}\mathbb{Z}/n\mathbb{Z})}
\DFormula{\simeq T_{\mathbb{Z}/n\mathbb{Z}}(H_{m}^{1}(X)(1)).}
\par\medskip
Variant (10.3.7). — For \(\ell \) prime to the characteristic p of k, one has\par
\par\medskip
\DFormula{H^{1}(X_{0},\mathbb{Z}_{\ell }(1))\simeq T_{\ell }(H_{m}^{1}(X_{0})(1)).}
\par\medskip
Construction (10.3.8). — Let \(X_{0}\) be a curve over \(\mathbb{C}\). We shall construct an isomorphism of mixed Hodge structures\par
\par\medskip
\DFormula{H^{1}(X_{0})(1)\simeq T(H_{m}^{1}(X_{0})(1)).}
\par\medskip
We may and shall suppose that \(X=X_{0}\).\par
\par\medskip
Construction (10.3.9). — (i) Let \(j_{\ast }^{m}\mathcal{O}_{X}^{\ast }\) be the subsheaf of meromorphic functions in \(j_{\ast }\mathcal{O}_{X}^{\ast }\). One has\par
\par\medskip
\DFormula{H^{1}(X,\mathbb{Z}(1))=H^{1}(\overline{X},[\mathcal{O}_{\overline{X}}\xrightarrow{\operatorname{exp}}j_{\ast }^{m}\mathcal{O}_{X}^{\ast }]).}
\par\medskip
\DFormula{(ii)\,H^{1}(X,\mathbb{C})=H^{1}(\overline{X},[\mathcal{O}_{\overline{X}}\xrightarrow{d}\overline{r}_{\ast }\Omega ^{1}_{\overline{X}′}(\operatorname{log}\,S)]).}
\par\medskip
(iii) The inclusion of \(H^{1}(X,\mathbb{Z})\) in \(H^{1}(X,\mathbb{C})\) is defined by the morphism of complexes\par
\par\medskip
\DAsset{P0067-A02}{small square: \(\mathcal{O}_{\overline{X}}\xrightarrow{\operatorname{exp}}j_{\ast }^{m}\mathcal{O}_{X}^{\ast }\) above \(\mathcal{O}_{\overline{X}}\to \overline{r}_{\ast }\Omega ^{1}_{\overline{X}′}(\operatorname{log}\) S), with vertical arrows \(\iota \) and df/f}
\par\medskip
Consider the commutative diagram\par
\par\medskip
\DAssetReference{P0067-A02}{three squares numbered \(\text{\textcircled{1}}\), \(\text{\textcircled{2}}\), and \(\text{\textcircled{3}}\), linking \(H^{1}(X,\mathbb{Z}(1))\), the hypercohomology of the exponential complexes on X and \(\overline{X}\), and then \(H^{1}(X,\mathbb{C})\) and the de Rham complexes; the horizontal arrows carry the printed isomorphisms}

\DArticleAnchor{0068}{68}{69}{735}{72}{P12}

The horizontal arrows of square 1 are isomorphisms, because the sequences\par
\par\medskip
\DFormula{0\to \mathbb{Z}(1)\to \mathcal{O}_{X}\xrightarrow{\operatorname{exp}}\mathcal{O}_{X}^{\ast }\to 0}
\par\medskip
\DFormula{0\to \mathbb{C}\to \mathcal{O}_{X}\xrightarrow{d}r_{\ast }\Omega ^{1}_{X′}\to 0}
\par\medskip
are exact. Those of square 2 are isomorphisms because \(R^{1}j_{\ast }\mathcal{O}_{X}=0\). Those of square 3 are isomorphisms because the morphisms of complexes defining them are quasi-isomorphisms. This diagram defines (10.3.9).\par
\par\medskip
(10.3.10) Recall that if \(d:F\to G\) is a morphism of sheaves of abelian groups on a space Z, then \(H^{1}(Z,[F\to G])\) identifies with the set of isomorphism classes of pairs \((P,\alpha )\), where P is an F-torsor \((=\) a principal homogeneous space under F) and \(\alpha \) is a trivialization of the G-torsor dP.\par
\par\medskip
In particular:\par
\par\medskip
a) \(H^{1}(\overline{X},[\mathcal{O}_{\overline{X}}^{\ast }\to \overline{r}_{\ast }\Omega ^{1}_{\overline{X}′}(\operatorname{log}\) S)]) identifies with the set of isomorphism classes of pairs \((\mathcal{L},\alpha )\), where \(\mathcal{L}\) is an invertible sheaf on \(\overline{X}\) and \(\alpha \) is a connection on \(\overline{r}^{\ast }\mathcal{L}\), holomorphic on X′ and allowed to have a simple pole at the points of S. Let a be the projection of \(\overline{X}\) onto \(\operatorname{Spec}(\mathbb{C})\) \((=\) a point). In the fppf topology, the sheaf \(R^{1}a_{\ast }([\mathcal{O}_{\overline{X}}^{\ast }\to \overline{r}_{\ast }\Omega ^{1}_{\overline{X}′}(\operatorname{log}\) S)]) is represented by a group scheme \(\operatorname{Pic}^{\natural }(X)\), whose set of points is \(H^{1}(\overline{X},[\mathcal{O}_{\overline{X}}^{\ast }\to \overline{r}_{\ast }\Omega ^{1}_{\overline{X}′}(\operatorname{log}\) S)]). \(\operatorname{Pic}^{\natural }(X)\) is an extension of a subgroup of \(\operatorname{Pic}(\overline{X})\), containing \(\operatorname{Pic}^{0}(\overline{X})\), by \(H^{0}(\overline{X}′,\Omega ^{1}_{\overline{X}′}(\operatorname{log}\) S)).\par
\par\medskip
One defines a map\par
\par\medskip
\DFormula{(10.3.10.1)\,\,\,\,\,\mathbb{Z}^{S}\to \operatorname{Pic}^{\natural }(X)}
\par\medskip
by associating to each divisor D of \(\overline{X}\) supported on S the invertible sheaf \(\mathcal{O}(D)\).\par
\par\medskip
b) \(H^{1}(X,\mathbb{C})=H^{1}(\overline{X},[\mathcal{O}_{\overline{X}}\to \overline{r}_{\ast }\Omega ^{1}_{\overline{X}′}(\operatorname{log}\) S)]) is the Lie algebra of \(\operatorname{Pic}^{\natural }(X)\). It is the set of isomorphism classes of pairs \((P,\alpha )\), where P is an \(\mathcal{O}_{\overline{X}}-torsor\) and \(\alpha \) is a connection on P as in a).\par
\par\medskip
c) \(H^{1}(X,\mathbb{Z}(1))=H^{1}(\overline{X},[\mathcal{O}_{\overline{X}}\to j_{\ast }^{m}\mathcal{O}_{X}^{\ast }])\) is the set of isomorphism classes of pairs \((P,\alpha )\), where P is an \(\mathcal{O}_{\overline{X}}-torsor\) and \(\alpha \) is an isomorphism from the invertible sheaf exp(P) to an invertible sheaf \(\mathcal{O}_{X}(D)\) (D supported on S). The map\par
\par\medskip
\DFormula{\operatorname{Aut}(P)=\mathbb{C}\to \operatorname{Aut}(\operatorname{exp}(P))=\mathbb{C}^{\ast }}
\par\medskip
is surjective. This again permits \(H^{1}(X,\mathbb{Z}(1))\) to be identified with the set of pairs (p,d), where p is an isomorphism class of \(\mathcal{O}_{\overline{X}}-torsors\), i.e. an element of Lie(Pic \(\overline{X})\), and \(d\in \operatorname{Ker}(\alpha )\) defines a divisor supported on S whose image in \(\operatorname{Pic}^{0}(\overline{X})\) is exp(p). In other words, an isomorphism has been defined\par
\par\medskip
\DFormula{(10.3.10.2)\,\,\,\,\,H^{1}(X,\mathbb{Z}(1))\simeq T_{\mathbb{Z}}(H_{m}^{1}(X)(1)).}
\par\medskip
It follows from part (i) of the next lemma that this isomorphism is compatible with the weight filtration.\par
\DAsset{P0068-A01}{exact sequences, torsor interpretation, and the group scheme \(\operatorname{Pic}^{\natural }(X)\)}
\DAsset{P0068-A02}{map \(\mathbb{Z}^{S}\to \operatorname{Pic}^{\natural }(X)\), Lie-algebra interpretation, and isomorphism (10.3.10.2)}

\DArticleAnchor{0069}{69}{70}{736}{73}{P12}

Lemma (10.3.11). — (i) One has\par
\par\medskip
\DFormula{W^{1}(H^{1}(X,\mathbb{Z}))=\operatorname{Im}(H^{1}(\overline{X},\mathbb{Z})\to H^{1}(X,\mathbb{Z}))}
\par\medskip
\DFormula{W^{0}(H^{1}(X,\mathbb{Z}))=\operatorname{Ker}(H^{1}(X,\mathbb{Z})\to H^{1}(\overline{X}′,\mathbb{Z})).}
\par\medskip
(ii) The spectral sequence defined by the stupid filtration of \([\mathcal{O}_{\overline{X}}\to \overline{r}_{\ast }\Omega ^{1}_{\overline{X}′}(\operatorname{log}\) S)] degenerates and converges to the Hodge filtration of \(H^{\ast }(X,\mathbb{C})\).\par
\par\medskip
Let us prove (i). It is enough to prove the first assertion: the second follows from (8.2.5). Compare the exact sequences\par
\par\medskip
\DAsset{P0069-A01}{source-printed diagram: \(H^{1}(\overline{X},\mathbb{Z})\to H^{1}(X,\mathbb{Z})\xrightarrow{\partial }H^{2}(\overline{X}\) mod \(X,\mathbb{Z})\), above \(H^{1}(\overline{X}′,\mathbb{Z})\to H^{1}(X′,\mathbb{Z})\xrightarrow{\partial }H^{2}(\overline{X}′\) mod \(X′,\mathbb{Z})\to H^{2}(\overline{X}′,\mathbb{Z})\), with vertical arrows}
\par\medskip
It follows from (8.3.7) and (3.2.17) that \(H^{2}(\overline{X}′\) mod \(X,\mathbb{Z})\), and hence \(H^{2}(\overline{X}\) mod \(X,\mathbb{Z})\), is purely of weight 2. In view of (8.2.4), this proves (i).\par
\par\medskip
To prove (ii), we must return to Definition (8.2.1) and use the fact that the simplicial scheme \(((X′/X)^{\Delta _{n}})_{n\geq 0}\) is smooth (the iterated \((n+1)-fold\) fiber product \((X′/X)^{\Delta _{n}}\) is the disjoint union of X′, embedded diagonally, and finitely many points) and admits \(((\overline{X}′/\overline{X})^{\Delta _{n}})_{n\geq 0}\) as a smooth compactification. Let \(\varepsilon :((\overline{X}′/\overline{X})^{\Delta _{n}})_{n\geq 0}\to \overline{X}\) be the augmentation morphism. Let (K,F) be the complex \(\mathcal{O}\to \Omega ^{1}(\operatorname{log}\) S) on \(((\overline{X}′/\overline{X})^{\Delta _{n}})_{n\geq 0}\), equipped with the stupid filtration. There is an augmentation morphism\par
\par\medskip
\DFormula{\alpha :[\mathcal{O}_{\overline{X}}\to \overline{r}_{\ast }\Omega ^{1}_{\overline{X}′}(\operatorname{log}\,S)]\to s\varepsilon _{\bullet \ast }K.}
\par\medskip
The diagram\par
\par\medskip
\DAsset{P0069-A02}{source-printed square: upper row \(H^{1}(X,\mathbb{C})=H^{1}(((\overline{X}′/\overline{X})^{\Delta _{n}})_{\bullet },K)\) \(\xleftarrow{\text{\textcircled{1}}}\) \(H^{1}(\overline{X},s\varepsilon _{\bullet \ast }K)\); lower row \(H^{1}(X,\mathbb{C})=H^{1}(\overline{X},[\mathcal{O}_{\overline{X}}\to \overline{r}_{\ast }\Omega ^{1}_{\overline{X}′}(\operatorname{log}\) S)]); the left side is vertical equality between the two \(H^{1}(X,\mathbb{C})\) terms; the right side is an upward arrow from the lower-right term to \(H^{1}(\overline{X},s\varepsilon _{\bullet \ast }K)\)}
\par\medskip
is commutative, and \(\text{\textcircled{1}}\) is an isomorphism because the \(\varepsilon _{n}\) are finite, so that \(R^{i}\varepsilon _{n\ast }=0\) for \(i>0\). The assertion now follows from the fact that \(\alpha \) is a filtered quasi-isomorphism.\par

\DArticleAnchor{0070}{70}{71}{737}{74}{P12}

(10.3.12) Let us complete Construction (10.3.8). The diagram\par
\par\medskip
\DAsset{P0070-A01}{source-printed diagram: \(H^{1}(X,\mathbb{Z}(1))\simeq T_{\mathbb{Z}}(H_{m}^{1}(X)(1))\to \operatorname{Lie}\) \(\operatorname{Pic}(\overline{X})=H^{1}(\overline{X},\mathcal{O}_{\overline{X}})\) above \(H^{1}(X,\mathbb{C})\simeq \operatorname{Lie}\) \(\operatorname{Pic}^{\natural }(X)\to \operatorname{Lie}\) \(\operatorname{Pic}(\overline{X})=H_{1}(\overline{X},\mathcal{O}_{\overline{X}})\), with vertical arrows; the lower-right cohomology index is printed as \(H_{1}\)}
\par\medskip
is commutative. By (10.3.11) (ii) and the definition of the Hodge filtrations, (10.3.10.2) is compatible with the Hodge filtration.\par
\par\medskip
Corollary (10.3.13). — \(H_{m}^{1}(X)(1)^{\natural }\) is the extension\par
\par\medskip
\DFormula{\operatorname{Ker}(\alpha )\to \operatorname{Pic}^{\natural }(X)^{0}}
\par\medskip
\DFormula{of\,H_{m}^{1}(X)(1)\,obtained\,from\,(10.3.10.1).}
\par\medskip
Indeed, there is a commutative diagram\par
\par\medskip
\DAsset{P0070-A02}{diagram: \(\operatorname{Ker}(\alpha )\to \operatorname{Pic}^{\natural }(X)^{0}\to \operatorname{Pic}^{0}(\overline{X})\) above \(H^{1}(X,\mathbb{Z})\to \operatorname{Lie}\) \(\operatorname{Pic}^{\natural }(X)\to \operatorname{Lie}\) \(\operatorname{Pic}(\overline{X})\), with three upward vertical arrows}
\par\medskip
and one applies (10.1.9).\par
\par\medskip
Exercise (10.3.14). — Verify that the diagram\par
\par\medskip
\DAssetReference{P0070-A02}{square: \(H^{1}(X,\mathbb{Z}(1))=T_{\mathbb{Z}}(H_{m}^{1}(X)(1))\) above \(H^{1}(X,\mu _{n})=T_{\mathbb{Z}/n\mathbb{Z}}(H_{m}^{1}(X)(1))\), with vertical arrows}
\par\medskip
is commutative.\par
\par\medskip
Remark (10.3.15). — The isomorphism deduced from (10.3.8) and (10.3.9) (ii), between \(H^{1}(\overline{X},\mathcal{O}_{\overline{X}}\to \overline{r}_{\ast }\Omega ^{1}_{\overline{X}′}(\operatorname{log}\) S)) (which may be called \(H^{1}_{DR}(X))\) and \(T_{DR}(H_{m}^{1}(X)(1))\), equal to Lie \(\operatorname{Pic}^{\natural }(X)\) by (10.3.13), was defined in a purely algebraic manner.\par

\DArticleAnchor{0071}{71}{72}{738}{75}{P12}

10.4. Translation of a theorem of Picard.\par
\par\medskip
(10.4.1) Let \(\mathcal{H}\) be a mixed Hodge structure such that \(h^{pq}=0\) for p or \(q<0\) (respectively for p or \(q>n)\). I shall provisionally denote by \(I(\mathcal{H})\) (respectively \(II_{n}(\mathcal{H}))\) the 1-motive defined by the largest mixed Hodge substructure of \((\mathcal{H}_{\mathbb{Z}}/torsion)(1)\) (respectively by the largest quotient of \((\mathcal{H}_{\mathbb{Z}}/torsion)(n))\) that is purely of type\par
\par\medskip
\DFormula{\lbrace (−1,−1),(−1,0),(0,−1),(0,0)\rbrace .}
\par\medskip
If X is a complex algebraic variety of dimension \(\leq N\), I conjecture that the 1-motives \(I(H^{n}(X,\mathbb{Z}))\), \(II_{n}(H^{n}(X,\mathbb{Z}))\) (for \(n\leq N)\), and \(II_{N}(H^{n}(X,\mathbb{Z}))\) (for \(n\geq N)\) admit a purely algebraic definition. The morphisms\par
\par\medskip
\DFormula{T_{\ell }(I(H^{n}(X,\mathbb{Z})))\to (H^{n}(X,\mathbb{Z}_{\ell })/torsion)(1)}
\DFormula{T_{\ell }(II_{n}(H^{n}(X,\mathbb{Z})))\leftarrow H^{n}(X,\mathbb{Z}_{\ell })(n)\,\,\,\,\,(for\,n\leq N)}
\DFormula{T_{\ell }(II_{N}(H^{n}(X,\mathbb{Z})))\leftarrow H^{n}(X,\mathbb{Z}_{\ell })(N)\,\,\,\,\,(for\,n\geq N)}
\par\medskip
and their analogues in de Rham cohomology should likewise admit a purely algebraic definition. This is what we have verified when \(\operatorname{dim}(X)=1\). Here is another example.\par
\par\medskip
(10.4.2) Let S be a smooth projective surface, U a dense (Zariski) open subset of S, and \(C=S−U\). Put\par
\par\medskip
\DFormula{H_{2}^{\operatorname{log}}(U,\mathbb{Z})=\operatorname{Ker}(H_{2}(U,\mathbb{Z})\to H_{2}(S,\mathbb{Z})).}
\par\medskip
According to [8] § 9, this group does not depend on the chosen smooth compactification S of U. Modulo finite groups, it is simply (3.2.17) \(W_{−3}H_{2}(U,\mathbb{Z})\) \((H_{2}(U,\mathbb{Q})\), dual to \(H^{2}(U,\mathbb{Q})\), is equipped with the dual mixed Hodge structure). One can verify that the 1-motive defined by \((H_{2}^{\operatorname{log}}(U,\mathbb{Z})/torsion)(−1)\) has the following description:\par
\par\medskip
a) By Poincaré duality, \(H_{2}(U,\mathbb{Z})=H_{c}^{2}(U,\mathbb{Z})\). The long exact cohomology sequence of (S,C) gives an exact sequence\par
\par\medskip
\DFormula{H^{1}(S,\mathbb{Z})\to H^{1}(C,\mathbb{Z})\to H_{2}^{\operatorname{log}}(U,\mathbb{Z})\to 0.}
\par\medskip
b) Let \(H_{m}^{1}(C)\) be the quotient of \(\operatorname{Pic}^{0}(C)\) by its unipotent radical. It is known that \(H^{1}(C,\mathbb{Z})=T_{\mathbb{Z}}(H_{m}^{1}(C))\) and \(H^{1}(S,\mathbb{Z})=T_{\mathbb{Z}}(\operatorname{Pic}^{0}(S))\). There is therefore a map\par
\par\medskip
\DFormula{H_{2}^{\operatorname{log}}(U,\mathbb{Z})\to T_{\mathbb{Z}}(H_{m}^{1}(C)/\operatorname{Pic}^{0}(S)),}
\par\medskip
which is an isomorphism modulo finite groups.\par
\par\medskip
c) The map deduced from b)\par
\par\medskip
\DFormula{(10.4.2.1)\,\,\,\,\,(H_{2}^{\operatorname{log}}(U,\mathbb{Z})/torsion)(−1)\to T_{\mathbb{Z}}(H_{m}^{1}(C)/\operatorname{Pic}^{0}(S))}
\par\medskip
is a morphism of mixed Hodge structures (and an isomorphism modulo finite groups).\par
\DAsset{P0071-A01}{§10.4, provisional 1-motive functors, and \(\ell -adic\) realization maps}
\DAsset{P0071-A02}{surface example, logarithmic homology, and formula (10.4.2.1)}

\DArticleAnchor{0072}{72}{73}{739}{76}{P12}

One can define in a purely algebraic manner the isomorphisms (modulo finite groups) deduced from (10.4.2.1) between the analogues of \(H_{2}^{\operatorname{log}}(U,\mathbb{Z})(−1)\) in \(\ell -adic\) or de Rham cohomology and the \(\ell -adic\) or de Rham realizations of \(H_{m}^{1}(C)/\operatorname{Pic}^{0}(S)\).\par
\par\medskip
(10.4.3) We shall not prove the preceding constructions and compatibilities. The result is implicit in the following theorem of E. Picard ([11], volume I, p. 62). Here P and A denote polynomials.\par
\par\medskip
Every residue of the double integral\par
\par\medskip
\DFormula{∬\,P(x,y)dxdy/A(x,y)}
\par\medskip
may be regarded as a logarithmic or cyclic period of an Abelian integral.\par
\par\medskip
I recall that:\par
\par\medskip
a) A k-fold integral on a nonsingular projective variety V of dimension d is a closed rational k-form on V: it is an element \(\alpha \in H^{0}(U,\Omega ^{k})\), for U a dense open subset of V, satisfying \(d\alpha =0\). Here \(k=2\) and \(V=ℙ^{2}\), so that \(k=\operatorname{dim}\) V and the condition \(d\alpha =0\) is vacuous.\par
\par\medskip
b) A k-fold integral \(\alpha \) defines an element, again denoted \(\alpha \), of \(H^{k}(U,\mathbb{C})\). The periods of \(\alpha \) are the numbers \(⟨\alpha ,c⟩\) for \(c\in H_{k}(U,\mathbb{Z})\). If c is defined by a singular chain C, then\par
\par\medskip
\DFormula{⟨\alpha ,c⟩=∫_{C}\,\alpha .}
\par\medskip
c) The residues of a k-fold integral \(\alpha \in H^{0}(U,\Omega ^{k})\) are the numbers \(⟨\alpha ,c⟩\) for\par
\par\medskip
\DFormula{c\in \operatorname{Ker}(H_{k}(U,\mathbb{Z})\to H_{k}(V,\mathbb{Z})).}
\par\medskip
In the case considered by Picard, the map \(H_{2}(U,\mathbb{Z})\to H_{2}(ℙ^{2},\mathbb{Z})\) is zero, and there is no need to distinguish between periods and residues. For the general case, cf. Lefschetz, [10], note I.\par
\par\medskip
d) By “logarithmic or cyclic period of an Abelian integral,” Picard means simply “period of a simple integral \(\beta \) on a curve Y.” It is understood that \(\beta \) and Y are determined rationally by P and A.\par
\par\medskip
(10.4.4) Let us translate this theorem. Let S and U be as in (10.4.2) (for example \(S=ℙ^{2}(\mathbb{C}))\), let a double integral \(\beta \in H^{0}(U,\Omega _{U}^{2})\) be given, and let \(c\in H_{2}^{\operatorname{log}}(U,\mathbb{Z})\) be a cycle. By definition\par
\par\medskip
\DFormula{⟨\beta ,c⟩=∫_{c}\,\beta }
\par\medskip
is a residue of \(\beta \). Let \(A=H_{m}^{1}(C)(1)/\operatorname{Pic}^{0}(S)\) (an extension of an abelian variety by a torus). The 2-form \(\beta \) defines a cohomology class \(cl(\beta )\in H_{DR}^{2}(U)\) and induces\par
\DAsset{P0072-A01}{continuation of (10.4.2), Picard’s double integral, and period/residue definitions}
\DAsset{P0072-A02}{opening of (10.4.4), de Rham class, and live continuation into article page 73}

\DArticleAnchor{0073}{73}{74}{740}{77}{P13}

a linear form \(\beta ′\) on \(H_{DR2}^{\operatorname{log}}(U)=T_{DR}(A)\). The correspondence \(\beta ↦\beta ′\) is rational. The class c defines \(c′\in T_{\mathbb{Z}}(A)\otimes \mathbb{Q}\), and one has\par
\par\medskip
\DFormula{(10.4.4.1)\,\,\,\,\,⟨\beta ,c⟩=2\pi i⟨\beta ′,c′⟩.}
\par\medskip
One can represent c′ by a cycle on C, and \(\beta ′\) as a simple integral on C.\par
\par\medskip
BIBLIOGRAPHY\par
\par\medskip
[1] A. Blanchard, Sur les variétés analytiques complexes, Ann. Sc. E.N.S., 73 (1956).\par
[2] T. Bloom and M. Herrera, De Rham cohomology of an analytic space, Inv. Math., 7 (1969), 275-296.\par
[3] P. Deligne, Théorème de Lefschetz et critères de dégénérescence de suites spectrales, Publ. Math. I.H.E.S., 35 (1968), 107-126.\par
[4] P. Deligne, Théorie de Hodge I, Actes du Congrès international des Mathématiciens (Nice, 1970), Gauthier-Villars, 1971, 1, 425-430.\par
[5] P. Deligne, Théorie de Hodge II, Publ. Math. I.H.E.S., 40 (1971), 5-58.\par
[6] P. A. Griffiths, Periods of integrals on algebraic manifolds. Summary of main results and discussion of open problems, Bull. Am. Math. Soc., 76 (1970), 228-296.\par
[7] P. A. Griffiths, On the periods of certain rational integrals : I, Annals of Math., 90 (1969), 460-495.\par
[8] A. Grothendieck, Le groupe de Brauer III : Exemples et compléments, in Dix exposés sur la cohomologie des schémas, North Holland Publ. Co., 88-188.\par
[9] A. Grothendieck, On the De Rham cohomology of algebraic varieties, Publ. Math. I.H.E.S., 29 (1966), 96-103.\par
[10] S. Lefschetz, L’analysis situs et la géométrie algébrique, Gauthier-Villars, 1924 (reprinted in Selected Papers, Chelsea Publ. Co., 1971).\par
[11] E. Picard and G. Simart, Théorie des fonctions algébriques de deux variables indépendantes, vol. I and II, Paris, Gauthier-Villars, 1895, 1906.\par
[12] D. G. Quillen, Notes on the homology of commutative rings, M.I.T., 1968.\par
[13] B. Saint-Donat (from notes by P. Deligne), Techniques de descente cohomologique, S.G.A. 4 V bis (to appear with S.G.A. 4, in the Lecture Notes, Springer-Verlag).\par
[14] G. Segal, Classifying spaces and spectral sequences, Publ. Math. I.H.E.S., 34 (1968), 105-112.\par
[15] J.-P. Serre, Géométrie algébrique et géométrie analytique, Ann. Inst. Fourier, 6 (1966), cited as G.A.G.A.\par
\par\medskip
ABBREVIATIONS\par
\par\medskip
S.G.A., Séminaire de Géométrie algébrique du Bois-Marie, 1 to 7, distributed by I.H.E.S.; 1, 3, 4, 6, and 7 I appeared in the Springer Lecture Notes, no. 224; 151, 152 and 153; 269, 270, 305; 225, 288 — S.G.A. 2 was published by North Holland Publ. Co.\par
\par\medskip
Manuscript received March 2, 1972.\par
\DAsset{P0073-A01}{completion of (10.4.4), logarithmic de Rham pairing, and formula (10.4.4.1)}
\DAsset{P0073-A02}{bibliography, abbreviations, and manuscript-received line}



\end{document}
